% Options for packages loaded elsewhere
% Options for packages loaded elsewhere
\PassOptionsToPackage{unicode,linktoc=all}{hyperref}
\PassOptionsToPackage{hyphens}{url}
\PassOptionsToPackage{dvipsnames,svgnames,x11names}{xcolor}
%
\documentclass[
  letterpaper,
]{scrbook}
\usepackage{xcolor}
\usepackage[top=30mm,left=20mm,heightrounded]{geometry}
\usepackage{amsmath,amssymb}
\setcounter{secnumdepth}{5}
\usepackage{iftex}
\ifPDFTeX
  \usepackage[T1]{fontenc}
  \usepackage[utf8]{inputenc}
  \usepackage{textcomp} % provide euro and other symbols
\else % if luatex or xetex
  \usepackage{unicode-math} % this also loads fontspec
  \defaultfontfeatures{Scale=MatchLowercase}
  \defaultfontfeatures[\rmfamily]{Ligatures=TeX,Scale=1}
\fi
\usepackage{lmodern}
\ifPDFTeX\else
  % xetex/luatex font selection
  \setmainfont[]{Calibri}
\fi
% Use upquote if available, for straight quotes in verbatim environments
\IfFileExists{upquote.sty}{\usepackage{upquote}}{}
\IfFileExists{microtype.sty}{% use microtype if available
  \usepackage[]{microtype}
  \UseMicrotypeSet[protrusion]{basicmath} % disable protrusion for tt fonts
}{}
\makeatletter
\@ifundefined{KOMAClassName}{% if non-KOMA class
  \IfFileExists{parskip.sty}{%
    \usepackage{parskip}
  }{% else
    \setlength{\parindent}{0pt}
    \setlength{\parskip}{6pt plus 2pt minus 1pt}}
}{% if KOMA class
  \KOMAoptions{parskip=half}}
\makeatother
% Make \paragraph and \subparagraph free-standing
\makeatletter
\ifx\paragraph\undefined\else
  \let\oldparagraph\paragraph
  \renewcommand{\paragraph}{
    \@ifstar
      \xxxParagraphStar
      \xxxParagraphNoStar
  }
  \newcommand{\xxxParagraphStar}[1]{\oldparagraph*{#1}\mbox{}}
  \newcommand{\xxxParagraphNoStar}[1]{\oldparagraph{#1}\mbox{}}
\fi
\ifx\subparagraph\undefined\else
  \let\oldsubparagraph\subparagraph
  \renewcommand{\subparagraph}{
    \@ifstar
      \xxxSubParagraphStar
      \xxxSubParagraphNoStar
  }
  \newcommand{\xxxSubParagraphStar}[1]{\oldsubparagraph*{#1}\mbox{}}
  \newcommand{\xxxSubParagraphNoStar}[1]{\oldsubparagraph{#1}\mbox{}}
\fi
\makeatother

\usepackage{color}
\usepackage{fancyvrb}
\newcommand{\VerbBar}{|}
\newcommand{\VERB}{\Verb[commandchars=\\\{\}]}
\DefineVerbatimEnvironment{Highlighting}{Verbatim}{commandchars=\\\{\}}
% Add ',fontsize=\small' for more characters per line
\newenvironment{Shaded}{}{}
\newcommand{\AlertTok}[1]{\textcolor[rgb]{1.00,0.33,0.33}{\textbf{#1}}}
\newcommand{\AnnotationTok}[1]{\textcolor[rgb]{0.42,0.45,0.49}{#1}}
\newcommand{\AttributeTok}[1]{\textcolor[rgb]{0.84,0.23,0.29}{#1}}
\newcommand{\BaseNTok}[1]{\textcolor[rgb]{0.00,0.36,0.77}{#1}}
\newcommand{\BuiltInTok}[1]{\textcolor[rgb]{0.84,0.23,0.29}{#1}}
\newcommand{\CharTok}[1]{\textcolor[rgb]{0.01,0.18,0.38}{#1}}
\newcommand{\CommentTok}[1]{\textcolor[rgb]{0.42,0.45,0.49}{#1}}
\newcommand{\CommentVarTok}[1]{\textcolor[rgb]{0.42,0.45,0.49}{#1}}
\newcommand{\ConstantTok}[1]{\textcolor[rgb]{0.00,0.36,0.77}{#1}}
\newcommand{\ControlFlowTok}[1]{\textcolor[rgb]{0.84,0.23,0.29}{#1}}
\newcommand{\DataTypeTok}[1]{\textcolor[rgb]{0.84,0.23,0.29}{#1}}
\newcommand{\DecValTok}[1]{\textcolor[rgb]{0.00,0.36,0.77}{#1}}
\newcommand{\DocumentationTok}[1]{\textcolor[rgb]{0.42,0.45,0.49}{#1}}
\newcommand{\ErrorTok}[1]{\textcolor[rgb]{1.00,0.33,0.33}{\underline{#1}}}
\newcommand{\ExtensionTok}[1]{\textcolor[rgb]{0.84,0.23,0.29}{\textbf{#1}}}
\newcommand{\FloatTok}[1]{\textcolor[rgb]{0.00,0.36,0.77}{#1}}
\newcommand{\FunctionTok}[1]{\textcolor[rgb]{0.44,0.26,0.76}{#1}}
\newcommand{\ImportTok}[1]{\textcolor[rgb]{0.01,0.18,0.38}{#1}}
\newcommand{\InformationTok}[1]{\textcolor[rgb]{0.42,0.45,0.49}{#1}}
\newcommand{\KeywordTok}[1]{\textcolor[rgb]{0.84,0.23,0.29}{#1}}
\newcommand{\NormalTok}[1]{\textcolor[rgb]{0.14,0.16,0.18}{#1}}
\newcommand{\OperatorTok}[1]{\textcolor[rgb]{0.14,0.16,0.18}{#1}}
\newcommand{\OtherTok}[1]{\textcolor[rgb]{0.44,0.26,0.76}{#1}}
\newcommand{\PreprocessorTok}[1]{\textcolor[rgb]{0.84,0.23,0.29}{#1}}
\newcommand{\RegionMarkerTok}[1]{\textcolor[rgb]{0.42,0.45,0.49}{#1}}
\newcommand{\SpecialCharTok}[1]{\textcolor[rgb]{0.00,0.36,0.77}{#1}}
\newcommand{\SpecialStringTok}[1]{\textcolor[rgb]{0.01,0.18,0.38}{#1}}
\newcommand{\StringTok}[1]{\textcolor[rgb]{0.01,0.18,0.38}{#1}}
\newcommand{\VariableTok}[1]{\textcolor[rgb]{0.89,0.38,0.04}{#1}}
\newcommand{\VerbatimStringTok}[1]{\textcolor[rgb]{0.01,0.18,0.38}{#1}}
\newcommand{\WarningTok}[1]{\textcolor[rgb]{1.00,0.33,0.33}{#1}}

\usepackage{longtable,booktabs,array}
\usepackage{calc} % for calculating minipage widths
% Correct order of tables after \paragraph or \subparagraph
\usepackage{etoolbox}
\makeatletter
\patchcmd\longtable{\par}{\if@noskipsec\mbox{}\fi\par}{}{}
\makeatother
% Allow footnotes in longtable head/foot
\IfFileExists{footnotehyper.sty}{\usepackage{footnotehyper}}{\usepackage{footnote}}
\makesavenoteenv{longtable}
\usepackage{graphicx}
\makeatletter
\newsavebox\pandoc@box
\newcommand*\pandocbounded[1]{% scales image to fit in text height/width
  \sbox\pandoc@box{#1}%
  \Gscale@div\@tempa{\textheight}{\dimexpr\ht\pandoc@box+\dp\pandoc@box\relax}%
  \Gscale@div\@tempb{\linewidth}{\wd\pandoc@box}%
  \ifdim\@tempb\p@<\@tempa\p@\let\@tempa\@tempb\fi% select the smaller of both
  \ifdim\@tempa\p@<\p@\scalebox{\@tempa}{\usebox\pandoc@box}%
  \else\usebox{\pandoc@box}%
  \fi%
}
% Set default figure placement to htbp
\def\fps@figure{htbp}
\makeatother


% definitions for citeproc citations
\NewDocumentCommand\citeproctext{}{}
\NewDocumentCommand\citeproc{mm}{%
  \begingroup\def\citeproctext{#2}\cite{#1}\endgroup}
\makeatletter
 % allow citations to break across lines
 \let\@cite@ofmt\@firstofone
 % avoid brackets around text for \cite:
 \def\@biblabel#1{}
 \def\@cite#1#2{{#1\if@tempswa , #2\fi}}
\makeatother
\newlength{\cslhangindent}
\setlength{\cslhangindent}{1.5em}
\newlength{\csllabelwidth}
\setlength{\csllabelwidth}{3em}
\newenvironment{CSLReferences}[2] % #1 hanging-indent, #2 entry-spacing
 {\begin{list}{}{%
  \setlength{\itemindent}{0pt}
  \setlength{\leftmargin}{0pt}
  \setlength{\parsep}{0pt}
  % turn on hanging indent if param 1 is 1
  \ifodd #1
   \setlength{\leftmargin}{\cslhangindent}
   \setlength{\itemindent}{-1\cslhangindent}
  \fi
  % set entry spacing
  \setlength{\itemsep}{#2\baselineskip}}}
 {\end{list}}
\usepackage{calc}
\newcommand{\CSLBlock}[1]{\hfill\break\parbox[t]{\linewidth}{\strut\ignorespaces#1\strut}}
\newcommand{\CSLLeftMargin}[1]{\parbox[t]{\csllabelwidth}{\strut#1\strut}}
\newcommand{\CSLRightInline}[1]{\parbox[t]{\linewidth - \csllabelwidth}{\strut#1\strut}}
\newcommand{\CSLIndent}[1]{\hspace{\cslhangindent}#1}



\setlength{\emergencystretch}{3em} % prevent overfull lines

\providecommand{\tightlist}{%
  \setlength{\itemsep}{0pt}\setlength{\parskip}{0pt}}



 


\makeatletter
\@ifpackageloaded{tcolorbox}{}{\usepackage[skins,breakable]{tcolorbox}}
\@ifpackageloaded{fontawesome5}{}{\usepackage{fontawesome5}}
\definecolor{quarto-callout-color}{HTML}{909090}
\definecolor{quarto-callout-note-color}{HTML}{0758E5}
\definecolor{quarto-callout-important-color}{HTML}{CC1914}
\definecolor{quarto-callout-warning-color}{HTML}{EB9113}
\definecolor{quarto-callout-tip-color}{HTML}{00A047}
\definecolor{quarto-callout-caution-color}{HTML}{FC5300}
\definecolor{quarto-callout-color-frame}{HTML}{acacac}
\definecolor{quarto-callout-note-color-frame}{HTML}{4582ec}
\definecolor{quarto-callout-important-color-frame}{HTML}{d9534f}
\definecolor{quarto-callout-warning-color-frame}{HTML}{f0ad4e}
\definecolor{quarto-callout-tip-color-frame}{HTML}{02b875}
\definecolor{quarto-callout-caution-color-frame}{HTML}{fd7e14}
\makeatother
\makeatletter
\@ifpackageloaded{bookmark}{}{\usepackage{bookmark}}
\makeatother
\makeatletter
\@ifpackageloaded{caption}{}{\usepackage{caption}}
\AtBeginDocument{%
\ifdefined\contentsname
  \renewcommand*\contentsname{Table of contents}
\else
  \newcommand\contentsname{Table of contents}
\fi
\ifdefined\listfigurename
  \renewcommand*\listfigurename{List of Figures}
\else
  \newcommand\listfigurename{List of Figures}
\fi
\ifdefined\listtablename
  \renewcommand*\listtablename{List of Tables}
\else
  \newcommand\listtablename{List of Tables}
\fi
\ifdefined\figurename
  \renewcommand*\figurename{Figure}
\else
  \newcommand\figurename{Figure}
\fi
\ifdefined\tablename
  \renewcommand*\tablename{Table}
\else
  \newcommand\tablename{Table}
\fi
}
\@ifpackageloaded{float}{}{\usepackage{float}}
\floatstyle{ruled}
\@ifundefined{c@chapter}{\newfloat{codelisting}{h}{lop}}{\newfloat{codelisting}{h}{lop}[chapter]}
\floatname{codelisting}{Listing}
\newcommand*\listoflistings{\listof{codelisting}{List of Listings}}
\makeatother
\makeatletter
\makeatother
\makeatletter
\@ifpackageloaded{caption}{}{\usepackage{caption}}
\@ifpackageloaded{subcaption}{}{\usepackage{subcaption}}
\makeatother
\usepackage{bookmark}
\IfFileExists{xurl.sty}{\usepackage{xurl}}{} % add URL line breaks if available
\urlstyle{same}
\hypersetup{
  pdftitle={IC Design of a Universal Biquad Filter},
  pdfauthor={Atakan Baydogan; Onno Kriens; Finn Ringelsiep; Alicia von Ahlen},
  colorlinks=true,
  linkcolor={blue},
  filecolor={Maroon},
  citecolor={Blue},
  urlcolor={Blue},
  pdfcreator={LaTeX via pandoc}}


\title{IC Design of a Universal Biquad Filter}
\author{Atakan Baydogan \and Onno Kriens \and Finn
Ringelsiep \and Alicia von Ahlen}
\date{2025-07-16}
\begin{document}
\frontmatter
\maketitle

\renewcommand*\contentsname{Table of contents}
{
\hypersetup{linkcolor=}
\setcounter{tocdepth}{2}
\tableofcontents
}

\mainmatter
\bookmarksetup{startatroot}

\chapter*{Abstract}\label{abstract}
\addcontentsline{toc}{chapter}{Abstract}

\markboth{Abstract}{Abstract}

\textbf{Keywords:} Filter design, Biquadratic filter, IC design, Open
Source Toolchain

This project explores the IC design of a biquad filter from mathematical
modelling over ideal circuit implementations to sizing and using several
differnet OTAs and a concluding physical implementation of an OTA. Two
different OTA models were incorporated into implementing the universal
biquad filter and compared with simulations to the theoretical model.
Additionally, a low pass filter was designed with a Gm-C filter to try
to realise the required behauviour that way. To experience the the
actual physical layouting of an IC chip, the 5T-OTA was implemented as a
physical design with the help of KLayout.

\bookmarksetup{startatroot}

\chapter{Introduction}\label{introduction}

This chapter provides a short motivation and overview for the IC design
process done during the lecture ``Analogue and Mixed-Signal Circuit
Design'' that Prof.~Dr.-Ing. M. Meiners gives in the graduate course
Electronic Engineering M.Sc. at City University of Applied Sciences
Bremen.

The chapter will start with the motivation behind the biquad IC filter
design, outline the scope of work of the project, and ends with giving
concrete specifications for the implemented filter.

\section{Motivation}\label{motivation}

The design and implementation of analog filters is a cornerstone in
signal processing, with applications ranging from audio processing to
communication systems. Among these, second order filters, like the
biquad filter, are versatile building blocks due to its ability to
realize four types of second order filters - low pass, high pass, band
pass, and band stop. This project focuses on the integrated circuit (IC)
design of a biquad filter, to get insight into the theoretical and
practical engineering considerations behind IC design.

For a deeper understanding of IC design, this project does not rely on
off-the-shelf operational amplifiers for the filter design, but aims to
implement the entire filter architecture at the transistor level. This
approach not only deepens the understanding of analog filter behavior
but also introduces the challenges and intricacies of IC design, such as
layout constraints, power efficiency, and stability.

This project demonstrates the design process of IC design from
theorectical modelling, over simulation and design constraints to
prototyping and to learn hands-on experience with tools used during the
design process.

\section{Scope of the project}\label{scope-of-the-project}

The scope of this project is supposed to follow a real-world design
flow, starting at a theorectical analysis of the specified filter and -
in the best case - end in a tape-out of a prototype. If that stage is
reached the prototype can be compared to the theorectical and simulation
results obtained during the design process and checked for
functionality.

As a tape-out of a prototype is fairly unrealistic in the time given,
the goal is to simulate the specified filter with templates for
operational amplifiers and base the IC layout on these templates.

All in all, this project includes a systems analysis if the specified
filter, simulation results with ideal components and real components,
taken from provided templates, and a physical layout prototype.

\section{Specifications}\label{specifications}

The main objective is to design a universal biquad filter, based on the
filter design proposed in the ASLK PRO Board Manual from Texas
Instruments (Rao and Ravikumar 2012). The biquad filter shall have the
following specifications:

\[
f_0 = 1\,kHz
\]

\[
Q = 10
\]

The circuit design is done in \textbf{Xschem} and the simulation in
\textbf{ngspice}. For the design on transistor level the 130nm CMOS
technology \textbf{SG13G2} is used. All these tools and PDk are
integrated into a docker image \textbf{IIC-OSIC-TOOLS} (Pretl and Zachl
2025) provided by Prof.~Dr.~Harald Pretl from Johannes Keplar
University.

This documentation provides a development report, which documents the
design process with the taken steps and decisions made.

\section{Open-Source}\label{open-source}

All the results of this report and development approach to design a
Biquad will be publicly available on GitHub. Everyone is invited and
should feel free to use, change, and share this work. This whole course
and project wouldn't be possible without the great Open-Source tools
provided by the amazing community of layout designers, enthusiasts, and
developers. Here is a list with just a few of these programs:
\texttt{IC-OSIC-TOOLS}, \texttt{IHP\ Open\ PDK}, \texttt{Linux},
\texttt{Docker}, \texttt{Xschem}, \texttt{ngspice}, \texttt{KLayout},
\texttt{Quarto}, \texttt{Vim}, \texttt{Pandoc}, \texttt{LaTeX},
\texttt{TeXLive}, \texttt{Python}, \texttt{Git}, \texttt{CoCalc},
\texttt{LibreOffice}.

\subsection{Anthem}\label{anthem}

In realms where code is free to fly, We build and share, our hearts
reach high. No walls, no locks, our wisdom streams, In open light, we
chase our dreams.

So sing the joy, the thrill, the spark, In open source, we find our arc.
Together strong, we rise, explore--- In lines of code, forevermore.

\bookmarksetup{startatroot}

\chapter{Theoretical Background}\label{theoretical-background}

This chapter introduces the theory and core concepts necessary for IC
design of a biquad filter. It is structured in a way, that it goes from
the big picture to the small components. First, biquad filters are
introduced with a focus on the universal biquad filter. After that,
operational amplifiers come into the the foreground, as biquad filters
make use of them in their circuits. Operational amplifiers are looked
upon from an IC design standpoint.

\section{Biquad filter}\label{biquad-filter}

The biquadratic filter, also known as the biquad filter, has its
earliest implementation in the 1960s but is still in use today, most
commonly in radio frequency receivers (Razavi 2018). In its application
in RF-technology, it is used to remove unwanted neighboring signals ans
noise (Razavi 2024). As biquad filter are second-order filters, they are
also used as building blocks for higher filter implementations, by
cascading them and adding first order filters (Rao and Ravikumar 2012).

\subsection{Universal biquad filter}\label{universal-biquad-filter}

For the filter design in this project, an universal biquad filter is
used. The universal biquad filter is biquad filter variant with four
operational amplifiers used in its design and the property of being able
to be used in four different filter variants. Depending on the output of
the universal biquad used, a low pass filter, high pass filter, band
pass filter, or band stop filter will be implemented. This can be seen
in Figure~\ref{fig-BiquadCircuit}. (Rao and Ravikumar 2012)

\begin{figure}

\centering{

\pandocbounded{\includegraphics[keepaspectratio]{images/sec_theorectical_background/biquad_circuit.png}}

}

\caption{\label{fig-BiquadCircuit}circuit design of an universal biquad
filter}

\end{figure}%

The universal biquad filter consists of two non-inverting amplifiers
working as adders in the circuit and two integrators. By setting \(R\)
and \(C\) to specific values, the resonance frequency can be chosen.
Other parameters adjustable in the universal biquad filter are the
quality factor \(Q\) and the low-frequency gain \(H_0\). The quality
factor and low-frequency gain determine the frequency response peaks of
the low pass filter and the band pass filter. (Rao and Ravikumar 2012)

\subsection{Characteristics}\label{characteristics}

As the universal biquad filter is a second order filter with the
specified outputs low pass, high pass, band pass, and band stop, each
filter option can be described with a tranfer function on system level.
The general second order transfer function is:

\begin{equation}\phantomsection\label{eq-TFSecondOrder}{
H(s) = \frac{a_1 s^2 + b_1 s + c_1}{a_2 s^2 + b_2 s + c_2}
}\end{equation}

The coeffcients can be choosen so, that different responses, like low
pass, high pass, band pass, and band stop are achieved.

In filter design a variant of this generalized transfer function is
often choosen because it is easier describe the system by quality factor
and angular frequency. Equation~\ref{eq-TFSecondOrderBiquad} is an
exemplatory low pass filter with a transfer function specified for
filter design. (Razavi 2018)

\begin{equation}\phantomsection\label{eq-TFSecondOrderBiquad}{
H(s) = \frac{\omega_n^2}{s^2 + \frac{\omega_n}{Q} s + \omega_n^2}
}\end{equation}

\(Q\) determines amoung other things the amount of peaking the transfer
function has at the chosen frequency. Figure~\ref{fig-lowPassDifferentQ}
shows this graphically, the amout of peaking increases with increasing
quality factor \(Q\).

\begin{codelisting}

\caption{\label{lst-lowpassdifferentq}}

\centering{

\begin{Shaded}
\begin{Highlighting}[]
\CommentTok{\# Behavioral Analysis Biquad Filter}

\ImportTok{import}\NormalTok{ numpy }\ImportTok{as}\NormalTok{ np}
\ImportTok{import}\NormalTok{ matplotlib.pyplot }\ImportTok{as}\NormalTok{ plt}

\CommentTok{\# Initial values}
\NormalTok{f0 }\OperatorTok{=} \FloatTok{1e3}  \CommentTok{\# Resonance frequency in Hz}
\NormalTok{w0 }\OperatorTok{=} \DecValTok{2} \OperatorTok{*}\NormalTok{ np.pi }\OperatorTok{*}\NormalTok{ f0  }\CommentTok{\# Angular frequency in rad/s}
\NormalTok{Q1 }\OperatorTok{=} \DecValTok{1}  \CommentTok{\# Quality factor}
\NormalTok{Q2 }\OperatorTok{=} \DecValTok{2}
\NormalTok{Q3 }\OperatorTok{=} \DecValTok{5}
\NormalTok{Q4 }\OperatorTok{=} \DecValTok{10}
\NormalTok{Q5 }\OperatorTok{=} \DecValTok{100}
\NormalTok{H0 }\OperatorTok{=} \DecValTok{1}  \CommentTok{\# Play around with this later}

\CommentTok{\# Logarithmic frequency axis}
\NormalTok{frequencies }\OperatorTok{=}\NormalTok{ np.logspace(}\DecValTok{2}\NormalTok{, }\DecValTok{4}\NormalTok{, }\DecValTok{10000}\NormalTok{)  }\CommentTok{\# Frequency from 10\^{}2 to 10\^{}4 Hz}
\NormalTok{s }\OperatorTok{=} \OtherTok{1j} \OperatorTok{*} \DecValTok{2} \OperatorTok{*}\NormalTok{ np.pi }\OperatorTok{*}\NormalTok{ frequencies  }\CommentTok{\# Laplace{-}Variable s = jω}

\CommentTok{\#\#\#\#\#\#\#\#\#\#\#\#\#\#\#\#\#\#\#\#\#\#\#\#\#\#\#\#\#\#\#\#\#\#\#\#\#\#\#\#\#\#\#\#}
\CommentTok{\# Transfer functions of Active Filters}
\CommentTok{\#\#\#\#\#\#\#\#\#\#\#\#\#\#\#\#\#\#\#\#\#\#\#\#\#\#\#\#\#\#\#\#\#\#\#\#\#\#\#\#\#\#\#\#}

\CommentTok{\#\#\# Numerator}
\CommentTok{\# Low Pass Filter}
\NormalTok{b\_lp }\OperatorTok{=}\NormalTok{ H0}

\CommentTok{\# High Pass Filter}
\CommentTok{\#b\_hp = (H0 * (s**2 / w0**2))}

\CommentTok{\# Band Pass Filter}
\CommentTok{\#b\_bp = ({-}H0 * (s / w0))}

\CommentTok{\# Band Stop Filter}
\CommentTok{\#b\_bs = {-}((1 + (s**2 / (w0**2))) * H0)}

\CommentTok{\# Denominator {-}\textgreater{} for all filters the same}
\NormalTok{a0 }\OperatorTok{=} \DecValTok{1}
\NormalTok{a1\_1 }\OperatorTok{=}\NormalTok{ (s }\OperatorTok{/}\NormalTok{ (w0 }\OperatorTok{*}\NormalTok{ Q1))}
\NormalTok{a1\_2 }\OperatorTok{=}\NormalTok{ (s }\OperatorTok{/}\NormalTok{ (w0 }\OperatorTok{*}\NormalTok{ Q2))}
\NormalTok{a1\_3 }\OperatorTok{=}\NormalTok{ (s }\OperatorTok{/}\NormalTok{ (w0 }\OperatorTok{*}\NormalTok{ Q3))}
\NormalTok{a1\_4 }\OperatorTok{=}\NormalTok{ (s }\OperatorTok{/}\NormalTok{ (w0 }\OperatorTok{*}\NormalTok{ Q4))}
\NormalTok{a1\_5 }\OperatorTok{=}\NormalTok{ (s }\OperatorTok{/}\NormalTok{ (w0 }\OperatorTok{*}\NormalTok{ Q5))}
\NormalTok{a2 }\OperatorTok{=}\NormalTok{ (s}\OperatorTok{**}\DecValTok{2} \OperatorTok{/}\NormalTok{ (w0}\OperatorTok{**}\DecValTok{2}\NormalTok{))}

\NormalTok{den1 }\OperatorTok{=}\NormalTok{ a0 }\OperatorTok{+}\NormalTok{ a1\_1 }\OperatorTok{+}\NormalTok{ a2}
\NormalTok{den2 }\OperatorTok{=}\NormalTok{ a0 }\OperatorTok{+}\NormalTok{ a1\_2 }\OperatorTok{+}\NormalTok{ a2}
\NormalTok{den3 }\OperatorTok{=}\NormalTok{ a0 }\OperatorTok{+}\NormalTok{ a1\_3 }\OperatorTok{+}\NormalTok{ a2}
\NormalTok{den4 }\OperatorTok{=}\NormalTok{ a0 }\OperatorTok{+}\NormalTok{ a1\_4 }\OperatorTok{+}\NormalTok{ a2}
\NormalTok{den5 }\OperatorTok{=}\NormalTok{ a0 }\OperatorTok{+}\NormalTok{ a1\_5 }\OperatorTok{+}\NormalTok{ a2}

\CommentTok{\#\#\#\#\#\#\#\#\#\#\#\#\#\#\#\#\#\#\#\#\#\#\#\#\#\#\#\#\#\#\#\#\#\#\#\#\#\#\#\#\#\#\#\#}
\CommentTok{\# Calculation of the transfer functions H(s)}
\CommentTok{\#\#\#\#\#\#\#\#\#\#\#\#\#\#\#\#\#\#\#\#\#\#\#\#\#\#\#\#\#\#\#\#\#\#\#\#\#\#\#\#\#\#\#\#}
\NormalTok{Hs\_lp\_1 }\OperatorTok{=}\NormalTok{ b\_lp }\OperatorTok{/}\NormalTok{ den1}
\NormalTok{Hs\_lp\_2 }\OperatorTok{=}\NormalTok{ b\_lp }\OperatorTok{/}\NormalTok{ den2}
\NormalTok{Hs\_lp\_3 }\OperatorTok{=}\NormalTok{ b\_lp }\OperatorTok{/}\NormalTok{ den3}
\NormalTok{Hs\_lp\_4 }\OperatorTok{=}\NormalTok{ b\_lp }\OperatorTok{/}\NormalTok{ den4}
\NormalTok{Hs\_lp\_5 }\OperatorTok{=}\NormalTok{ b\_lp }\OperatorTok{/}\NormalTok{ den5}
\CommentTok{\#Hs\_hp = b\_hp / den}
\CommentTok{\#Hs\_bp = b\_bp / den}
\CommentTok{\#Hs\_bs = b\_bs / den}

\CommentTok{\# Bode Diagram}
\NormalTok{fig, axs }\OperatorTok{=}\NormalTok{ plt.subplots(}\DecValTok{2}\NormalTok{)}
\CommentTok{\#fig.suptitle("frequency response of biquad filter")}

\CommentTok{\# Low Pass Filter}
\NormalTok{axs[}\DecValTok{0}\NormalTok{].semilogx(frequencies, }\DecValTok{20} \OperatorTok{*}\NormalTok{ np.log10(np.}\BuiltInTok{abs}\NormalTok{(Hs\_lp\_1)), label}\OperatorTok{=}\StringTok{\textquotesingle{}$Q = 1$\textquotesingle{}}\NormalTok{)}
\NormalTok{axs[}\DecValTok{1}\NormalTok{].semilogx(frequencies, np.unwrap(np.angle(Hs\_lp\_1)) }\OperatorTok{*}\NormalTok{ (}\DecValTok{180} \OperatorTok{/}\NormalTok{ np.pi), label}\OperatorTok{=}\StringTok{\textquotesingle{}$Q = 1$\textquotesingle{}}\NormalTok{)}
\NormalTok{axs[}\DecValTok{0}\NormalTok{].semilogx(frequencies, }\DecValTok{20} \OperatorTok{*}\NormalTok{ np.log10(np.}\BuiltInTok{abs}\NormalTok{(Hs\_lp\_2)), label}\OperatorTok{=}\StringTok{\textquotesingle{}$Q = 2$\textquotesingle{}}\NormalTok{)}
\NormalTok{axs[}\DecValTok{1}\NormalTok{].semilogx(frequencies, np.unwrap(np.angle(Hs\_lp\_2)) }\OperatorTok{*}\NormalTok{ (}\DecValTok{180} \OperatorTok{/}\NormalTok{ np.pi), label}\OperatorTok{=}\StringTok{\textquotesingle{}$Q = 2$\textquotesingle{}}\NormalTok{)}
\NormalTok{axs[}\DecValTok{0}\NormalTok{].semilogx(frequencies, }\DecValTok{20} \OperatorTok{*}\NormalTok{ np.log10(np.}\BuiltInTok{abs}\NormalTok{(Hs\_lp\_3)), label}\OperatorTok{=}\StringTok{\textquotesingle{}$Q = 5$\textquotesingle{}}\NormalTok{)}
\NormalTok{axs[}\DecValTok{1}\NormalTok{].semilogx(frequencies, np.unwrap(np.angle(Hs\_lp\_3)) }\OperatorTok{*}\NormalTok{ (}\DecValTok{180} \OperatorTok{/}\NormalTok{ np.pi), label}\OperatorTok{=}\StringTok{\textquotesingle{}$Q = 5$\textquotesingle{}}\NormalTok{)}
\NormalTok{axs[}\DecValTok{0}\NormalTok{].semilogx(frequencies, }\DecValTok{20} \OperatorTok{*}\NormalTok{ np.log10(np.}\BuiltInTok{abs}\NormalTok{(Hs\_lp\_4)), label}\OperatorTok{=}\StringTok{\textquotesingle{}$Q = 10$\textquotesingle{}}\NormalTok{)}
\NormalTok{axs[}\DecValTok{1}\NormalTok{].semilogx(frequencies, np.unwrap(np.angle(Hs\_lp\_4)) }\OperatorTok{*}\NormalTok{ (}\DecValTok{180} \OperatorTok{/}\NormalTok{ np.pi), label}\OperatorTok{=}\StringTok{\textquotesingle{}$Q = 10$\textquotesingle{}}\NormalTok{)}
\NormalTok{axs[}\DecValTok{0}\NormalTok{].semilogx(frequencies, }\DecValTok{20} \OperatorTok{*}\NormalTok{ np.log10(np.}\BuiltInTok{abs}\NormalTok{(Hs\_lp\_5)), label}\OperatorTok{=}\StringTok{\textquotesingle{}$Q = 100$\textquotesingle{}}\NormalTok{)}
\NormalTok{axs[}\DecValTok{1}\NormalTok{].semilogx(frequencies, np.unwrap(np.angle(Hs\_lp\_5)) }\OperatorTok{*}\NormalTok{ (}\DecValTok{180} \OperatorTok{/}\NormalTok{ np.pi), label}\OperatorTok{=}\StringTok{\textquotesingle{}$Q = 100$\textquotesingle{}}\NormalTok{)}
\CommentTok{\textquotesingle{}\textquotesingle{}\textquotesingle{}}
\CommentTok{\# High Pass Filter}
\CommentTok{axs[0].semilogx(frequencies, 20 * np.log10(np.abs(Hs\_hp)), label=\textquotesingle{}high pass\textquotesingle{})}
\CommentTok{axs[1].semilogx(frequencies, np.unwrap(np.angle(Hs\_hp)) * (180 / np.pi), label=\textquotesingle{}high pass\textquotesingle{})}

\CommentTok{\# Band Pass Filter}
\CommentTok{axs[0].semilogx(frequencies, 20 * np.log10(np.abs(Hs\_bp)), label=\textquotesingle{}band pass\textquotesingle{})}
\CommentTok{axs[1].semilogx(frequencies, np.unwrap(np.angle(Hs\_bp)) * (180 / np.pi), label=\textquotesingle{}band pass\textquotesingle{})}

\CommentTok{\# Band Stop Filter}
\CommentTok{axs[0].semilogx(frequencies, 20 * np.log10(np.abs(Hs\_bs)), label=\textquotesingle{}band stop\textquotesingle{})}
\CommentTok{axs[1].semilogx(frequencies, (np.angle(Hs\_bs)) * (180 / np.pi), label=\textquotesingle{}band stop\textquotesingle{})}
\CommentTok{\textquotesingle{}\textquotesingle{}\textquotesingle{}}
\CommentTok{\#axs[0].title("amplitude response")}
\NormalTok{axs[}\DecValTok{0}\NormalTok{].set\_xlabel(}\StringTok{"frequency/Hz"}\NormalTok{)}
\NormalTok{axs[}\DecValTok{0}\NormalTok{].set\_ylabel(}\StringTok{"amplitude/dB"}\NormalTok{)}
\CommentTok{\#axs[0].set\_ylim({-}50, 25)}
\NormalTok{axs[}\DecValTok{0}\NormalTok{].grid(}\VariableTok{True}\NormalTok{, which}\OperatorTok{=}\StringTok{"both"}\NormalTok{, ls}\OperatorTok{=}\StringTok{"{-}{-}"}\NormalTok{)}
\NormalTok{axs[}\DecValTok{0}\NormalTok{].legend(loc}\OperatorTok{=}\DecValTok{1}\NormalTok{)}

\CommentTok{\#axs[1].title("phase response")}
\NormalTok{axs[}\DecValTok{1}\NormalTok{].set\_xlabel(}\StringTok{"frequency/Hz"}\NormalTok{)}
\NormalTok{axs[}\DecValTok{1}\NormalTok{].set\_ylabel(}\StringTok{"phase/deg"}\NormalTok{)}
\NormalTok{axs[}\DecValTok{1}\NormalTok{].grid(}\VariableTok{True}\NormalTok{, which}\OperatorTok{=}\StringTok{"both"}\NormalTok{, ls}\OperatorTok{=}\StringTok{"{-}{-}"}\NormalTok{)}
\NormalTok{axs[}\DecValTok{1}\NormalTok{].legend()}

\NormalTok{plt.tight\_layout()}
\CommentTok{\#plt.show()}

\NormalTok{plt.savefig(}\StringTok{\textquotesingle{}../images/sec\_theorectical\_background/lowPassDifferentQ.png\textquotesingle{}}\NormalTok{, }\BuiltInTok{format}\OperatorTok{=}\StringTok{\textquotesingle{}png\textquotesingle{}}\NormalTok{, dpi}\OperatorTok{=}\DecValTok{1000}\NormalTok{)}
\end{Highlighting}
\end{Shaded}

}

\end{codelisting}%

\begin{figure}

\centering{

\pandocbounded{\includegraphics[keepaspectratio]{content/../images/sec_theorectical_background/lowPassDifferentQ.png}}

}

\caption{\label{fig-lowPassDifferentQ}Low pass filter with different
quality factors}

\end{figure}%

The height of the peak can be calculated with:

\begin{equation}\phantomsection\label{eq-calculatePeak}{
A_{peak} = \frac{Q}{\sqrt{1 - \frac{1}{4Q^2}}}
}\end{equation}

For \(Q = 100\) the peak is \(A_{peak} = 100.001\), which converted into
dB is \(A_{peak,dB} = 40\,dB\), as is shown in
Figure~\ref{fig-lowPassDifferentQ}.

\section{Operational Amplifier}\label{operational-amplifier}

As seen in Figure~\ref{fig-BiquadCircuit}, a Biquad Filter consists of
four Operational Amplifiers. To get a better understanding of these,
this chapter will discuss the main aspects of OPAMPs. There are
different types of operational amplifiers that differ, for example, by
their low- or high-impedance inputs and outputs. According to Schmid,
there exist nine different types of the Opamps, but only four are mainly
used (Schmid 2000). This is because almost always, the non-inverting
(positive) input is designed as a high-impedance voltage input. The
inverting (negative) input can either be a high-impedance voltage input
or a low-impedance current input, depending on the type. Accordingly,
the output can be either a low-impedance voltage output or a
high-impedance current output. This results in four basic
configurations, as shown in the accompanying table
Table~\ref{tbl-opamp_compare}. (Images are taken from Wikipedia (2025))

\begin{longtable}[]{@{}
  >{\raggedright\arraybackslash}p{(\linewidth - 4\tabcolsep) * \real{0.1556}}
  >{\raggedright\arraybackslash}p{(\linewidth - 4\tabcolsep) * \real{0.4222}}
  >{\raggedright\arraybackslash}p{(\linewidth - 4\tabcolsep) * \real{0.4222}}@{}}
\caption{Four typicall Operational
Amplifiers}\label{tbl-opamp_compare}\tabularnewline
\toprule\noalign{}
\begin{minipage}[b]{\linewidth}\raggedright
\end{minipage} & \begin{minipage}[b]{\linewidth}\raggedright
Voltage output
\end{minipage} & \begin{minipage}[b]{\linewidth}\raggedright
Current output
\end{minipage} \\
\midrule\noalign{}
\endfirsthead
\toprule\noalign{}
\begin{minipage}[b]{\linewidth}\raggedright
\end{minipage} & \begin{minipage}[b]{\linewidth}\raggedright
Voltage output
\end{minipage} & \begin{minipage}[b]{\linewidth}\raggedright
Current output
\end{minipage} \\
\midrule\noalign{}
\endhead
\bottomrule\noalign{}
\endlastfoot
\textbf{Voltage input} & Voltage-Feedback Amplifiers & Operational
Transconductance Amplifier \\
&
\pandocbounded{\includegraphics[keepaspectratio]{content/../images/sec_theorectical_background/OpV_VV.png}}
&
\pandocbounded{\includegraphics[keepaspectratio]{content/../images/sec_theorectical_background/OpV_VC.png}} \\
\textbf{Current input} & Current-Feedback Amplifiers & Current
Amplifier \\
&
\pandocbounded{\includegraphics[keepaspectratio]{content/../images/sec_theorectical_background/OpV_CV.png}}
&
\pandocbounded{\includegraphics[keepaspectratio]{content/../images/sec_theorectical_background/OpV_CC.png}} \\
\end{longtable}

\begin{tcolorbox}[enhanced jigsaw, title=\textcolor{quarto-callout-tip-color}{\faLightbulb}\hspace{0.5em}{Tip}, colback=white, bottomrule=.15mm, coltitle=black, colbacktitle=quarto-callout-tip-color!10!white, toptitle=1mm, colframe=quarto-callout-tip-color-frame, rightrule=.15mm, arc=.35mm, leftrule=.75mm, bottomtitle=1mm, titlerule=0mm, toprule=.15mm, opacitybacktitle=0.6, opacityback=0, left=2mm, breakable]

As a general rule, the simplest circuit that can do a job is usually the
best choice. (H. Pretl and Michael Koefinger 2025)

\end{tcolorbox}

\subsection{Voltage-Feedback Amplifiers (VFA)}\label{sec-VFA}

When discussing operational amplifiers (OPAMPs), most sources refer to
the Voltage-Feedback Amplifier. These VFAs are
\textbf{voltage-controlled voltage sources}, essentially acting as
voltage boosters. They are characterized by a high-impedance input for
both the non-inverting and inverting terminals, and a low-impedance
voltage output. To realize various desired circuits, such as
amplification, integration, addition/subtraction, etc\ldots, these
functions should ideally be achieved only through the surrounding
circuitry. To meet this requirement, three main requirements need be
satisfied:

\begin{itemize}
\tightlist
\item
  \textbf{Extremely High Voltage Gain}: Typically ranging from 60 to 120
  dB (or gain factor of \(10^4\) to \(10^6\)), this gain should be
  available over a wide frequency range.
\item
  \textbf{High Impedance at Differential Inputs}: Ensuring minimal
  loading of the signal source.
\item
  \textbf{Low Impedance at the Output}: Allowing the amplifier to drive
  various loads without significant voltage drop.
\end{itemize}

\begin{tcolorbox}[enhanced jigsaw, title=\textcolor{quarto-callout-tip-color}{\faLightbulb}\hspace{0.5em}{Difference between a voltage source and a current source}, colback=white, bottomrule=.15mm, coltitle=black, colbacktitle=quarto-callout-tip-color!10!white, toptitle=1mm, colframe=quarto-callout-tip-color-frame, rightrule=.15mm, arc=.35mm, leftrule=.75mm, bottomtitle=1mm, titlerule=0mm, toprule=.15mm, opacitybacktitle=0.6, opacityback=0, left=2mm, breakable]

\textbf{Voltage source}\\
A voltage source creates a constant voltage output by changing the
current.

\textbf{Current source} A current source creates a constant current by
changing the voltage.

Both of this sources can be explained by Ohm's Law: \(R = \frac{U}{I}\).

For example the voltage source: There is no influence of the load, so
\(R\) may change and its value is unknown to the source. So to keep the
voltage stable we can only change the current.

\end{tcolorbox}

\subsection{Operational Transconductance Amplifier
(OTA)}\label{operational-transconductance-amplifier-ota}

In the design process of the biquad only OTAs will be used, so the focus
of this chapter will be on them. The operational transconductance
amplifier puts out out a current proportional to its input voltage,
unlike Voltage Feedback Amplifiers Section~\ref{sec-VFA}. In other
words, an OTA is a \textbf{voltage controlled current source}.

\begin{figure}

\centering{

\pandocbounded{\includegraphics[keepaspectratio]{content/../images/sec_theorectical_background/ota_symbol.png}}

}

\caption{\label{fig-otasym}Schematic symbol of an OTA (taken from
Wangenheim (2007))}

\end{figure}%

As seen in the figure Figure~\ref{fig-otasym}, the OTA has the two
differental inputs and a current output. On top of that it has a biasing
current input \(I_{bias}\) which controls the transconductance \(g_m\)
of the OTA.

\subsubsection{Transconductance}\label{transconductance}

Transconductance is a fundamental parameter that describes the
relationship between the input voltage and the output current in
electronic devices, particularly in transistors. It is defined as the
ratio of the change in output current to the change in input voltage,
under conditions where all other variables are held constant.
Mathematically, it is expressed as:

\[
g_m = \frac{\Delta I_{out}}{\Delta V_{in}}
\]

where \(g_m\) is the transconductance, \(\Delta I_{out}\) is the change
in output current, and \(\Delta V_{in}\) is the change in input voltage.
The unit of transconductance is the siemens (\(S\)), which is equivalent
to amperes per volt (\(A\)/\(V\)). Historically, it was also measured in
``mho'', which is ``ohm'' spelled backwards, reflecting its inverse
relationship to resistance.

In field-effect transistors (FETs) the transconductance determines the
device's ability to amplify signals. A higher transconductance value
indicates a stronger amplification capability, as a small change in
input voltage can result in a significant change in output current.

\begin{tcolorbox}[enhanced jigsaw, title=\textcolor{quarto-callout-note-color}{\faInfo}\hspace{0.5em}{Origin of the term transconductance}, colback=white, bottomrule=.15mm, coltitle=black, colbacktitle=quarto-callout-note-color!10!white, toptitle=1mm, colframe=quarto-callout-note-color-frame, rightrule=.15mm, arc=.35mm, leftrule=.75mm, bottomtitle=1mm, titlerule=0mm, toprule=.15mm, opacitybacktitle=0.6, opacityback=0, left=2mm, breakable]

The term transconductance originates from the concept of
\textbf{transfer conductance}. It combines the ideas of
\textbf{transfer}, indicating the transfer of a signal from the input to
the output, and \textbf{conductance}, which is the inverse of resistance
and measures how easily a material conducts electric current. In
essence, transconductance describes how effectively a device can convert
a voltage change at its input into a proportional current change at its
output. Kids (2025)

\end{tcolorbox}

\subsubsection{5-Transistor OTA}\label{transistor-ota}

H. Pretl and Michael Koefinger (2025) introduces a basic 5-Transistor
OTA as a design proposal for a real circuits.

\begin{figure}

\centering{

\pandocbounded{\includegraphics[keepaspectratio]{content/../images/sec_theorectical_background/basic_ota_output.png}}

}

\caption{\label{fig-5tota}H. Pretl and Michael Koefinger (2025) 's
5-Transistor OTA}

\end{figure}%

It consists of few elements which are going to be described in the next
chapters. The basic function of this OTA can be described with: The
transistors \(M_1\) and \(M_2\) form a differential pair, which is
biased by the current source \(M_5\). Transistors \(M_5\) and \(M_6\)
create a current mirror, where the input bias current
\(I_{\text{bias}}\) sets the bias current for the Operational
Transconductance Amplifier. To keep the biasing current stable, a
current mirror (Section~\ref{sec-current_mirror}) is used instead of
applying the \(I_{\text{bias}}\) directly to \(M_5\). The differential
pair (Section~\ref{sec-diff_pair}) \(M_1\) and \(M_2\) is loaded by the
current mirror \(M_3\) and \(M_4\), which mirrors the drain current of
\(M_1\) to the right side of the circuit. The combined currents from
\(M_4\) and \(M_2\) at the output node result in the generation of an
output current. To achieve this desired result, it is important to
ensure that \(M_{1,2}\) and \(M_{3,4}\) are symmetric, meaning they
should have identical \(W\) (width) and \(L\) (length) dimensions. To
set the right bias current, \(M_{5,6}\) should be chosen accordingly.

\subsection{Current Mirror}\label{sec-current_mirror}

A current mirror consists of at least two transistors, one working as
the input (reference) source and the other ones as copies of it. As the
name suggests, the primary function of a current mirror is to copy the
input current to the output, ensuring that the output current is an
exact duplicate of the input. Importantly, the flow of the input current
is unaffected by the size of the mirrored output current or any
variations in it.

\begin{figure}

\centering{

\includegraphics[width=0.4\linewidth,height=\textheight,keepaspectratio]{content/../images/sec_theorectical_background/current_mirror.png}

}

\caption{\label{fig-current_mirror}A current mirror (taken from Nerds Do
Stuff (2024))}

\end{figure}%

This behavior is made possible by two key characteristics of the current
mirror: its relatively low input resistance and its relatively high
output resistance. The low input resistance ensures the input current
remains stable regardless of drive conditions, and the high output
resistance maintains a constant output current independent of load
variations. (Analog University (2024))

A current mirror is also capable of copying one current and converting
it to multiple different outputs. This is possible by attaching more
transistors with their gate \(G\) to the drain \(D\) of the reference
transistor. By scaling the width and length of each output transistor
differently, multiple scales of the input current can be created.
Hereby, a big advantage is when using MOSFETs, due to the fact that no
current is flowing through the gate. In the case of BJTs, a compensation
circuit is added. (H. Pretl and Michael Koefinger (2025))

\subsection{Differential Pair}\label{sec-diff_pair}

A differential pair, also referred to as a differential amplifier, is an
electronic circuit designed to compare the difference between two input
voltages while minimizing any voltage that is common to both inputs. It
consists of two transistors and can have two inputs and two outputs, as
seen in Figure~\ref{fig-diff_pair}.

\begin{figure}

\centering{

\pandocbounded{\includegraphics[keepaspectratio]{content/../images/sec_theorectical_background/differential_pair.png}}

}

\caption{\label{fig-diff_pair}A differential pair (taken from H. Pretl
and Michael Koefinger (2025))}

\end{figure}%

Its primary function is to amplify the difference between two input
signals while rejecting any common-mode signals, which are identical in
both amplitude and phase at the two inputs. This property makes the
differential pair highly effective in noise reduction and signal
processing applications. In a typical differential pair using MOSFETs,
the two transistors are matched in terms of their electrical
characteristics, such as threshold voltage \(V_{th}\) and
transconductance \(g_m\). The inputs are applied to the gates \(G\) of
the MOSFETs, while the outputs are taken from the drains \(D\). A
current source is usually connected to the sources \(S\) of the MOSFETs
to provide a constant bias current, here noted as \(g_{\text{tail}}\) or
like in Figure~\ref{fig-5tota} as \(I_{\text{bias}}\).

The differential gain of the pair is determined by the transconductance
\(g_m\) of the MOSFETs. The common-mode rejection ratio (CMRR) is a
measure of the differential pair's ability to reject common-mode
signals. It is defined as the ratio of the differential gain to the
common-mode gain and is typically expressed in decibels (dB). The
performance of the differential pair is also influenced by the tail
current source, represented by \(g_{\text{tail}}\). The tail current
source provides a constant bias current \(I_{\text{bias}}\) that flows
through both transistors, ensuring that the total current through the
differential pair remains constant regardless of the input voltages.
(Razavi (2024))

\begin{tcolorbox}[enhanced jigsaw, title=\textcolor{quarto-callout-important-color}{\faExclamation}\hspace{0.5em}{Role in an OTA}, colback=white, bottomrule=.15mm, coltitle=black, colbacktitle=quarto-callout-important-color!10!white, toptitle=1mm, colframe=quarto-callout-important-color-frame, rightrule=.15mm, arc=.35mm, leftrule=.75mm, bottomtitle=1mm, titlerule=0mm, toprule=.15mm, opacitybacktitle=0.6, opacityback=0, left=2mm, breakable]

When used in OTAs, the differential pair serves as the input stage,
where the differential input voltage is converted into a differential
current.

\end{tcolorbox}

\section{MOSFET (Metal-Oxide-Semiconductor Field-Effect-Transistor
)}\label{mosfet-metal-oxide-semiconductor-field-effect-transistor}

To dive in one more level, the parts an operational amplifier is build
of are transistors. Precisely in this design process they are (MOSFETs).
In integrated circuit design, where pre-built components are not
available, MOSFETs can be designed from scratch. This design flexibility
allows for the adjustment of the MOSFET's physical dimensions,
specifically its width \(W\) and length \(L\). By doing this there is a
lot of room for freedom and experimental space. The length \(L\)
significantly influences MOSFET performance, enabling a trade-off
between speed and output conductance. The width \(W\) of a MOSFET acts
as a scaling factor that adjusts the charge density within the channel,
enabling the control of current levels to meet specific design
requirements.

\begin{tcolorbox}[enhanced jigsaw, title=\textcolor{quarto-callout-warning-color}{\faExclamationTriangle}\hspace{0.5em}{Circuit symbol of MOSFETs}, colback=white, bottomrule=.15mm, coltitle=black, colbacktitle=quarto-callout-warning-color!10!white, toptitle=1mm, colframe=quarto-callout-warning-color-frame, rightrule=.15mm, arc=.35mm, leftrule=.75mm, bottomtitle=1mm, titlerule=0mm, toprule=.15mm, opacitybacktitle=0.6, opacityback=0, left=2mm, breakable]

In the literature, the circuit symbol of MOSFETs is not uniform. Some
denote them with a small arrow between the gate (\(G\)) and source
(\(S\)), which is not very precise due to the fact MOSFETs are
symmetric. This means that the drain (\(D\)) and source (\(S\)) can be
interchanged, and the names are only defined to make circuits clearer.
Therefore, an accurate symbol for the n-MOS (shown in
Figure~\ref{fig-nmos}) and p-MOS (shown in Figure~\ref{fig-pmos}) is
shown in the following figures taken from H. Pretl and Michael Koefinger
(2025):

\begin{figure}[H]

\centering{

\pandocbounded{\includegraphics[keepaspectratio]{content/../images/sec_theorectical_background/nmos_symbol.png}}

}

\caption{\label{fig-nmos}n-MOSFET circuit symbol}

\end{figure}%

\begin{figure}[H]

\centering{

\pandocbounded{\includegraphics[keepaspectratio]{content/../images/sec_theorectical_background/pmos_symbol.png}}

}

\caption{\label{fig-pmos}p-MOSFET circuit symbol}

\end{figure}%

However, in this development report, the symbols are not strictly used
like shown here. But all of the used figures are based on this symmetric
model of a MOSFET.

Besides, in many situations the bulk \(B\) is connected to the source
\(S\) terminal. Thus the bulk \(B\) is not mentioned in the circuit
diagrams.

\end{tcolorbox}

\subsection{MOSFET Large-Signal
Regions}\label{mosfet-large-signal-regions}

A MOSFET can operate in three different modes: Saturation, Linear (or
Triode) region, or Cutoff. The modes depend on the device's threshold
voltage \(V_{\text{th}}\), gate-to-source voltage \(V_{\text{GS}}\), and
drain-to-source voltage \(V_{\text{DS}}\).

\textbf{Saturation}\\
The saturation region is the active mode used for analog amplification,
where the MOSFET behaves like a voltage-controlled current source. A
conductive channel forms between the source and drain. This state is
reached when the condition
\(V_{\text{DS}} \geq (V_{\text{GS}} - V_{\text{th}})\) is met. In
saturation, the drain-source current (\(I_D\)) stabilizes or
``saturates'' and becomes nearly independent of \(V_{\text{DS}}\).
However, it can still be precisely controlled via \(V_{\text{GS}}\) due
to the phenomenon known as
\href{https://en.wikipedia.org/wiki/Channel_length_modulation}{``pinch-off''}.

\textbf{Linear Region}\\
The Linear, or Triode, region occurs when a MOSFET is conducting and
acts more like a resistor or voltage-controlled device rather than a
current source. This region is defined by the condition
\(V_{\text{DS}} < (V_{\text{GS}} - V_{\text{th}})\) and
\(V_{\text{GS}} > V_{\text{th}}\). In this mode, both \(V_{\text{GS}}\)
and \(V_{\text{DS}}\) control the \(I_D\), allowing the device to be
used effectively for switching and amplification in low-voltage
applications.

\textbf{Cutoff Region}\\
In the cutoff region, the MOSFET is effectively turned off, with no
conductive channel formed between the drain and source. This mode is
reached when the gate-to-source voltage is less than the threshold
voltage (\(V_{\text{GS}} < V_{\text{th}}\)). In cutoff, the drain-source
current \(I_D\) is nearly zero, which means the MOSFET is not conducting
and behaves like an open switch.

\begin{tcolorbox}[enhanced jigsaw, title=\textcolor{quarto-callout-warning-color}{\faExclamationTriangle}\hspace{0.5em}{Warning}, colback=white, bottomrule=.15mm, coltitle=black, colbacktitle=quarto-callout-warning-color!10!white, toptitle=1mm, colframe=quarto-callout-warning-color-frame, rightrule=.15mm, arc=.35mm, leftrule=.75mm, bottomtitle=1mm, titlerule=0mm, toprule=.15mm, opacitybacktitle=0.6, opacityback=0, left=2mm, breakable]

In the large-signal analysis, the transitions between operating modes
occur gradually, making it challenging to pinpoint an exact moment when
one mode shifts to another. Consequently, the modes are used in a more
approximate manner, reflecting the continuous nature of these
transitions rather than distinct, sharply defined states.

\end{tcolorbox}

\subsection{Small-Signal
Representation}\label{small-signal-representation}

By applying different voltages to a MOSFET, different behaviours appear.
But there is one thing which is very unliked by engineers and others,
all these transfer functions are non linear. So making any calculations
by hand is nearly impossible. That's why in practise the biasing method
exists. This is done by applying a dc voltage to the terminals of the
MOSFET and during operation only a small signal change appears. With
this technique an \textbf{Operating Point} is created, which enables to
assume a linear behaviour in its small signal area (like in
Figure~\ref{fig-mosfetoppoint}). Because the MOSFET is put into an small
operation point and all the changes applied are very small, this
linearisation area is called \emph{Small-signal model}.

\begin{figure}

\centering{

\pandocbounded{\includegraphics[keepaspectratio]{content/../images/sec_theorectical_background/mosfet_oppoint.png}}

}

\caption{\label{fig-mosfetoppoint}MOSFET drain current
vs.~drain-to-source voltage for several values}

\end{figure}%

\begin{tcolorbox}[enhanced jigsaw, title=\textcolor{quarto-callout-important-color}{\faExclamation}\hspace{0.5em}{Important variables}, colback=white, bottomrule=.15mm, coltitle=black, colbacktitle=quarto-callout-important-color!10!white, toptitle=1mm, colframe=quarto-callout-important-color-frame, rightrule=.15mm, arc=.35mm, leftrule=.75mm, bottomtitle=1mm, titlerule=0mm, toprule=.15mm, opacitybacktitle=0.6, opacityback=0, left=2mm, breakable]

\begin{itemize}
\tightlist
\item
  \textbf{\(I_D\)}: The drain current
\item
  \textbf{\(I_{\text{bias}}\)}: The biasing current. The current to
  adjust the operating point
\item
  \textbf{\(V_{GS}\)}: The gate-to-source voltage
\item
  \textbf{\(V_{DS}\)}: The drain-to-source voltage
\item
  \textbf{\(V_{th}\)}: The threshold voltage. A critical parameter that
  defines the minimum gate-to-source voltage \(V_{GS}\) required to form
  a conductive channel between the source and drain
\item
  \textbf{\(V_{OV}\)}: The overdrive voltage. the difference between the
  gate-to-sourcevoltage \(V_{GS}\) and the threshold voltage \(V_{th}\).
\item
  \textbf{\(g_m\)}: The transconductance. The ratio between the output
  current and input voltage
\end{itemize}

\end{tcolorbox}

\begin{figure}

\centering{

\pandocbounded{\includegraphics[keepaspectratio]{content/../images/sec_theorectical_background/mosfet_smsi.png}}

}

\caption{\label{fig-mosfetsmsi}MOSFET-Small Signal Model}

\end{figure}%

\subsubsection{Capacitances in MOSFETs}\label{capacitances-in-mosfets}

\textbf{Gate-Source Capacitance (\(C_{gs}\)):}\\
\(C_{gs}\) is the capacitance between the gate and source terminals.
This arises due to the physical overlap of the gate and source regions,
as well as the insulating oxide layer between them. \(C_{gs}\) affects
both the input impedance and the switching speed. Higher \(C_{gs}\)
values can slow down the switching process because more charge needs to
be moved to change the gate voltage.

\textbf{Gate-Drain Capacitance (\(C_{gd}\)):}\\
\(C_{gd}\) is the capacitance between the gate and drain terminals.
People also know it as the Miller capacitance because of its significant
impact on the Miller effect. The
\href{https://en.wikipedia.org/wiki/Miller_effect}{Miller effect}
amplifies the effective value of \(C_{gd}\), making it a critical factor
in determining the high-frequency behaviour of the MOSFET. This
capacitance can cause phase shifts and affect the stability of the
circuit, particularly in amplifiers.

\textbf{Drain-Source Capacitance (\(C_{ds}\)):}\\
\(C_{ds}\) is the capacitance between the drain and source terminals. It
arises due to the depletion region between the drain and source regions.
\(C_{ds}\) affects the output impedance and frequency response of the
circuit. Higher drain-source capacitance can lead to increased output
capacitance, affecting the switching speed and introducing additional
losses.

\subsubsection{Transconductances in
MOSFET}\label{transconductances-in-mosfet}

\textbf{Output Conductance (\(g_{ds})\)}\\
\(g_{ds}\) represents the small-signal conductance between the drain and
source terminals of the MOSFET. It expresses how the drain current
\(I_D\) changes with respect to the drain-source voltage \(V_{ds}\) when
the gate-source voltage \(V_{gs}\) is held constant. Mathematically, it
is defined as follows:

\(g_{ds}=\frac{\partial I_D}{\partial V_{ds}}\bigg|_{V_{gs}=\text{constant}}\)

\(g_{ds}\) affects the output impedance of the MOSFET. A higher
\(g_{ds}\) results in a lower output impedance.In saturation region,
\(g_{ds}\) is typically small, which helps maintain the linearity of the
device. However, in the triode region, \(g_{ds}\) can be larger, leading
to non-linear behaviour.

\textbf{Mutual Transconductance - Gate-Source Voltage (\(g_mv_{gs}\))}\\
The transconductance \(g_m\) of a MOSFET quantifies how the drain
current \(I_D\) varies in response to changes in the gate-source voltage
\(V_{gs}\), while keeping the drain-source voltage \(V_{ds}\) constant:

\(g_m=\frac{\partial I_D}{\partial V_{gs}}\bigg|_{V_{ds}=\text{constant}}\)

The term \(g_mv_{gs}\) represents the small-signal current generated by
the transconductance \(g_m\) in response to a small change in the
gate-source voltage \(v_{gs}\). A higher \(g_m\) results in a higher
voltage gain.

\begin{tcolorbox}[enhanced jigsaw, title=\textcolor{quarto-callout-note-color}{\faInfo}\hspace{0.5em}{Summary of \(g_mv_{gs}\)}, colback=white, bottomrule=.15mm, coltitle=black, colbacktitle=quarto-callout-note-color!10!white, toptitle=1mm, colframe=quarto-callout-note-color-frame, rightrule=.15mm, arc=.35mm, leftrule=.75mm, bottomtitle=1mm, titlerule=0mm, toprule=.15mm, opacitybacktitle=0.6, opacityback=0, left=2mm, breakable]

In the saturation region, the output conductance \(g_{ds}\) is generally
much smaller than the transconductance \(g_m\). This characteristic
enables the MOSFET to function effectively as a voltage-controlled
current source. However, as the MOSFET moves into the triode region,
\(g_{ds}\) increases, causing the device to behave more like a resistor.
The output conductance \(g_{ds}\) influences the output impedance and
linearity of the device, while the product \(g_mv_{gs}\) controls the
gain and frequency response of the circuit.

\end{tcolorbox}

\bookmarksetup{startatroot}

\chapter{Characterisation}\label{characterisation}

This chapter is about characterising the biquad filter with the help of
computational software like python and and comparing that with ideal
circuit simulations. The comparison between a mathematical approach and
a basic circuit implementation, gives a first impression if the filter
is operational.

\section{Behauvioural Analysis and macro
modelling}\label{behauvioural-analysis-and-macro-modelling}

The behauvioural analysis is done through macro modelling the universal
biquad filter as a system. The system can be described with transfer
functions and modelled with python.

\subsection{Transfer Functions and frequency
response}\label{transfer-functions-and-frequency-response}

The ASLK PRO Manual (Rao and Ravikumar 2012) provides the transfer
functions of the four filter outputs: low pass, high pass, band pass,
and band stop. The transfer functions are adaptations of the general
second order transfer function as seen in
Equation~\ref{eq-TFSecondOrder}. (Razavi 2018)

In the following transfer function the input and output voltage are
referenced according to Figure~\ref{fig-BiquadCircuit}. The sections
only contain their specific transfer function and frequency response.

\textbf{Low pass}

The output if the low pass filter is marked in the circuit
(Figure~\ref{fig-BiquadCircuit}) as \(LPF\) and corresponds to
\(V_{03}\) in the transfer function Equation~\ref{eq-TFLowpass}.

\begin{equation}\phantomsection\label{eq-TFLowpass}{
\frac{V_{03}}{V_i} = \frac{H_0}{\left( 1 + \frac{s}{\omega_0 Q} + \frac{s^2}{\omega_0^2} \right)}
}\end{equation}

Figure~\ref{fig-freqResponseLowpass} shows the amplitude and phase
response of the low pass filter. The required frequency \(f_0 = 1\,kHz\)
and quality factor \(Q = 10\) recognizable in the bode plot. As the
dc-gain was chosen to be \(H_0 = 1\), the low pass filter has a
amplitude amplification of 1 in the lower frequencies.

\begin{codelisting}

\caption{\label{lst-freqresponselowpass}}

\centering{

\begin{Shaded}
\begin{Highlighting}[]
\CommentTok{\# Behavioral Analysis Biquad Filter}

\ImportTok{import}\NormalTok{ numpy }\ImportTok{as}\NormalTok{ np}
\ImportTok{import}\NormalTok{ matplotlib.pyplot }\ImportTok{as}\NormalTok{ plt}

\CommentTok{\# Initial values}
\NormalTok{f0 }\OperatorTok{=} \FloatTok{1e3}  \CommentTok{\# Resonance frequency in Hz}
\NormalTok{w0 }\OperatorTok{=} \DecValTok{2} \OperatorTok{*}\NormalTok{ np.pi }\OperatorTok{*}\NormalTok{ f0  }\CommentTok{\# Angular frequency in rad/s}
\NormalTok{Q }\OperatorTok{=} \DecValTok{10}  \CommentTok{\# Quality factor}
\NormalTok{H0 }\OperatorTok{=} \DecValTok{1}  \CommentTok{\# Play around with this later}

\CommentTok{\# Logarithmic frequency axis}
\NormalTok{frequencies }\OperatorTok{=}\NormalTok{ np.logspace(}\DecValTok{2}\NormalTok{, }\DecValTok{4}\NormalTok{, }\DecValTok{10000}\NormalTok{)  }\CommentTok{\# Frequency from 10\^{}2 to 10\^{}4 Hz}
\NormalTok{s }\OperatorTok{=} \OtherTok{1j} \OperatorTok{*} \DecValTok{2} \OperatorTok{*}\NormalTok{ np.pi }\OperatorTok{*}\NormalTok{ frequencies  }\CommentTok{\# Laplace{-}Variable s = jω}

\CommentTok{\#\#\#\#\#\#\#\#\#\#\#\#\#\#\#\#\#\#\#\#\#\#\#\#\#\#\#\#\#\#\#\#\#\#\#\#\#\#\#\#\#\#\#\#}
\CommentTok{\# Transfer functions of Active Filters}
\CommentTok{\#\#\#\#\#\#\#\#\#\#\#\#\#\#\#\#\#\#\#\#\#\#\#\#\#\#\#\#\#\#\#\#\#\#\#\#\#\#\#\#\#\#\#\#}

\CommentTok{\#\#\# Numerator}
\CommentTok{\# Low Pass Filter}
\NormalTok{b\_lp }\OperatorTok{=}\NormalTok{ H0}

\CommentTok{\# High Pass Filter}
\NormalTok{b\_hp }\OperatorTok{=}\NormalTok{ (H0 }\OperatorTok{*}\NormalTok{ (s}\OperatorTok{**}\DecValTok{2} \OperatorTok{/}\NormalTok{ w0}\OperatorTok{**}\DecValTok{2}\NormalTok{))}

\CommentTok{\# Band Pass Filter}
\NormalTok{b\_bp }\OperatorTok{=}\NormalTok{ (}\OperatorTok{{-}}\NormalTok{H0 }\OperatorTok{*}\NormalTok{ (s }\OperatorTok{/}\NormalTok{ w0))}

\CommentTok{\# Band Stop Filter}
\NormalTok{b\_bs }\OperatorTok{=} \OperatorTok{{-}}\NormalTok{((}\DecValTok{1} \OperatorTok{+}\NormalTok{ (s}\OperatorTok{**}\DecValTok{2} \OperatorTok{/}\NormalTok{ (w0}\OperatorTok{**}\DecValTok{2}\NormalTok{))) }\OperatorTok{*}\NormalTok{ H0)}

\CommentTok{\# Denominator {-}\textgreater{} for all filters the same}
\NormalTok{a0 }\OperatorTok{=} \DecValTok{1}
\NormalTok{a1 }\OperatorTok{=}\NormalTok{ (s }\OperatorTok{/}\NormalTok{ (w0 }\OperatorTok{*}\NormalTok{ Q))}
\NormalTok{a2 }\OperatorTok{=}\NormalTok{ (s}\OperatorTok{**}\DecValTok{2} \OperatorTok{/}\NormalTok{ (w0}\OperatorTok{**}\DecValTok{2}\NormalTok{))}

\NormalTok{den }\OperatorTok{=}\NormalTok{ a0 }\OperatorTok{+}\NormalTok{ a1 }\OperatorTok{+}\NormalTok{ a2}

\CommentTok{\#\#\#\#\#\#\#\#\#\#\#\#\#\#\#\#\#\#\#\#\#\#\#\#\#\#\#\#\#\#\#\#\#\#\#\#\#\#\#\#\#\#\#\#}
\CommentTok{\# Calculation of the transfer functions H(s)}
\CommentTok{\#\#\#\#\#\#\#\#\#\#\#\#\#\#\#\#\#\#\#\#\#\#\#\#\#\#\#\#\#\#\#\#\#\#\#\#\#\#\#\#\#\#\#\#}
\NormalTok{Hs\_lp }\OperatorTok{=}\NormalTok{ b\_lp }\OperatorTok{/}\NormalTok{ den}
\NormalTok{Hs\_hp }\OperatorTok{=}\NormalTok{ b\_hp }\OperatorTok{/}\NormalTok{ den}
\NormalTok{Hs\_bp }\OperatorTok{=}\NormalTok{ b\_bp }\OperatorTok{/}\NormalTok{ den}
\NormalTok{Hs\_bs }\OperatorTok{=}\NormalTok{ b\_bs }\OperatorTok{/}\NormalTok{ den}

\CommentTok{\# Bode Diagram}
\NormalTok{fig, axs }\OperatorTok{=}\NormalTok{ plt.subplots(}\DecValTok{2}\NormalTok{)}
\CommentTok{\#fig.suptitle("frequency response of biquad filter")}

\CommentTok{\# Low Pass Filter}
\NormalTok{axs[}\DecValTok{0}\NormalTok{].semilogx(frequencies, }\DecValTok{20} \OperatorTok{*}\NormalTok{ np.log10(np.}\BuiltInTok{abs}\NormalTok{(Hs\_lp)), label}\OperatorTok{=}\StringTok{\textquotesingle{}low pass\textquotesingle{}}\NormalTok{)}
\NormalTok{axs[}\DecValTok{1}\NormalTok{].semilogx(frequencies, np.unwrap(np.angle(Hs\_lp)) }\OperatorTok{*}\NormalTok{ (}\DecValTok{180} \OperatorTok{/}\NormalTok{ np.pi), label}\OperatorTok{=}\StringTok{\textquotesingle{}low pass\textquotesingle{}}\NormalTok{)}
\CommentTok{\textquotesingle{}\textquotesingle{}\textquotesingle{}}
\CommentTok{\# High Pass Filter}
\CommentTok{axs[0].semilogx(frequencies, 20 * np.log10(np.abs(Hs\_hp)), label=\textquotesingle{}high pass\textquotesingle{})}
\CommentTok{axs[1].semilogx(frequencies, np.unwrap(np.angle(Hs\_hp)) * (180 / np.pi), label=\textquotesingle{}high pass\textquotesingle{})}

\CommentTok{\# Band Pass Filter}
\CommentTok{axs[0].semilogx(frequencies, 20 * np.log10(np.abs(Hs\_bp)), label=\textquotesingle{}band pass\textquotesingle{})}
\CommentTok{axs[1].semilogx(frequencies, np.unwrap(np.angle(Hs\_bp)) * (180 / np.pi), label=\textquotesingle{}band pass\textquotesingle{})}

\CommentTok{\# Band Stop Filter}
\CommentTok{axs[0].semilogx(frequencies, 20 * np.log10(np.abs(Hs\_bs)), label=\textquotesingle{}band stop\textquotesingle{})}
\CommentTok{axs[1].semilogx(frequencies, (np.angle(Hs\_bs)) * (180 / np.pi), label=\textquotesingle{}band stop\textquotesingle{})}
\CommentTok{\textquotesingle{}\textquotesingle{}\textquotesingle{}}
\CommentTok{\#axs[0].title("amplitude response")}
\NormalTok{axs[}\DecValTok{0}\NormalTok{].set\_xlabel(}\StringTok{"frequency/Hz"}\NormalTok{)}
\NormalTok{axs[}\DecValTok{0}\NormalTok{].set\_ylabel(}\StringTok{"amplitude/dB"}\NormalTok{)}
\NormalTok{axs[}\DecValTok{0}\NormalTok{].set\_ylim(}\OperatorTok{{-}}\DecValTok{50}\NormalTok{, }\DecValTok{25}\NormalTok{)}
\NormalTok{axs[}\DecValTok{0}\NormalTok{].grid(}\VariableTok{True}\NormalTok{, which}\OperatorTok{=}\StringTok{"both"}\NormalTok{, ls}\OperatorTok{=}\StringTok{"{-}{-}"}\NormalTok{)}
\CommentTok{\#axs[0].legend(loc=1)}

\CommentTok{\#axs[1].title("phase response")}
\NormalTok{axs[}\DecValTok{1}\NormalTok{].set\_xlabel(}\StringTok{"frequency/Hz"}\NormalTok{)}
\NormalTok{axs[}\DecValTok{1}\NormalTok{].set\_ylabel(}\StringTok{"phase/deg"}\NormalTok{)}
\NormalTok{axs[}\DecValTok{1}\NormalTok{].grid(}\VariableTok{True}\NormalTok{, which}\OperatorTok{=}\StringTok{"both"}\NormalTok{, ls}\OperatorTok{=}\StringTok{"{-}{-}"}\NormalTok{)}
\CommentTok{\#axs[1].legend()}

\NormalTok{plt.tight\_layout()}
\CommentTok{\#plt.show()}

\NormalTok{plt.savefig(}\StringTok{"../images/sec\_characterisation/freqResponseLowpass.png"}\NormalTok{, }\BuiltInTok{format}\OperatorTok{=}\StringTok{"png"}\NormalTok{,dpi}\OperatorTok{=}\DecValTok{1000}\NormalTok{)}
\end{Highlighting}
\end{Shaded}

}

\end{codelisting}%

\begin{figure}

\centering{

\pandocbounded{\includegraphics[keepaspectratio]{content/../images/sec_characterisation/freqResponseLowpass.png}}

}

\caption{\label{fig-freqResponseLowpass}Frequency response of the low
pass filter}

\end{figure}%

With knowing the dc gain \(H_0 = 1\) and quality factot \(Q = 10\), the
amplitude of the peak can be calculated as seen in
Equation~\ref{eq-calculatePeak}. Expression the value in dB, gives peak
amplitude of \(A_{peak,dB} = 20.002\,dB\) which corresponds to the peak
value seen in Figure~\ref{fig-freqResponseLowpass}.

\textbf{High pass}

The output if the high pass filter is marked in the circuit
(Figure~\ref{fig-BiquadCircuit}) as \(HPF\) and corresponds to
\(V_{01}\) in the transfer function Equation~\ref{eq-TFHighpass}.

\begin{equation}\phantomsection\label{eq-TFHighpass}{
\frac{V_{01}}{V_i} = \frac{\left( H_0 \cdot \frac{s^2}{\omega_0^2} \right)}{\left( 1 + \frac{s}{\omega_0 Q} + \frac{s^2}{\omega_0^2} \right)} 
}\end{equation}

Figure~\ref{fig-freqResponseHighpass} shows the amplitude and phase
response of the high pass filter. The required frequency
\(f_0 = 1\,kHz\) and quality factor \(Q = 10\) recognizable in the bode
plot. As the dc-gain was chosen to be \(H_0 = 1\), the low pass filter
has a amplitude amplification of 1 in the higher frequencies.

\begin{codelisting}

\caption{\label{lst-freqresponsehighpass}}

\centering{

\begin{Shaded}
\begin{Highlighting}[]
\CommentTok{\# Behavioral Analysis Biquad Filter}

\ImportTok{import}\NormalTok{ numpy }\ImportTok{as}\NormalTok{ np}
\ImportTok{import}\NormalTok{ matplotlib.pyplot }\ImportTok{as}\NormalTok{ plt}

\CommentTok{\# Initial values}
\NormalTok{f0 }\OperatorTok{=} \FloatTok{1e3}  \CommentTok{\# Resonance frequency in Hz}
\NormalTok{w0 }\OperatorTok{=} \DecValTok{2} \OperatorTok{*}\NormalTok{ np.pi }\OperatorTok{*}\NormalTok{ f0  }\CommentTok{\# Angular frequency in rad/s}
\NormalTok{Q }\OperatorTok{=} \DecValTok{10}  \CommentTok{\# Quality factor}
\NormalTok{H0 }\OperatorTok{=} \DecValTok{1}  \CommentTok{\# Play around with this later}

\CommentTok{\# Logarithmic frequency axis}
\NormalTok{frequencies }\OperatorTok{=}\NormalTok{ np.logspace(}\DecValTok{2}\NormalTok{, }\DecValTok{4}\NormalTok{, }\DecValTok{10000}\NormalTok{)  }\CommentTok{\# Frequency from 10\^{}2 to 10\^{}4 Hz}
\NormalTok{s }\OperatorTok{=} \OtherTok{1j} \OperatorTok{*} \DecValTok{2} \OperatorTok{*}\NormalTok{ np.pi }\OperatorTok{*}\NormalTok{ frequencies  }\CommentTok{\# Laplace{-}Variable s = jω}

\CommentTok{\#\#\#\#\#\#\#\#\#\#\#\#\#\#\#\#\#\#\#\#\#\#\#\#\#\#\#\#\#\#\#\#\#\#\#\#\#\#\#\#\#\#\#\#}
\CommentTok{\# Transfer functions of Active Filters}
\CommentTok{\#\#\#\#\#\#\#\#\#\#\#\#\#\#\#\#\#\#\#\#\#\#\#\#\#\#\#\#\#\#\#\#\#\#\#\#\#\#\#\#\#\#\#\#}

\CommentTok{\#\#\# Numerator}
\CommentTok{\# Low Pass Filter}
\NormalTok{b\_lp }\OperatorTok{=}\NormalTok{ H0}

\CommentTok{\# High Pass Filter}
\NormalTok{b\_hp }\OperatorTok{=}\NormalTok{ (H0 }\OperatorTok{*}\NormalTok{ (s}\OperatorTok{**}\DecValTok{2} \OperatorTok{/}\NormalTok{ w0}\OperatorTok{**}\DecValTok{2}\NormalTok{))}

\CommentTok{\# Band Pass Filter}
\NormalTok{b\_bp }\OperatorTok{=}\NormalTok{ (}\OperatorTok{{-}}\NormalTok{H0 }\OperatorTok{*}\NormalTok{ (s }\OperatorTok{/}\NormalTok{ w0))}

\CommentTok{\# Band Stop Filter}
\NormalTok{b\_bs }\OperatorTok{=} \OperatorTok{{-}}\NormalTok{((}\DecValTok{1} \OperatorTok{+}\NormalTok{ (s}\OperatorTok{**}\DecValTok{2} \OperatorTok{/}\NormalTok{ (w0}\OperatorTok{**}\DecValTok{2}\NormalTok{))) }\OperatorTok{*}\NormalTok{ H0)}

\CommentTok{\# Denominator {-}\textgreater{} for all filters the same}
\NormalTok{a0 }\OperatorTok{=} \DecValTok{1}
\NormalTok{a1 }\OperatorTok{=}\NormalTok{ (s }\OperatorTok{/}\NormalTok{ (w0 }\OperatorTok{*}\NormalTok{ Q))}
\NormalTok{a2 }\OperatorTok{=}\NormalTok{ (s}\OperatorTok{**}\DecValTok{2} \OperatorTok{/}\NormalTok{ (w0}\OperatorTok{**}\DecValTok{2}\NormalTok{))}

\NormalTok{den }\OperatorTok{=}\NormalTok{ a0 }\OperatorTok{+}\NormalTok{ a1 }\OperatorTok{+}\NormalTok{ a2}

\CommentTok{\#\#\#\#\#\#\#\#\#\#\#\#\#\#\#\#\#\#\#\#\#\#\#\#\#\#\#\#\#\#\#\#\#\#\#\#\#\#\#\#\#\#\#\#}
\CommentTok{\# Calculation of the transfer functions H(s)}
\CommentTok{\#\#\#\#\#\#\#\#\#\#\#\#\#\#\#\#\#\#\#\#\#\#\#\#\#\#\#\#\#\#\#\#\#\#\#\#\#\#\#\#\#\#\#\#}
\NormalTok{Hs\_lp }\OperatorTok{=}\NormalTok{ b\_lp }\OperatorTok{/}\NormalTok{ den}
\NormalTok{Hs\_hp }\OperatorTok{=}\NormalTok{ b\_hp }\OperatorTok{/}\NormalTok{ den}
\NormalTok{Hs\_bp }\OperatorTok{=}\NormalTok{ b\_bp }\OperatorTok{/}\NormalTok{ den}
\NormalTok{Hs\_bs }\OperatorTok{=}\NormalTok{ b\_bs }\OperatorTok{/}\NormalTok{ den}

\CommentTok{\# Bode Diagram}
\NormalTok{fig, axs }\OperatorTok{=}\NormalTok{ plt.subplots(}\DecValTok{2}\NormalTok{)}
\CommentTok{\#fig.suptitle("frequency response of biquad filter")}

\CommentTok{\# Low Pass Filter}
\CommentTok{\#axs[0].semilogx(frequencies, 20 * np.log10(np.abs(Hs\_lp)), label=\textquotesingle{}low pass\textquotesingle{})}
\CommentTok{\#axs[1].semilogx(frequencies, np.unwrap(np.angle(Hs\_lp)) * (180 / np.pi), label=\textquotesingle{}low pass\textquotesingle{})}

\CommentTok{\# High Pass Filter}
\NormalTok{axs[}\DecValTok{0}\NormalTok{].semilogx(frequencies, }\DecValTok{20} \OperatorTok{*}\NormalTok{ np.log10(np.}\BuiltInTok{abs}\NormalTok{(Hs\_hp)), label}\OperatorTok{=}\StringTok{\textquotesingle{}high pass\textquotesingle{}}\NormalTok{)}
\NormalTok{axs[}\DecValTok{1}\NormalTok{].semilogx(frequencies, np.unwrap(np.angle(Hs\_hp)) }\OperatorTok{*}\NormalTok{ (}\DecValTok{180} \OperatorTok{/}\NormalTok{ np.pi), label}\OperatorTok{=}\StringTok{\textquotesingle{}high pass\textquotesingle{}}\NormalTok{)}
\CommentTok{\textquotesingle{}\textquotesingle{}\textquotesingle{}}
\CommentTok{\# Band Pass Filter}
\CommentTok{axs[0].semilogx(frequencies, 20 * np.log10(np.abs(Hs\_bp)), label=\textquotesingle{}band pass\textquotesingle{})}
\CommentTok{axs[1].semilogx(frequencies, np.unwrap(np.angle(Hs\_bp)) * (180 / np.pi), label=\textquotesingle{}band pass\textquotesingle{})}

\CommentTok{\# Band Stop Filter}
\CommentTok{axs[0].semilogx(frequencies, 20 * np.log10(np.abs(Hs\_bs)), label=\textquotesingle{}band stop\textquotesingle{})}
\CommentTok{axs[1].semilogx(frequencies, (np.angle(Hs\_bs)) * (180 / np.pi), label=\textquotesingle{}band stop\textquotesingle{})}
\CommentTok{\textquotesingle{}\textquotesingle{}\textquotesingle{}}
\CommentTok{\#axs[0].title("amplitude response")}
\NormalTok{axs[}\DecValTok{0}\NormalTok{].set\_xlabel(}\StringTok{"frequency/Hz"}\NormalTok{)}
\NormalTok{axs[}\DecValTok{0}\NormalTok{].set\_ylabel(}\StringTok{"amplitude/dB"}\NormalTok{)}
\NormalTok{axs[}\DecValTok{0}\NormalTok{].set\_ylim(}\OperatorTok{{-}}\DecValTok{50}\NormalTok{, }\DecValTok{25}\NormalTok{)}
\NormalTok{axs[}\DecValTok{0}\NormalTok{].grid(}\VariableTok{True}\NormalTok{, which}\OperatorTok{=}\StringTok{"both"}\NormalTok{, ls}\OperatorTok{=}\StringTok{"{-}{-}"}\NormalTok{)}
\CommentTok{\#axs[0].legend(loc=1)}

\CommentTok{\#axs[1].title("phase response")}
\NormalTok{axs[}\DecValTok{1}\NormalTok{].set\_xlabel(}\StringTok{"frequency/Hz"}\NormalTok{)}
\NormalTok{axs[}\DecValTok{1}\NormalTok{].set\_ylabel(}\StringTok{"phase/deg"}\NormalTok{)}
\NormalTok{axs[}\DecValTok{1}\NormalTok{].grid(}\VariableTok{True}\NormalTok{, which}\OperatorTok{=}\StringTok{"both"}\NormalTok{, ls}\OperatorTok{=}\StringTok{"{-}{-}"}\NormalTok{)}
\CommentTok{\#axs[1].legend()}

\NormalTok{plt.tight\_layout()}
\CommentTok{\#plt.show()}

\NormalTok{plt.savefig(}\StringTok{"../images/sec\_characterisation/freqResponseHighpass.png"}\NormalTok{, }\BuiltInTok{format}\OperatorTok{=}\StringTok{"png"}\NormalTok{, dpi}\OperatorTok{=}\DecValTok{1000}\NormalTok{)}
\end{Highlighting}
\end{Shaded}

}

\end{codelisting}%

\begin{figure}

\centering{

\pandocbounded{\includegraphics[keepaspectratio]{content/../images/sec_characterisation/freqResponseHighpass.png}}

}

\caption{\label{fig-freqResponseHighpass}Frequency response of the high
pass filter}

\end{figure}%

\textbf{Band pass}

The output for the band pass filter is marked as \(BPF\) in
Figure~\ref{fig-BiquadCircuit}. This denotes the point that is
referenced in Equation~\ref{eq-TFBandpass} as \(V_{02}\).

\begin{equation}\phantomsection\label{eq-TFBandpass}{
\frac{V_{02}}{V_i} = \frac{\left( - H_0 \cdot \frac{s}{\omega_0} \right)}{\left( 1 + \frac{s}{\omega_0 Q} + \frac{s^2}{\omega_0^2} \right)}
}\end{equation}

The band pass shown in Figure~\ref{fig-freqResponseBandpass} has its
center frequency at \(1\,kHz\) as set in the requirements. Similarly to
the low pass filter in Figure~\ref{fig-freqResponseLowpass} and the high
pass filter in Figure~\ref{fig-freqResponseHighpass} the amplitude
response peaks at this frequency, with its peak influenced by the
quality factor.

\begin{codelisting}

\caption{\label{lst-freqresponsebandpass}}

\centering{

\begin{Shaded}
\begin{Highlighting}[]
\CommentTok{\# Behavioral Analysis Biquad Filter}

\ImportTok{import}\NormalTok{ numpy }\ImportTok{as}\NormalTok{ np}
\ImportTok{import}\NormalTok{ matplotlib.pyplot }\ImportTok{as}\NormalTok{ plt}

\CommentTok{\# Initial values}
\NormalTok{f0 }\OperatorTok{=} \FloatTok{1e3}  \CommentTok{\# Resonance frequency in Hz}
\NormalTok{w0 }\OperatorTok{=} \DecValTok{2} \OperatorTok{*}\NormalTok{ np.pi }\OperatorTok{*}\NormalTok{ f0  }\CommentTok{\# Angular frequency in rad/s}
\NormalTok{Q }\OperatorTok{=} \DecValTok{10}  \CommentTok{\# Quality factor}
\NormalTok{H0 }\OperatorTok{=} \DecValTok{1}  \CommentTok{\# Play around with this later}

\CommentTok{\# Logarithmic frequency axis}
\NormalTok{frequencies }\OperatorTok{=}\NormalTok{ np.logspace(}\DecValTok{2}\NormalTok{, }\DecValTok{4}\NormalTok{, }\DecValTok{10000}\NormalTok{)  }\CommentTok{\# Frequency from 10\^{}2 to 10\^{}4 Hz}
\NormalTok{s }\OperatorTok{=} \OtherTok{1j} \OperatorTok{*} \DecValTok{2} \OperatorTok{*}\NormalTok{ np.pi }\OperatorTok{*}\NormalTok{ frequencies  }\CommentTok{\# Laplace{-}Variable s = jω}

\CommentTok{\#\#\#\#\#\#\#\#\#\#\#\#\#\#\#\#\#\#\#\#\#\#\#\#\#\#\#\#\#\#\#\#\#\#\#\#\#\#\#\#\#\#\#\#}
\CommentTok{\# Transfer functions of Active Filters}
\CommentTok{\#\#\#\#\#\#\#\#\#\#\#\#\#\#\#\#\#\#\#\#\#\#\#\#\#\#\#\#\#\#\#\#\#\#\#\#\#\#\#\#\#\#\#\#}

\CommentTok{\#\#\# Numerator}
\CommentTok{\# Low Pass Filter}
\NormalTok{b\_lp }\OperatorTok{=}\NormalTok{ H0}

\CommentTok{\# High Pass Filter}
\NormalTok{b\_hp }\OperatorTok{=}\NormalTok{ (H0 }\OperatorTok{*}\NormalTok{ (s}\OperatorTok{**}\DecValTok{2} \OperatorTok{/}\NormalTok{ w0}\OperatorTok{**}\DecValTok{2}\NormalTok{))}

\CommentTok{\# Band Pass Filter}
\NormalTok{b\_bp }\OperatorTok{=}\NormalTok{ (}\OperatorTok{{-}}\NormalTok{H0 }\OperatorTok{*}\NormalTok{ (s }\OperatorTok{/}\NormalTok{ w0))}

\CommentTok{\# Band Stop Filter}
\NormalTok{b\_bs }\OperatorTok{=} \OperatorTok{{-}}\NormalTok{((}\DecValTok{1} \OperatorTok{+}\NormalTok{ (s}\OperatorTok{**}\DecValTok{2} \OperatorTok{/}\NormalTok{ (w0}\OperatorTok{**}\DecValTok{2}\NormalTok{))) }\OperatorTok{*}\NormalTok{ H0)}

\CommentTok{\# Denominator {-}\textgreater{} for all filters the same}
\NormalTok{a0 }\OperatorTok{=} \DecValTok{1}
\NormalTok{a1 }\OperatorTok{=}\NormalTok{ (s }\OperatorTok{/}\NormalTok{ (w0 }\OperatorTok{*}\NormalTok{ Q))}
\NormalTok{a2 }\OperatorTok{=}\NormalTok{ (s}\OperatorTok{**}\DecValTok{2} \OperatorTok{/}\NormalTok{ (w0}\OperatorTok{**}\DecValTok{2}\NormalTok{))}

\NormalTok{den }\OperatorTok{=}\NormalTok{ a0 }\OperatorTok{+}\NormalTok{ a1 }\OperatorTok{+}\NormalTok{ a2}

\CommentTok{\#\#\#\#\#\#\#\#\#\#\#\#\#\#\#\#\#\#\#\#\#\#\#\#\#\#\#\#\#\#\#\#\#\#\#\#\#\#\#\#\#\#\#\#}
\CommentTok{\# Calculation of the transfer functions H(s)}
\CommentTok{\#\#\#\#\#\#\#\#\#\#\#\#\#\#\#\#\#\#\#\#\#\#\#\#\#\#\#\#\#\#\#\#\#\#\#\#\#\#\#\#\#\#\#\#}
\NormalTok{Hs\_lp }\OperatorTok{=}\NormalTok{ b\_lp }\OperatorTok{/}\NormalTok{ den}
\NormalTok{Hs\_hp }\OperatorTok{=}\NormalTok{ b\_hp }\OperatorTok{/}\NormalTok{ den}
\NormalTok{Hs\_bp }\OperatorTok{=}\NormalTok{ b\_bp }\OperatorTok{/}\NormalTok{ den}
\NormalTok{Hs\_bs }\OperatorTok{=}\NormalTok{ b\_bs }\OperatorTok{/}\NormalTok{ den}

\CommentTok{\# Bode Diagram}
\NormalTok{fig, axs }\OperatorTok{=}\NormalTok{ plt.subplots(}\DecValTok{2}\NormalTok{)}
\CommentTok{\#fig.suptitle("frequency response of biquad filter")}
\CommentTok{\textquotesingle{}\textquotesingle{}\textquotesingle{}}
\CommentTok{\# Low Pass Filter}
\CommentTok{axs[0].semilogx(frequencies, 20 * np.log10(np.abs(Hs\_lp)), label=\textquotesingle{}low pass\textquotesingle{})}
\CommentTok{axs[1].semilogx(frequencies, np.unwrap(np.angle(Hs\_lp)) * (180 / np.pi), label=\textquotesingle{}low pass\textquotesingle{})}

\CommentTok{\# High Pass Filter}
\CommentTok{axs[0].semilogx(frequencies, 20 * np.log10(np.abs(Hs\_hp)), label=\textquotesingle{}high pass\textquotesingle{})}
\CommentTok{axs[1].semilogx(frequencies, np.unwrap(np.angle(Hs\_hp)) * (180 / np.pi), label=\textquotesingle{}high pass\textquotesingle{})}
\CommentTok{\textquotesingle{}\textquotesingle{}\textquotesingle{}}
\CommentTok{\# Band Pass Filter}
\NormalTok{axs[}\DecValTok{0}\NormalTok{].semilogx(frequencies, }\DecValTok{20} \OperatorTok{*}\NormalTok{ np.log10(np.}\BuiltInTok{abs}\NormalTok{(Hs\_bp)), label}\OperatorTok{=}\StringTok{\textquotesingle{}band pass\textquotesingle{}}\NormalTok{)}
\NormalTok{axs[}\DecValTok{1}\NormalTok{].semilogx(frequencies, np.unwrap(np.angle(Hs\_bp)) }\OperatorTok{*}\NormalTok{ (}\DecValTok{180} \OperatorTok{/}\NormalTok{ np.pi), label}\OperatorTok{=}\StringTok{\textquotesingle{}band pass\textquotesingle{}}\NormalTok{)}

\CommentTok{\# Band Stop Filter}
\CommentTok{\#axs[0].semilogx(frequencies, 20 * np.log10(np.abs(Hs\_bs)), label=\textquotesingle{}band stop\textquotesingle{})}
\CommentTok{\#axs[1].semilogx(frequencies, (np.angle(Hs\_bs)) * (180 / np.pi), label=\textquotesingle{}band stop\textquotesingle{})}

\CommentTok{\#axs[0].title("amplitude response")}
\NormalTok{axs[}\DecValTok{0}\NormalTok{].set\_xlabel(}\StringTok{"frequency/Hz"}\NormalTok{)}
\NormalTok{axs[}\DecValTok{0}\NormalTok{].set\_ylabel(}\StringTok{"amplitude/dB"}\NormalTok{)}
\NormalTok{axs[}\DecValTok{0}\NormalTok{].set\_ylim(}\OperatorTok{{-}}\DecValTok{50}\NormalTok{, }\DecValTok{25}\NormalTok{)}
\NormalTok{axs[}\DecValTok{0}\NormalTok{].grid(}\VariableTok{True}\NormalTok{, which}\OperatorTok{=}\StringTok{"both"}\NormalTok{, ls}\OperatorTok{=}\StringTok{"{-}{-}"}\NormalTok{)}
\CommentTok{\#axs[0].legend(loc=1)}

\CommentTok{\#axs[1].title("phase response")}
\NormalTok{axs[}\DecValTok{1}\NormalTok{].set\_xlabel(}\StringTok{"frequency/Hz"}\NormalTok{)}
\NormalTok{axs[}\DecValTok{1}\NormalTok{].set\_ylabel(}\StringTok{"phase/deg"}\NormalTok{)}
\NormalTok{axs[}\DecValTok{1}\NormalTok{].grid(}\VariableTok{True}\NormalTok{, which}\OperatorTok{=}\StringTok{"both"}\NormalTok{, ls}\OperatorTok{=}\StringTok{"{-}{-}"}\NormalTok{)}
\CommentTok{\#axs[1].legend()}

\NormalTok{plt.tight\_layout()}
\CommentTok{\#plt.show()}

\NormalTok{plt.savefig(}\StringTok{"../images/sec\_characterisation/freqResponseBandpass.png"}\NormalTok{, }\BuiltInTok{format}\OperatorTok{=}\StringTok{"png"}\NormalTok{, dpi}\OperatorTok{=}\DecValTok{1000}\NormalTok{)}
\end{Highlighting}
\end{Shaded}

}

\end{codelisting}%

\begin{figure}

\centering{

\pandocbounded{\includegraphics[keepaspectratio]{content/../images/sec_characterisation/freqResponseBandpass.png}}

}

\caption{\label{fig-freqResponseBandpass}Frequency response of the band
pass filter}

\end{figure}%

\textbf{Band stop}

The output for the band stop filter is marked in
Figure~\ref{fig-BiquadCircuit} as \(BSF\) and in the transfer function
as \(V_{04}\).

\begin{equation}\phantomsection\label{eq-TFBandstopFalse}{
\frac{V_{04}}{V_i} = \frac{\left( 1 + \frac{s^2}{\omega_0^2} \right) \cdot H_0}{\left( 1 + \frac{s}{\omega_0 Q} + \frac{s^2}{\omega_0^2} \right)}
}\end{equation}

(Renner 2025) argues that Equation~\ref{eq-TFBandstopFalse} from the
ASLK PRO Manual (Rao and Ravikumar 2012) is incorrect, as using that
equation produces inconsistent results. Using the negated form of
Equation~\ref{eq-TFBandstopFalse} as seen in
Equation~\ref{eq-TFBandstopCorrect} seems to produce the correct output.
Therefore Equation~\ref{eq-TFBandstopCorrect} will be used for further
analysis.

\begin{equation}\phantomsection\label{eq-TFBandstopCorrect}{
\frac{V_{04}}{V_i} = - \frac{\left( 1 + \frac{s^2}{\omega_0^2} \right) \cdot H_0}{\left( 1 + \frac{s}{\omega_0 Q} + \frac{s^2}{\omega_0^2} \right)}
}\end{equation}

Figure~\ref{fig-freqResponseBandstop} shows the frequency response of
the band stop, with its center frequency at \(1\,kHz\).

\begin{codelisting}

\caption{\label{lst-freqresponsebandstop}}

\centering{

\begin{Shaded}
\begin{Highlighting}[]
\CommentTok{\# Behavioral Analysis Biquad Filter}

\ImportTok{import}\NormalTok{ numpy }\ImportTok{as}\NormalTok{ np}
\ImportTok{import}\NormalTok{ matplotlib.pyplot }\ImportTok{as}\NormalTok{ plt}

\CommentTok{\# Initial values}
\NormalTok{f0 }\OperatorTok{=} \FloatTok{1e3}  \CommentTok{\# Resonance frequency in Hz}
\NormalTok{w0 }\OperatorTok{=} \DecValTok{2} \OperatorTok{*}\NormalTok{ np.pi }\OperatorTok{*}\NormalTok{ f0  }\CommentTok{\# Angular frequency in rad/s}
\NormalTok{Q }\OperatorTok{=} \DecValTok{10}  \CommentTok{\# Quality factor}
\NormalTok{H0 }\OperatorTok{=} \DecValTok{1}  \CommentTok{\# Play around with this later}

\CommentTok{\# Logarithmic frequency axis}
\NormalTok{frequencies }\OperatorTok{=}\NormalTok{ np.logspace(}\DecValTok{2}\NormalTok{, }\DecValTok{4}\NormalTok{, }\DecValTok{10000}\NormalTok{)  }\CommentTok{\# Frequency from 10\^{}2 to 10\^{}4 Hz}
\NormalTok{s }\OperatorTok{=} \OtherTok{1j} \OperatorTok{*} \DecValTok{2} \OperatorTok{*}\NormalTok{ np.pi }\OperatorTok{*}\NormalTok{ frequencies  }\CommentTok{\# Laplace{-}Variable s = jω}

\CommentTok{\#\#\#\#\#\#\#\#\#\#\#\#\#\#\#\#\#\#\#\#\#\#\#\#\#\#\#\#\#\#\#\#\#\#\#\#\#\#\#\#\#\#\#\#}
\CommentTok{\# Transfer functions of Active Filters}
\CommentTok{\#\#\#\#\#\#\#\#\#\#\#\#\#\#\#\#\#\#\#\#\#\#\#\#\#\#\#\#\#\#\#\#\#\#\#\#\#\#\#\#\#\#\#\#}

\CommentTok{\#\#\# Numerator}
\CommentTok{\# Low Pass Filter}
\NormalTok{b\_lp }\OperatorTok{=}\NormalTok{ H0}

\CommentTok{\# High Pass Filter}
\NormalTok{b\_hp }\OperatorTok{=}\NormalTok{ (H0 }\OperatorTok{*}\NormalTok{ (s}\OperatorTok{**}\DecValTok{2} \OperatorTok{/}\NormalTok{ w0}\OperatorTok{**}\DecValTok{2}\NormalTok{))}

\CommentTok{\# Band Pass Filter}
\NormalTok{b\_bp }\OperatorTok{=}\NormalTok{ (}\OperatorTok{{-}}\NormalTok{H0 }\OperatorTok{*}\NormalTok{ (s }\OperatorTok{/}\NormalTok{ w0))}

\CommentTok{\# Band Stop Filter}
\NormalTok{b\_bs }\OperatorTok{=} \OperatorTok{{-}}\NormalTok{((}\DecValTok{1} \OperatorTok{+}\NormalTok{ (s}\OperatorTok{**}\DecValTok{2} \OperatorTok{/}\NormalTok{ (w0}\OperatorTok{**}\DecValTok{2}\NormalTok{))) }\OperatorTok{*}\NormalTok{ H0)}

\CommentTok{\# Denominator {-}\textgreater{} for all filters the same}
\NormalTok{a0 }\OperatorTok{=} \DecValTok{1}
\NormalTok{a1 }\OperatorTok{=}\NormalTok{ (s }\OperatorTok{/}\NormalTok{ (w0 }\OperatorTok{*}\NormalTok{ Q))}
\NormalTok{a2 }\OperatorTok{=}\NormalTok{ (s}\OperatorTok{**}\DecValTok{2} \OperatorTok{/}\NormalTok{ (w0}\OperatorTok{**}\DecValTok{2}\NormalTok{))}

\NormalTok{den }\OperatorTok{=}\NormalTok{ a0 }\OperatorTok{+}\NormalTok{ a1 }\OperatorTok{+}\NormalTok{ a2}

\CommentTok{\#\#\#\#\#\#\#\#\#\#\#\#\#\#\#\#\#\#\#\#\#\#\#\#\#\#\#\#\#\#\#\#\#\#\#\#\#\#\#\#\#\#\#\#}
\CommentTok{\# Calculation of the transfer functions H(s)}
\CommentTok{\#\#\#\#\#\#\#\#\#\#\#\#\#\#\#\#\#\#\#\#\#\#\#\#\#\#\#\#\#\#\#\#\#\#\#\#\#\#\#\#\#\#\#\#}
\NormalTok{Hs\_lp }\OperatorTok{=}\NormalTok{ b\_lp }\OperatorTok{/}\NormalTok{ den}
\NormalTok{Hs\_hp }\OperatorTok{=}\NormalTok{ b\_hp }\OperatorTok{/}\NormalTok{ den}
\NormalTok{Hs\_bp }\OperatorTok{=}\NormalTok{ b\_bp }\OperatorTok{/}\NormalTok{ den}
\NormalTok{Hs\_bs }\OperatorTok{=}\NormalTok{ b\_bs }\OperatorTok{/}\NormalTok{ den}

\CommentTok{\# Bode Diagram}
\NormalTok{fig, axs }\OperatorTok{=}\NormalTok{ plt.subplots(}\DecValTok{2}\NormalTok{)}
\CommentTok{\#fig.suptitle("frequency response of biquad filter")}
\CommentTok{\textquotesingle{}\textquotesingle{}\textquotesingle{}}
\CommentTok{\# Low Pass Filter}
\CommentTok{axs[0].semilogx(frequencies, 20 * np.log10(np.abs(Hs\_lp)), label=\textquotesingle{}low pass\textquotesingle{})}
\CommentTok{axs[1].semilogx(frequencies, np.unwrap(np.angle(Hs\_lp)) * (180 / np.pi), label=\textquotesingle{}low pass\textquotesingle{})}

\CommentTok{\# High Pass Filter}
\CommentTok{axs[0].semilogx(frequencies, 20 * np.log10(np.abs(Hs\_hp)), label=\textquotesingle{}high pass\textquotesingle{})}
\CommentTok{axs[1].semilogx(frequencies, np.unwrap(np.angle(Hs\_hp)) * (180 / np.pi), label=\textquotesingle{}high pass\textquotesingle{})}

\CommentTok{\# Band Pass Filter}
\CommentTok{axs[0].semilogx(frequencies, 20 * np.log10(np.abs(Hs\_bp)), label=\textquotesingle{}band pass\textquotesingle{})}
\CommentTok{axs[1].semilogx(frequencies, np.unwrap(np.angle(Hs\_bp)) * (180 / np.pi), label=\textquotesingle{}band pass\textquotesingle{})}
\CommentTok{\textquotesingle{}\textquotesingle{}\textquotesingle{}}
\CommentTok{\# Band Stop Filter}
\NormalTok{axs[}\DecValTok{0}\NormalTok{].semilogx(frequencies, }\DecValTok{20} \OperatorTok{*}\NormalTok{ np.log10(np.}\BuiltInTok{abs}\NormalTok{(Hs\_bs)), label}\OperatorTok{=}\StringTok{\textquotesingle{}band stop\textquotesingle{}}\NormalTok{)}
\NormalTok{axs[}\DecValTok{1}\NormalTok{].semilogx(frequencies, (np.angle(Hs\_bs)) }\OperatorTok{*}\NormalTok{ (}\DecValTok{180} \OperatorTok{/}\NormalTok{ np.pi), label}\OperatorTok{=}\StringTok{\textquotesingle{}band stop\textquotesingle{}}\NormalTok{)}

\CommentTok{\#axs[0].title("amplitude response")}
\NormalTok{axs[}\DecValTok{0}\NormalTok{].set\_xlabel(}\StringTok{"frequency/Hz"}\NormalTok{)}
\NormalTok{axs[}\DecValTok{0}\NormalTok{].set\_ylabel(}\StringTok{"amplitude/dB"}\NormalTok{)}
\NormalTok{axs[}\DecValTok{0}\NormalTok{].set\_ylim(}\OperatorTok{{-}}\DecValTok{50}\NormalTok{, }\DecValTok{25}\NormalTok{)}
\NormalTok{axs[}\DecValTok{0}\NormalTok{].grid(}\VariableTok{True}\NormalTok{, which}\OperatorTok{=}\StringTok{"both"}\NormalTok{, ls}\OperatorTok{=}\StringTok{"{-}{-}"}\NormalTok{)}
\CommentTok{\#axs[0].legend(loc=1)}

\CommentTok{\#axs[1].title("phase response")}
\NormalTok{axs[}\DecValTok{1}\NormalTok{].set\_xlabel(}\StringTok{"frequency/Hz"}\NormalTok{)}
\NormalTok{axs[}\DecValTok{1}\NormalTok{].set\_ylabel(}\StringTok{"phase/deg"}\NormalTok{)}
\NormalTok{axs[}\DecValTok{1}\NormalTok{].grid(}\VariableTok{True}\NormalTok{, which}\OperatorTok{=}\StringTok{"both"}\NormalTok{, ls}\OperatorTok{=}\StringTok{"{-}{-}"}\NormalTok{)}
\CommentTok{\#axs[1].legend()}

\NormalTok{plt.tight\_layout()}
\CommentTok{\#plt.show()}

\NormalTok{plt.savefig(}\StringTok{"../images/sec\_characterisation/freqResponseBandstop.png"}\NormalTok{, }\BuiltInTok{format}\OperatorTok{=}\StringTok{"png"}\NormalTok{, dpi}\OperatorTok{=}\DecValTok{1000}\NormalTok{)}
\end{Highlighting}
\end{Shaded}

}

\end{codelisting}%

\begin{figure}

\centering{

\pandocbounded{\includegraphics[keepaspectratio]{content/../images/sec_characterisation/freqResponseBandstop.png}}

}

\caption{\label{fig-freqResponseBandstop}Frequency response of the band
stop filter}

\end{figure}%

\textbf{Comparison}

Figure~\ref{fig-freqResponseFilter} shows all four frequency responses
together in oen graph. It shows nicely that all three pass filters peak
at the same freqeuncy and the same height. Comparing this plot with the
one from (Rao and Ravikumar 2012) gives reason to argue that the filter
design on a theoretic system level should work.

\begin{codelisting}

\caption{\label{lst-freqresponsefilter}}

\centering{

\begin{Shaded}
\begin{Highlighting}[]
\CommentTok{\# Behavioral Analysis Biquad Filter}

\ImportTok{import}\NormalTok{ numpy }\ImportTok{as}\NormalTok{ np}
\ImportTok{import}\NormalTok{ matplotlib.pyplot }\ImportTok{as}\NormalTok{ plt}

\CommentTok{\# Initial values}
\NormalTok{f0 }\OperatorTok{=} \FloatTok{1e3}  \CommentTok{\# Resonance frequency in Hz}
\NormalTok{R }\OperatorTok{=} \FloatTok{1e3}
\NormalTok{C }\OperatorTok{=} \FloatTok{159e{-}9}
\NormalTok{w0 }\OperatorTok{=} \DecValTok{2} \OperatorTok{*}\NormalTok{ np.pi }\OperatorTok{*}\NormalTok{ f0  }\CommentTok{\# Angular frequency in rad/s}
\NormalTok{w0 }\OperatorTok{=} \DecValTok{1} \OperatorTok{/}\NormalTok{ (R}\OperatorTok{*}\NormalTok{C)}
\NormalTok{Q }\OperatorTok{=} \DecValTok{10}  \CommentTok{\# Quality factor}
\NormalTok{H0 }\OperatorTok{=} \DecValTok{1}  \CommentTok{\# Play around with this later}

\CommentTok{\# Logarithmic frequency axis}
\NormalTok{frequencies }\OperatorTok{=}\NormalTok{ np.logspace(}\DecValTok{2}\NormalTok{, }\DecValTok{4}\NormalTok{, }\DecValTok{10000}\NormalTok{)  }\CommentTok{\# Frequency from 10\^{}2 to 10\^{}4 Hz}
\NormalTok{s }\OperatorTok{=} \OtherTok{1j} \OperatorTok{*} \DecValTok{2} \OperatorTok{*}\NormalTok{ np.pi }\OperatorTok{*}\NormalTok{ frequencies  }\CommentTok{\# Laplace{-}Variable s = jω}

\CommentTok{\#\#\#\#\#\#\#\#\#\#\#\#\#\#\#\#\#\#\#\#\#\#\#\#\#\#\#\#\#\#\#\#\#\#\#\#\#\#\#\#\#\#\#\#}
\CommentTok{\# Transfer functions of Active Filters}
\CommentTok{\#\#\#\#\#\#\#\#\#\#\#\#\#\#\#\#\#\#\#\#\#\#\#\#\#\#\#\#\#\#\#\#\#\#\#\#\#\#\#\#\#\#\#\#}

\CommentTok{\#\#\# Numerator}
\CommentTok{\# Low Pass Filter}
\NormalTok{b\_lp }\OperatorTok{=}\NormalTok{ H0}

\CommentTok{\# High Pass Filter}
\NormalTok{b\_hp }\OperatorTok{=}\NormalTok{ (H0 }\OperatorTok{*}\NormalTok{ (s}\OperatorTok{**}\DecValTok{2} \OperatorTok{/}\NormalTok{ w0}\OperatorTok{**}\DecValTok{2}\NormalTok{))}

\CommentTok{\# Band Pass Filter}
\NormalTok{b\_bp }\OperatorTok{=}\NormalTok{ (}\OperatorTok{{-}}\NormalTok{H0 }\OperatorTok{*}\NormalTok{ (s }\OperatorTok{/}\NormalTok{ w0))}

\CommentTok{\# Band Stop Filter}
\NormalTok{b\_bs }\OperatorTok{=} \OperatorTok{{-}}\NormalTok{((}\DecValTok{1} \OperatorTok{+}\NormalTok{ (s}\OperatorTok{**}\DecValTok{2} \OperatorTok{/}\NormalTok{ (w0}\OperatorTok{**}\DecValTok{2}\NormalTok{))) }\OperatorTok{*}\NormalTok{ H0)}

\CommentTok{\# Denominator {-}\textgreater{} for all filters the same}
\NormalTok{a0 }\OperatorTok{=} \DecValTok{1}
\NormalTok{a1 }\OperatorTok{=}\NormalTok{ (s }\OperatorTok{/}\NormalTok{ (w0 }\OperatorTok{*}\NormalTok{ Q))}
\NormalTok{a2 }\OperatorTok{=}\NormalTok{ (s}\OperatorTok{**}\DecValTok{2} \OperatorTok{/}\NormalTok{ (w0}\OperatorTok{**}\DecValTok{2}\NormalTok{))}

\NormalTok{den }\OperatorTok{=}\NormalTok{ a0 }\OperatorTok{+}\NormalTok{ a1 }\OperatorTok{+}\NormalTok{ a2}

\CommentTok{\#\#\#\#\#\#\#\#\#\#\#\#\#\#\#\#\#\#\#\#\#\#\#\#\#\#\#\#\#\#\#\#\#\#\#\#\#\#\#\#\#\#\#\#}
\CommentTok{\# Calculation of the transfer functions H(s)}
\CommentTok{\#\#\#\#\#\#\#\#\#\#\#\#\#\#\#\#\#\#\#\#\#\#\#\#\#\#\#\#\#\#\#\#\#\#\#\#\#\#\#\#\#\#\#\#}
\NormalTok{Hs\_lp }\OperatorTok{=}\NormalTok{ b\_lp }\OperatorTok{/}\NormalTok{ den}
\NormalTok{Hs\_hp }\OperatorTok{=}\NormalTok{ b\_hp }\OperatorTok{/}\NormalTok{ den}
\NormalTok{Hs\_bp }\OperatorTok{=}\NormalTok{ b\_bp }\OperatorTok{/}\NormalTok{ den}
\NormalTok{Hs\_bs }\OperatorTok{=}\NormalTok{ b\_bs }\OperatorTok{/}\NormalTok{ den}

\CommentTok{\# Bode Diagram}
\NormalTok{fig, axs }\OperatorTok{=}\NormalTok{ plt.subplots(}\DecValTok{2}\NormalTok{)}
\CommentTok{\#fig.suptitle("frequency response of biquad filter")}

\CommentTok{\# Low Pass Filter}
\NormalTok{axs[}\DecValTok{0}\NormalTok{].semilogx(frequencies, }\DecValTok{20} \OperatorTok{*}\NormalTok{ np.log10(np.}\BuiltInTok{abs}\NormalTok{(Hs\_lp)), label}\OperatorTok{=}\StringTok{\textquotesingle{}low pass\textquotesingle{}}\NormalTok{)}
\NormalTok{axs[}\DecValTok{1}\NormalTok{].semilogx(frequencies, np.unwrap(np.angle(Hs\_lp)) }\OperatorTok{*}\NormalTok{ (}\DecValTok{180} \OperatorTok{/}\NormalTok{ np.pi), label}\OperatorTok{=}\StringTok{\textquotesingle{}low pass\textquotesingle{}}\NormalTok{)}

\CommentTok{\# High Pass Filter}
\NormalTok{axs[}\DecValTok{0}\NormalTok{].semilogx(frequencies, }\DecValTok{20} \OperatorTok{*}\NormalTok{ np.log10(np.}\BuiltInTok{abs}\NormalTok{(Hs\_hp)), label}\OperatorTok{=}\StringTok{\textquotesingle{}high pass\textquotesingle{}}\NormalTok{)}
\NormalTok{axs[}\DecValTok{1}\NormalTok{].semilogx(frequencies, np.unwrap(np.angle(Hs\_hp)) }\OperatorTok{*}\NormalTok{ (}\DecValTok{180} \OperatorTok{/}\NormalTok{ np.pi), label}\OperatorTok{=}\StringTok{\textquotesingle{}high pass\textquotesingle{}}\NormalTok{)}

\CommentTok{\# Band Pass Filter}
\NormalTok{axs[}\DecValTok{0}\NormalTok{].semilogx(frequencies, }\DecValTok{20} \OperatorTok{*}\NormalTok{ np.log10(np.}\BuiltInTok{abs}\NormalTok{(Hs\_bp)), label}\OperatorTok{=}\StringTok{\textquotesingle{}band pass\textquotesingle{}}\NormalTok{)}
\NormalTok{axs[}\DecValTok{1}\NormalTok{].semilogx(frequencies, np.unwrap(np.angle(Hs\_bp)) }\OperatorTok{*}\NormalTok{ (}\DecValTok{180} \OperatorTok{/}\NormalTok{ np.pi), label}\OperatorTok{=}\StringTok{\textquotesingle{}band pass\textquotesingle{}}\NormalTok{)}

\CommentTok{\# Band Stop Filter}
\NormalTok{axs[}\DecValTok{0}\NormalTok{].semilogx(frequencies, }\DecValTok{20} \OperatorTok{*}\NormalTok{ np.log10(np.}\BuiltInTok{abs}\NormalTok{(Hs\_bs)), label}\OperatorTok{=}\StringTok{\textquotesingle{}band stop\textquotesingle{}}\NormalTok{)}
\NormalTok{axs[}\DecValTok{1}\NormalTok{].semilogx(frequencies, (np.angle(Hs\_bs)) }\OperatorTok{*}\NormalTok{ (}\DecValTok{180} \OperatorTok{/}\NormalTok{ np.pi), label}\OperatorTok{=}\StringTok{\textquotesingle{}band stop\textquotesingle{}}\NormalTok{)}

\CommentTok{\#axs[0].title("amplitude response")}
\NormalTok{axs[}\DecValTok{0}\NormalTok{].set\_xlabel(}\StringTok{"frequency/Hz"}\NormalTok{)}
\NormalTok{axs[}\DecValTok{0}\NormalTok{].set\_ylabel(}\StringTok{"amplitude/dB"}\NormalTok{)}
\NormalTok{axs[}\DecValTok{0}\NormalTok{].set\_ylim(}\OperatorTok{{-}}\DecValTok{50}\NormalTok{, }\DecValTok{40}\NormalTok{)}
\NormalTok{axs[}\DecValTok{0}\NormalTok{].grid(}\VariableTok{True}\NormalTok{, which}\OperatorTok{=}\StringTok{"both"}\NormalTok{, ls}\OperatorTok{=}\StringTok{"{-}{-}"}\NormalTok{)}
\NormalTok{axs[}\DecValTok{0}\NormalTok{].legend(loc}\OperatorTok{=}\DecValTok{1}\NormalTok{)}

\CommentTok{\#axs[1].title("phase response")}
\NormalTok{axs[}\DecValTok{1}\NormalTok{].set\_xlabel(}\StringTok{"frequency/Hz"}\NormalTok{)}
\NormalTok{axs[}\DecValTok{1}\NormalTok{].set\_ylabel(}\StringTok{"phase/deg"}\NormalTok{)}
\NormalTok{axs[}\DecValTok{1}\NormalTok{].grid(}\VariableTok{True}\NormalTok{, which}\OperatorTok{=}\StringTok{"both"}\NormalTok{, ls}\OperatorTok{=}\StringTok{"{-}{-}"}\NormalTok{)}
\NormalTok{axs[}\DecValTok{1}\NormalTok{].legend()}

\NormalTok{plt.tight\_layout()}
\CommentTok{\#plt.show()}

\NormalTok{plt.savefig(}\StringTok{"../images/sec\_characterisation/freqResponseFilter.png"}\NormalTok{, }\BuiltInTok{format}\OperatorTok{=}\StringTok{"png"}\NormalTok{, dpi}\OperatorTok{=}\DecValTok{1000}\NormalTok{)}
\end{Highlighting}
\end{Shaded}

}

\end{codelisting}%

\begin{figure}

\centering{

\pandocbounded{\includegraphics[keepaspectratio]{content/../images/sec_characterisation/freqResponseFilter.png}}

}

\caption{\label{fig-freqResponseFilter}Behauvioural analysis of biquad
filter}

\end{figure}%

\subsection{Stability}\label{stability}

The stability of the biquad is checked at different hierarchical levels.
The first analysis considers the system from a theorectical standpoint
with transfer functions, and checks if conceptual design of the biquad
filter is stable. On a component level the stability of the integrators
and adders is analyzed, to verify that the chosen values for resistors
and capcitors do not induce oscillations through the feedback loop.

\subsubsection{System stability}\label{system-stability}

A system is stable if its impulse response is absolutley integrateable.
In case of a given transfer function, this can also be checked by
calculating the poles of the transfer function. If all the poles lay in
the left half of the s-plane, the system is considered stable. There is
a special case where single poles can lay on the \(j\omega\)-axis, on
their own or in combination with poles in the left half of the s-plane.
Systems which fall under that, are called marginally stable. (Fliege
1991)

\subsubsection{Pole-zero plot}\label{pole-zero-plot}

\begin{codelisting}

\caption{\label{lst-polezerostability}}

\centering{

\begin{Shaded}
\begin{Highlighting}[]
\ImportTok{import}\NormalTok{ numpy }\ImportTok{as}\NormalTok{ np}
\ImportTok{import}\NormalTok{ matplotlib.pyplot }\ImportTok{as}\NormalTok{ plt}
\ImportTok{from}\NormalTok{ scipy.signal }\ImportTok{import}\NormalTok{ tf2zpk}

\CommentTok{\# Given values}
\NormalTok{f }\OperatorTok{=} \FloatTok{1e3}
\NormalTok{w }\OperatorTok{=} \DecValTok{2} \OperatorTok{*}\NormalTok{ np.pi }\OperatorTok{*}\NormalTok{ f}
\NormalTok{R }\OperatorTok{=} \FloatTok{1e3}
\NormalTok{C }\OperatorTok{=} \DecValTok{1} \OperatorTok{/}\NormalTok{ (w }\OperatorTok{*}\NormalTok{ R)}
\NormalTok{Q }\OperatorTok{=} \DecValTok{10}
\NormalTok{H0 }\OperatorTok{=} \DecValTok{1}

\CommentTok{\# Calculate w0}
\NormalTok{w0 }\OperatorTok{=} \DecValTok{1} \OperatorTok{/}\NormalTok{ (R }\OperatorTok{*}\NormalTok{ C)}

\CommentTok{\# Transfer function coefficients}
\NormalTok{a2 }\OperatorTok{=} \DecValTok{1} \OperatorTok{/}\NormalTok{ w0}\OperatorTok{**}\DecValTok{2}
\NormalTok{a1 }\OperatorTok{=} \DecValTok{1} \OperatorTok{/}\NormalTok{ (w0 }\OperatorTok{*}\NormalTok{ Q)}
\NormalTok{a0 }\OperatorTok{=} \DecValTok{1}

\CommentTok{\# Define transfer functions manually as (numerator, denominator) pairs}
\NormalTok{systems }\OperatorTok{=}\NormalTok{ \{}
    \StringTok{\textquotesingle{}Low pass filter\textquotesingle{}}\NormalTok{: ([H0], [a2, a1, a0]),}
    \StringTok{\textquotesingle{}High pass filter\textquotesingle{}}\NormalTok{: ([H0 }\OperatorTok{/}\NormalTok{ w0}\OperatorTok{**}\DecValTok{2}\NormalTok{, }\DecValTok{0}\NormalTok{, }\DecValTok{0}\NormalTok{], [a2, a1, a0]),}
    \StringTok{\textquotesingle{}Band pass filter\textquotesingle{}}\NormalTok{: ([}\OperatorTok{{-}}\NormalTok{H0 }\OperatorTok{/}\NormalTok{ w0, }\DecValTok{0}\NormalTok{], [a2, a1, a0]),}
    \StringTok{\textquotesingle{}Band stop filter\textquotesingle{}}\NormalTok{: ([H0 }\OperatorTok{/}\NormalTok{ w0}\OperatorTok{**}\DecValTok{2}\NormalTok{, }\DecValTok{0}\NormalTok{, H0], [a2, a1, a0])}
\NormalTok{\}}

\CommentTok{\# Function to plot pole{-}zero map}
\KeywordTok{def}\NormalTok{ plot\_pzmap(num, den, title, subplot\_pos):}
\NormalTok{    zeros, poles, \_ }\OperatorTok{=}\NormalTok{ tf2zpk(num, den)}
\NormalTok{    plt.subplot(}\DecValTok{2}\NormalTok{, }\DecValTok{2}\NormalTok{, subplot\_pos)}
\NormalTok{    plt.plot(np.real(zeros), np.imag(zeros), }\StringTok{\textquotesingle{}go\textquotesingle{}}\NormalTok{, label}\OperatorTok{=}\StringTok{\textquotesingle{}Zeros\textquotesingle{}}\NormalTok{)}
\NormalTok{    plt.plot(np.real(poles), np.imag(poles), }\StringTok{\textquotesingle{}rx\textquotesingle{}}\NormalTok{, label}\OperatorTok{=}\StringTok{\textquotesingle{}Poles\textquotesingle{}}\NormalTok{)}
\NormalTok{    plt.axhline(}\DecValTok{0}\NormalTok{, color}\OperatorTok{=}\StringTok{\textquotesingle{}gray\textquotesingle{}}\NormalTok{, lw}\OperatorTok{=}\FloatTok{0.5}\NormalTok{)}
\NormalTok{    plt.axvline(}\DecValTok{0}\NormalTok{, color}\OperatorTok{=}\StringTok{\textquotesingle{}gray\textquotesingle{}}\NormalTok{, lw}\OperatorTok{=}\FloatTok{0.5}\NormalTok{)}
\NormalTok{    plt.title(title)}
\NormalTok{    plt.xlabel(}\StringTok{\textquotesingle{}$}\ErrorTok{\textbackslash{}}\StringTok{sigma$\textquotesingle{}}\NormalTok{)}
\NormalTok{    plt.ylabel(}\StringTok{\textquotesingle{}$j}\ErrorTok{\textbackslash{}}\StringTok{omega$\textquotesingle{}}\NormalTok{)}
\NormalTok{    plt.xlim([}\OperatorTok{{-}}\DecValTok{3500}\NormalTok{, }\DecValTok{1500}\NormalTok{])}
\NormalTok{    plt.ylim([}\OperatorTok{{-}}\DecValTok{10000}\NormalTok{, }\DecValTok{10000}\NormalTok{])}
\NormalTok{    plt.grid(}\VariableTok{True}\NormalTok{)}
\NormalTok{    plt.legend(loc}\OperatorTok{=}\StringTok{\textquotesingle{}upper right\textquotesingle{}}\NormalTok{)}

\CommentTok{\# Plot all systems}
\NormalTok{plt.figure(figsize}\OperatorTok{=}\NormalTok{(}\DecValTok{12}\NormalTok{, }\DecValTok{10}\NormalTok{))}
\ControlFlowTok{for}\NormalTok{ i, (title, (num, den)) }\KeywordTok{in} \BuiltInTok{enumerate}\NormalTok{(systems.items(), }\DecValTok{1}\NormalTok{):}
\NormalTok{    plot\_pzmap(num, den, title, i)}

\NormalTok{plt.tight\_layout()}
\CommentTok{\#plt.show()}

\NormalTok{plt.savefig(}\StringTok{"../images/sec\_characterisation/poleZeroStability.png"}\NormalTok{, }\BuiltInTok{format}\OperatorTok{=}\StringTok{"png"}\NormalTok{, dpi}\OperatorTok{=}\DecValTok{1000}\NormalTok{)}
\end{Highlighting}
\end{Shaded}

}

\end{codelisting}%

\begin{figure}

\centering{

\pandocbounded{\includegraphics[keepaspectratio]{content/../images/sec_characterisation/poleZeroStability.png}}

}

\caption{\label{fig-poleZeroStability}Pole-zero plot for all transfer
functions}

\end{figure}%

Figure~\ref{fig-poleZeroStability} shows the pole-zero plots of all four
filters, low pass, high pass, band pass and band stop. In all four plots
the poles are located in the left half of the s-plane and the system can
therefore theoretically be classified as stable. (Razavi 2018) confirms
this, as the article explains that with \(Q \rightarrow \infty\) the
poles of the system approach the \(j\omega\) axis and the system becomes
unstable.

This analysis only considers the system as a mathematical model and as a
whole. Further considerations regarding the stability of the components,
integrators and adders, have to be done.

\subsubsection{Component stability}\label{component-stability}

Circuits with opamps often have feedback loops, meaning that the output
of the operational amplifier is somehow connected to the inverted input
of the opamp. These feedback loops become problematic when the feedback
signal is in phase with the input signal, as positive feedback is
created and the circuit is working as an oscillator. (Reisch 2007)

The stability of the non-inverting amplifier can be verified by
calculating the phase reserve \(\alpha\) of the circuit. If \(f_k\) is
the frequency where the feedback gain is equal to 1 and \(\varphi_k\) is
the corresponding phase to that frequency, then the phase reserve is
calculated by:

\[
\alpha = 180° - \varphi_k
\]

For circuits to be considered stable, the phase reserve has to be
positive. To reduce overshoots during the transient response, it is
customary to have a phase reserve of \(\alpha > 45°\). (Reisch 2007)

\begin{figure}

\centering{

\pandocbounded{\includegraphics[keepaspectratio]{images/sec_characterisation/reisch_stability_example_trans.png}}

}

\caption{\label{fig-reischExampleStabilityTransient}Example of a
transcient response of a circuit with a phase reserve of
\(\alpha = 5.7°\) (taken from Reisch 2007)}

\end{figure}%

Figure Figure~\ref{fig-reischExampleStabilityTransient} shows the
transient response of a circuit with a phase reserve of
\(\alpha = 5.7°\). The overshoots are clearly visible and number of the
overshoots per puls are larger then the customary ``one over, one
under''-rule. As the phase reserve is positive, the figure shows that
even though the transient response is not ideal, the oscillations are
attenuated and the circuit is can be considered as stable.

In practical application, the phase reserve can be graphically
determined with the help of bode diagrams. The bode diagram of the
circuit with an open feedback loop is simulated, so that the frequency
\(f_k\) can be read out. This is the frequency where the feedback gain
is 1 or 0 dB. The corresponding frequency to that, is the phase of the
feedback gain \(\varphi_k\), the difference between \(-180°\) and
\(\varphi_k\) is the phase reserve \(\alpha\). (Reisch 2007)

For the analysis of component stability the 5t-ota design from (Pretl,
Koefinger, and Dorrer 2025) was used.
Figure~\ref{fig-circuitStabilityAdder} and
Figure~\ref{fig-circuitStabilityIntegrator} show the schematics that
simulated the stability analysis for the components. In both cases the
AC source was inserted in the open feedback loop and the output of the
OTA was used to analyse the stability.

\begin{figure}

\centering{

\pandocbounded{\includegraphics[keepaspectratio]{content/../images/sec_characterisation/stability_adder_circuit.png}}

}

\caption{\label{fig-circuitStabilityAdder}Circuit schematic of the
stability analysis of the adder}

\end{figure}%

\begin{figure}

\centering{

\pandocbounded{\includegraphics[keepaspectratio]{content/../images/sec_characterisation/stability_integrator_circuit.png}}

}

\caption{\label{fig-circuitStabilityIntegrator}Circuit schematic of the
stability analysis of the integrator}

\end{figure}%

In the following figures Figure~\ref{fig-stabilityAdder} and
Figure~\ref{fig-stabilityIntegrator} this stability analysis method was
used to determine the stability over the phase reserve.

\begin{codelisting}

\caption{\label{lst-stabilityadder}}

\centering{

\begin{Shaded}
\begin{Highlighting}[]
\CommentTok{\# Stability analysis adder}

\ImportTok{import}\NormalTok{ numpy }\ImportTok{as}\NormalTok{ np}
\ImportTok{import}\NormalTok{ matplotlib.pyplot }\ImportTok{as}\NormalTok{ plt}
\ImportTok{import}\NormalTok{ sys}
\NormalTok{sys.path.insert(}\DecValTok{0}\NormalTok{, }\StringTok{\textquotesingle{}../simulation\textquotesingle{}}\NormalTok{)}
\ImportTok{import}\NormalTok{ ltspy3}

\NormalTok{sd}\OperatorTok{=}\NormalTok{ltspy3.SimData(}\StringTok{\textquotesingle{}../simulation/stability\_adder.raw\textquotesingle{}}\NormalTok{,[}\StringTok{b\textquotesingle{}v(v\_out)\textquotesingle{}}\NormalTok{,}\StringTok{b\textquotesingle{}frequency\textquotesingle{}}\NormalTok{])}

\NormalTok{nvout }\OperatorTok{=}\NormalTok{ sd.variables.index(}\StringTok{b\textquotesingle{}v(v\_out)\textquotesingle{}}\NormalTok{)}
\NormalTok{nfrequency }\OperatorTok{=}\NormalTok{ sd.variables.index(}\StringTok{b\textquotesingle{}frequency\textquotesingle{}}\NormalTok{)}

\NormalTok{fig, (ax1, ax2) }\OperatorTok{=}\NormalTok{ plt.subplots(}\DecValTok{2}\NormalTok{, }\DecValTok{1}\NormalTok{, figsize}\OperatorTok{=}\NormalTok{(}\DecValTok{10}\NormalTok{, }\DecValTok{6}\NormalTok{), sharex}\OperatorTok{=}\VariableTok{True}\NormalTok{)}

\NormalTok{ax1.semilogx(sd.values[nfrequency],}\DecValTok{20}\OperatorTok{*}\NormalTok{np.log10(}\BuiltInTok{abs}\NormalTok{(sd.values[nvout])))}
\NormalTok{ax1.set\_ylabel(}\StringTok{"Magnitude/dB"}\NormalTok{)}
\NormalTok{ax1.axvline(}\FloatTok{3e5}\NormalTok{,color}\OperatorTok{=}\StringTok{\textquotesingle{}red\textquotesingle{}}\NormalTok{,linestyle}\OperatorTok{=}\StringTok{\textquotesingle{}{-}{-}\textquotesingle{}}\NormalTok{)}
\NormalTok{ax1.grid(}\VariableTok{True}\NormalTok{, which}\OperatorTok{=}\StringTok{"both"}\NormalTok{, ls}\OperatorTok{=}\StringTok{"{-}{-}"}\NormalTok{)}

\NormalTok{ax2.semilogx(sd.values[nfrequency],np.angle(sd.values[nvout], deg}\OperatorTok{=}\VariableTok{True}\NormalTok{))}
\NormalTok{ax2.axvline(}\FloatTok{3e5}\NormalTok{,color}\OperatorTok{=}\StringTok{\textquotesingle{}red\textquotesingle{}}\NormalTok{,linestyle}\OperatorTok{=}\StringTok{\textquotesingle{}{-}{-}\textquotesingle{}}\NormalTok{)}
\NormalTok{ax2.axhline(}\FloatTok{82.75}\NormalTok{,color}\OperatorTok{=}\StringTok{\textquotesingle{}red\textquotesingle{}}\NormalTok{,linestyle}\OperatorTok{=}\StringTok{\textquotesingle{}{-}{-}\textquotesingle{}}\NormalTok{)}
\NormalTok{ax2.set\_ylabel(}\StringTok{"Phase/deg"}\NormalTok{)}
\NormalTok{ax2.set\_xlabel(}\StringTok{"Frequency/Hz"}\NormalTok{)}


\NormalTok{plt.grid(}\VariableTok{True}\NormalTok{, which}\OperatorTok{=}\StringTok{"both"}\NormalTok{, ls}\OperatorTok{=}\StringTok{"{-}{-}"}\NormalTok{)}
\NormalTok{plt.tight\_layout()}
\CommentTok{\#plt.show()}

\NormalTok{plt.savefig(}\StringTok{"../images/sec\_characterisation/stabilityAdder.png"}\NormalTok{, }\BuiltInTok{format}\OperatorTok{=}\StringTok{"png"}\NormalTok{, dpi}\OperatorTok{=}\DecValTok{1000}\NormalTok{)}
\end{Highlighting}
\end{Shaded}

}

\end{codelisting}%

\begin{figure}

\centering{

\pandocbounded{\includegraphics[keepaspectratio]{content/../images/sec_characterisation/stabilityAdder.png}}

}

\caption{\label{fig-stabilityAdder}Stability analysis of the adder}

\end{figure}%

The frequency response of the adder shows a phase \(\varphi_k = 83°\),
when the gain is 1. Calculation the phase reserve from that

\[\alpha_{add} = 180° - \varphi_k = 180° - 83° = 97° > 45°
\]

gives a phase reserves, that indicates stability for this component.

\begin{codelisting}

\caption{\label{lst-stabilityintegrator}}

\centering{

\begin{Shaded}
\begin{Highlighting}[]
\CommentTok{\# Stability analysis integrator}

\ImportTok{import}\NormalTok{ numpy }\ImportTok{as}\NormalTok{ np}
\ImportTok{import}\NormalTok{ matplotlib.pyplot }\ImportTok{as}\NormalTok{ plt}
\ImportTok{import}\NormalTok{ sys}
\NormalTok{sys.path.insert(}\DecValTok{0}\NormalTok{, }\StringTok{\textquotesingle{}../simulation\textquotesingle{}}\NormalTok{)}
\ImportTok{import}\NormalTok{ ltspy3}

\NormalTok{sd}\OperatorTok{=}\NormalTok{ltspy3.SimData(}\StringTok{\textquotesingle{}../simulation/stability\_integrator.raw\textquotesingle{}}\NormalTok{,[}\StringTok{b\textquotesingle{}v(v\_out)\textquotesingle{}}\NormalTok{,}\StringTok{b\textquotesingle{}frequency\textquotesingle{}}\NormalTok{])}

\NormalTok{nvout }\OperatorTok{=}\NormalTok{ sd.variables.index(}\StringTok{b\textquotesingle{}v(v\_out)\textquotesingle{}}\NormalTok{)}
\NormalTok{nfrequency }\OperatorTok{=}\NormalTok{ sd.variables.index(}\StringTok{b\textquotesingle{}frequency\textquotesingle{}}\NormalTok{)}

\NormalTok{fig, (ax1, ax2) }\OperatorTok{=}\NormalTok{ plt.subplots(}\DecValTok{2}\NormalTok{, }\DecValTok{1}\NormalTok{, figsize}\OperatorTok{=}\NormalTok{(}\DecValTok{10}\NormalTok{, }\DecValTok{6}\NormalTok{), sharex}\OperatorTok{=}\VariableTok{True}\NormalTok{)}

\NormalTok{ax1.semilogx(sd.values[nfrequency],}\DecValTok{20}\OperatorTok{*}\NormalTok{np.log10(}\BuiltInTok{abs}\NormalTok{(sd.values[nvout])))}
\NormalTok{ax1.set\_ylabel(}\StringTok{"Magnitude/dB"}\NormalTok{)}
\NormalTok{ax1.axvline(}\FloatTok{7e5}\NormalTok{,color}\OperatorTok{=}\StringTok{\textquotesingle{}red\textquotesingle{}}\NormalTok{,linestyle}\OperatorTok{=}\StringTok{\textquotesingle{}{-}{-}\textquotesingle{}}\NormalTok{)}
\NormalTok{ax1.grid(}\VariableTok{True}\NormalTok{, which}\OperatorTok{=}\StringTok{"both"}\NormalTok{, ls}\OperatorTok{=}\StringTok{"{-}{-}"}\NormalTok{)}

\NormalTok{ax2.semilogx(sd.values[nfrequency],np.angle(sd.values[nvout], deg}\OperatorTok{=}\VariableTok{True}\NormalTok{))}
\NormalTok{ax2.axvline(}\FloatTok{7e5}\NormalTok{,color}\OperatorTok{=}\StringTok{\textquotesingle{}red\textquotesingle{}}\NormalTok{,linestyle}\OperatorTok{=}\StringTok{\textquotesingle{}{-}{-}\textquotesingle{}}\NormalTok{)}
\NormalTok{ax2.axhline(}\DecValTok{68}\NormalTok{,color}\OperatorTok{=}\StringTok{\textquotesingle{}red\textquotesingle{}}\NormalTok{,linestyle}\OperatorTok{=}\StringTok{\textquotesingle{}{-}{-}\textquotesingle{}}\NormalTok{)}
\NormalTok{ax2.set\_ylabel(}\StringTok{"Phase/deg"}\NormalTok{)}
\NormalTok{ax2.set\_xlabel(}\StringTok{"Frequency/Hz"}\NormalTok{)}

\NormalTok{plt.grid(}\VariableTok{True}\NormalTok{, which}\OperatorTok{=}\StringTok{"both"}\NormalTok{, ls}\OperatorTok{=}\StringTok{"{-}{-}"}\NormalTok{)}
\NormalTok{plt.tight\_layout()}
\CommentTok{\#plt.show()}

\NormalTok{plt.savefig(}\StringTok{"../images/sec\_characterisation/stabilityIntegrator.png"}\NormalTok{, }\BuiltInTok{format}\OperatorTok{=}\StringTok{"png"}\NormalTok{, dpi}\OperatorTok{=}\DecValTok{1000}\NormalTok{)}
\end{Highlighting}
\end{Shaded}

}

\end{codelisting}%

\begin{figure}

\centering{

\pandocbounded{\includegraphics[keepaspectratio]{content/../images/sec_characterisation/stabilityIntegrator.png}}

}

\caption{\label{fig-stabilityIntegrator}Stability analysis of the
integrator}

\end{figure}%

The stability of the integrator circuit, as seen in
Figure~\ref{fig-stabilityIntegrator}, shows a intersection of the
magnitude plot with the \(0\,dB\) line at about \(f_k = 70\,kHz\), which
corresponds to a phase of \(\varphi_k = 68°\). This would leave a phase
reserve of:

\[
\alpha_{int} = 180° - \varphi_k = 112° > 45°
\]

Therefore the integrator would be stable.

\subsection{Ideal Opamp}\label{ideal-opamp}

To check the behauviour of the implemented circuit against the modelled
behauviour of the transfer function, the universal biquad was built as
an ideal circuit with voltage-regulated current sources instead of OTAs.
This simulation of the circuit verfies that the circuit implementation
of the biquadratic filter with OTAs can fulfill the requirements at
least in the ideal case.

\begin{figure}

\centering{

\pandocbounded{\includegraphics[keepaspectratio]{content/../images/sec_characterisation/biquad_univ_circuit.png}}

}

\caption{\label{fig-CircuitBiquadIdeal}Schematic of an universal biquad
filter with ideal OTAs}

\end{figure}%

Figure~\ref{fig-CircuitBiquadIdeal} depicts the schematic of an
universal biquad filter, where the OTAs are idealised with voltage
controlled current sources. An ac simulation of this circuit can be
viewed in Figure~\ref{fig-idealCir}.

\begin{codelisting}

\caption{\label{lst-idealcir}}

\centering{

\begin{Shaded}
\begin{Highlighting}[]
\CommentTok{\# plot Ideal (voltage controlled current? source) biquad}
\ImportTok{import}\NormalTok{ numpy }\ImportTok{as}\NormalTok{ np}
\ImportTok{import}\NormalTok{ matplotlib.pyplot }\ImportTok{as}\NormalTok{ plt}
\ImportTok{import}\NormalTok{ sys}
\NormalTok{sys.path.insert(}\DecValTok{0}\NormalTok{, }\StringTok{\textquotesingle{}../simulation\textquotesingle{}}\NormalTok{)}
\ImportTok{import}\NormalTok{ ltspy3}

\NormalTok{sd}\OperatorTok{=}\NormalTok{ltspy3.SimData(}\StringTok{\textquotesingle{}../simulation/biquad\_univ.raw\textquotesingle{}}\NormalTok{)}

\NormalTok{nvoutLPF }\OperatorTok{=}\NormalTok{ sd.variables.index(}\StringTok{b\textquotesingle{}v(lpf)\textquotesingle{}}\NormalTok{)}
\NormalTok{nvoutHPF }\OperatorTok{=}\NormalTok{ sd.variables.index(}\StringTok{b\textquotesingle{}v(hpf)\textquotesingle{}}\NormalTok{)}
\NormalTok{nvoutBPF }\OperatorTok{=}\NormalTok{ sd.variables.index(}\StringTok{b\textquotesingle{}v(bpf)\textquotesingle{}}\NormalTok{)}
\NormalTok{nvoutBSF }\OperatorTok{=}\NormalTok{ sd.variables.index(}\StringTok{b\textquotesingle{}v(bsf)\textquotesingle{}}\NormalTok{)}
\NormalTok{nfrequency }\OperatorTok{=}\NormalTok{ sd.variables.index(}\StringTok{b\textquotesingle{}frequency\textquotesingle{}}\NormalTok{)}

\NormalTok{fig, (ax1, ax2) }\OperatorTok{=}\NormalTok{ plt.subplots(}\DecValTok{2}\NormalTok{, }\DecValTok{1}\NormalTok{, figsize}\OperatorTok{=}\NormalTok{(}\DecValTok{10}\NormalTok{, }\DecValTok{6}\NormalTok{), sharex}\OperatorTok{=}\VariableTok{True}\NormalTok{)}

\NormalTok{ax1.semilogx(sd.values[nfrequency],}\DecValTok{20}\OperatorTok{*}\NormalTok{np.log10(}\BuiltInTok{abs}\NormalTok{(sd.values[nvoutLPF])),label}\OperatorTok{=}\StringTok{\textquotesingle{}lpf\textquotesingle{}}\NormalTok{)}
\NormalTok{ax1.semilogx(sd.values[nfrequency],}\DecValTok{20}\OperatorTok{*}\NormalTok{np.log10(}\BuiltInTok{abs}\NormalTok{(sd.values[nvoutHPF])),label}\OperatorTok{=}\StringTok{\textquotesingle{}hpf\textquotesingle{}}\NormalTok{)}
\NormalTok{ax1.semilogx(sd.values[nfrequency],}\DecValTok{20}\OperatorTok{*}\NormalTok{np.log10(}\BuiltInTok{abs}\NormalTok{(sd.values[nvoutBPF])),label}\OperatorTok{=}\StringTok{\textquotesingle{}bpf\textquotesingle{}}\NormalTok{)}
\NormalTok{ax1.semilogx(sd.values[nfrequency],}\DecValTok{20}\OperatorTok{*}\NormalTok{np.log10(}\BuiltInTok{abs}\NormalTok{(sd.values[nvoutBSF])),label}\OperatorTok{=}\StringTok{\textquotesingle{}bsf\textquotesingle{}}\NormalTok{)}
\NormalTok{ax1.set\_xlim([}\FloatTok{10e1}\NormalTok{,}\FloatTok{10e3}\NormalTok{])}
\NormalTok{ax1.set\_ylim([}\OperatorTok{{-}}\DecValTok{40}\NormalTok{,}\DecValTok{20}\NormalTok{])}
\NormalTok{ax1.set\_ylabel(}\StringTok{"Magnitude/dB"}\NormalTok{)}
\NormalTok{ax1.grid(}\VariableTok{True}\NormalTok{, which}\OperatorTok{=}\StringTok{"both"}\NormalTok{, ls}\OperatorTok{=}\StringTok{"{-}{-}"}\NormalTok{)}
\NormalTok{ax1.legend()}

\NormalTok{ax2.semilogx(sd.values[nfrequency],np.angle(sd.values[nvoutLPF], deg}\OperatorTok{=}\VariableTok{True}\NormalTok{),label}\OperatorTok{=}\StringTok{\textquotesingle{}lpf\textquotesingle{}}\NormalTok{)}
\NormalTok{ax2.semilogx(sd.values[nfrequency],np.angle(sd.values[nvoutHPF], deg}\OperatorTok{=}\VariableTok{True}\NormalTok{),label}\OperatorTok{=}\StringTok{\textquotesingle{}hpf\textquotesingle{}}\NormalTok{)}
\NormalTok{ax2.semilogx(sd.values[nfrequency],np.angle(sd.values[nvoutBPF], deg}\OperatorTok{=}\VariableTok{True}\NormalTok{),label}\OperatorTok{=}\StringTok{\textquotesingle{}bpf\textquotesingle{}}\NormalTok{)}
\NormalTok{ax2.semilogx(sd.values[nfrequency],np.angle(sd.values[nvoutBSF], deg}\OperatorTok{=}\VariableTok{True}\NormalTok{),label}\OperatorTok{=}\StringTok{\textquotesingle{}bsf\textquotesingle{}}\NormalTok{)}
\NormalTok{ax2.set\_ylabel(}\StringTok{"Phase/deg"}\NormalTok{)}
\NormalTok{ax2.set\_xlabel(}\StringTok{"Frequency/Hz"}\NormalTok{)}
\NormalTok{ax2.legend()}

\NormalTok{plt.grid(}\VariableTok{True}\NormalTok{, which}\OperatorTok{=}\StringTok{"both"}\NormalTok{, ls}\OperatorTok{=}\StringTok{"{-}{-}"}\NormalTok{)}
\NormalTok{plt.tight\_layout()}
\CommentTok{\#plt.show()}

\NormalTok{plt.savefig(}\StringTok{"../images/sec\_characterisation/idealCir.png"}\NormalTok{, }\BuiltInTok{format}\OperatorTok{=}\StringTok{"png"}\NormalTok{, dpi}\OperatorTok{=}\DecValTok{1000}\NormalTok{)}
\end{Highlighting}
\end{Shaded}

}

\end{codelisting}%

\begin{figure}

\centering{

\pandocbounded{\includegraphics[keepaspectratio]{content/../images/sec_characterisation/idealCir.png}}

}

\caption{\label{fig-idealCir}Simulation of an idealised biquad}

\end{figure}%

Figure~\ref{fig-compTfIdealAmplitude} and
Figure~\ref{fig-compTfIdealPhase} compares the system-theoretic analysis
with transfer functions with the simulated idealised of the filter
design with each other. Both plots show very nicely that the amplitude
response as well as the phase response of all four filters match with
their simulated and calculated responses.

\begin{codelisting}

\caption{\label{lst-comptfidealamplitude}}

\centering{

\begin{Shaded}
\begin{Highlighting}[]
\ImportTok{import}\NormalTok{ numpy }\ImportTok{as}\NormalTok{ np}
\ImportTok{import}\NormalTok{ matplotlib.pyplot }\ImportTok{as}\NormalTok{ plt}
\ImportTok{import}\NormalTok{ sys}
\NormalTok{sys.path.insert(}\DecValTok{0}\NormalTok{, }\StringTok{\textquotesingle{}../simulation\textquotesingle{}}\NormalTok{)}
\ImportTok{import}\NormalTok{ ltspy3}

\NormalTok{sd}\OperatorTok{=}\NormalTok{ltspy3.SimData(}\StringTok{\textquotesingle{}../simulation/biquad\_univ.raw\textquotesingle{}}\NormalTok{)}

\NormalTok{nvoutLPF }\OperatorTok{=}\NormalTok{ sd.variables.index(}\StringTok{b\textquotesingle{}v(lpf)\textquotesingle{}}\NormalTok{)}
\NormalTok{nvoutHPF }\OperatorTok{=}\NormalTok{ sd.variables.index(}\StringTok{b\textquotesingle{}v(hpf)\textquotesingle{}}\NormalTok{)}
\NormalTok{nvoutBPF }\OperatorTok{=}\NormalTok{ sd.variables.index(}\StringTok{b\textquotesingle{}v(bpf)\textquotesingle{}}\NormalTok{)}
\NormalTok{nvoutBSF }\OperatorTok{=}\NormalTok{ sd.variables.index(}\StringTok{b\textquotesingle{}v(bsf)\textquotesingle{}}\NormalTok{)}
\NormalTok{nfrequency }\OperatorTok{=}\NormalTok{ sd.variables.index(}\StringTok{b\textquotesingle{}frequency\textquotesingle{}}\NormalTok{)}

\CommentTok{\#behauvioural moddling with tfs}

\CommentTok{\# Initial values}
\NormalTok{f0 }\OperatorTok{=} \DecValTok{1000}  \CommentTok{\# Resonance frequency in Hz}
\NormalTok{w0 }\OperatorTok{=} \DecValTok{2} \OperatorTok{*}\NormalTok{ np.pi }\OperatorTok{*}\NormalTok{ f0  }\CommentTok{\# Angular frequency in rad/s}
\NormalTok{Q }\OperatorTok{=} \DecValTok{10}  \CommentTok{\# Quality factor}
\NormalTok{H0 }\OperatorTok{=} \DecValTok{1}  \CommentTok{\# Play around with this later}

\CommentTok{\# Logarithmic frequency axis}
\NormalTok{frequencies }\OperatorTok{=}\NormalTok{ np.logspace(}\DecValTok{0}\NormalTok{, }\DecValTok{5}\NormalTok{, }\DecValTok{10000}\NormalTok{)  }\CommentTok{\# Frequency from 10\^{}2 to 10\^{}4 Hz}
\NormalTok{s }\OperatorTok{=} \OtherTok{1j} \OperatorTok{*} \DecValTok{2} \OperatorTok{*}\NormalTok{ np.pi }\OperatorTok{*}\NormalTok{ frequencies  }\CommentTok{\# Laplace{-}Variable s = jω}

\CommentTok{\#\#\#\#\#\#\#\#\#\#\#\#\#\#\#\#\#\#\#\#\#\#\#\#\#\#\#\#\#\#\#\#\#\#\#\#\#\#\#\#\#\#\#\#}
\CommentTok{\# Transfer functions of Active Filters}
\CommentTok{\#\#\#\#\#\#\#\#\#\#\#\#\#\#\#\#\#\#\#\#\#\#\#\#\#\#\#\#\#\#\#\#\#\#\#\#\#\#\#\#\#\#\#\#}


\CommentTok{\#\#\# Numerator}
\CommentTok{\# Low Pass Filter}
\NormalTok{b\_lp }\OperatorTok{=}\NormalTok{ H0}

\CommentTok{\# High Pass Filter}
\NormalTok{b\_hp }\OperatorTok{=}\NormalTok{ (H0 }\OperatorTok{*}\NormalTok{ (s}\OperatorTok{**}\DecValTok{2} \OperatorTok{/}\NormalTok{ w0}\OperatorTok{**}\DecValTok{2}\NormalTok{))}

\CommentTok{\# Band Pass Filter}
\NormalTok{b\_bp }\OperatorTok{=}\NormalTok{ (}\OperatorTok{{-}}\NormalTok{H0 }\OperatorTok{*}\NormalTok{ (s }\OperatorTok{/}\NormalTok{ w0))}

\CommentTok{\# Band Stop Filter}
\NormalTok{b\_bs }\OperatorTok{=} \OperatorTok{{-}}\NormalTok{(}\DecValTok{1} \OperatorTok{+}\NormalTok{ (s}\OperatorTok{**}\DecValTok{2} \OperatorTok{/}\NormalTok{ (w0}\OperatorTok{**}\DecValTok{2}\NormalTok{))) }\OperatorTok{*}\NormalTok{ H0}

\CommentTok{\# Denominator {-}\textgreater{} for all filters the same}
\NormalTok{a0 }\OperatorTok{=} \DecValTok{1}
\NormalTok{a1 }\OperatorTok{=}\NormalTok{ (s }\OperatorTok{/}\NormalTok{ (w0 }\OperatorTok{*}\NormalTok{ Q))}
\NormalTok{a2 }\OperatorTok{=}\NormalTok{ (s}\OperatorTok{**}\DecValTok{2} \OperatorTok{/}\NormalTok{ (w0}\OperatorTok{**}\DecValTok{2}\NormalTok{))}

\NormalTok{den }\OperatorTok{=}\NormalTok{ a0 }\OperatorTok{+}\NormalTok{ a1 }\OperatorTok{+}\NormalTok{ a2}

\CommentTok{\#\#\#\#\#\#\#\#\#\#\#\#\#\#\#\#\#\#\#\#\#\#\#\#\#\#\#\#\#\#\#\#\#\#\#\#\#\#\#\#\#\#\#\#}
\CommentTok{\# Calculation of the transfer functions H(s)}
\CommentTok{\#\#\#\#\#\#\#\#\#\#\#\#\#\#\#\#\#\#\#\#\#\#\#\#\#\#\#\#\#\#\#\#\#\#\#\#\#\#\#\#\#\#\#\#}
\NormalTok{Hs\_lp }\OperatorTok{=}\NormalTok{ b\_lp }\OperatorTok{/}\NormalTok{ den}
\NormalTok{Hs\_hp }\OperatorTok{=}\NormalTok{ b\_hp }\OperatorTok{/}\NormalTok{ den}
\NormalTok{Hs\_bp }\OperatorTok{=}\NormalTok{ b\_bp }\OperatorTok{/}\NormalTok{ den}
\NormalTok{Hs\_bs }\OperatorTok{=}\NormalTok{ b\_bs }\OperatorTok{/}\NormalTok{ den}

\CommentTok{\#mag plots}

\NormalTok{fig, axs }\OperatorTok{=}\NormalTok{ plt.subplots(}\DecValTok{2}\NormalTok{, }\DecValTok{2}\NormalTok{, sharex}\OperatorTok{=}\VariableTok{True}\NormalTok{, sharey}\OperatorTok{=}\VariableTok{True}\NormalTok{)}

\NormalTok{ax1, ax2, ax3, ax4 }\OperatorTok{=}\NormalTok{ axs.flatten()}

\CommentTok{\# Plotting each}
\NormalTok{ax1.semilogx(frequencies, }\DecValTok{20} \OperatorTok{*}\NormalTok{ np.log10(np.}\BuiltInTok{abs}\NormalTok{(Hs\_lp)), label}\OperatorTok{=}\StringTok{\textquotesingle{}tf\textquotesingle{}}\NormalTok{)}
\NormalTok{ax1.semilogx(sd.values[nfrequency], }\DecValTok{20} \OperatorTok{*}\NormalTok{ np.log10(}\BuiltInTok{abs}\NormalTok{(sd.values[nvoutLPF])), label}\OperatorTok{=}\StringTok{\textquotesingle{}ngspice\textquotesingle{}}\NormalTok{)}
\NormalTok{ax1.set\_ylabel(}\StringTok{"magnitude/dB"}\NormalTok{)}
\NormalTok{ax1.set\_title(}\StringTok{"low pass filter"}\NormalTok{)}
\NormalTok{ax1.set\_xlim([}\DecValTok{10}\NormalTok{,}\FloatTok{10e4}\NormalTok{])}
\NormalTok{ax1.set\_ylim([}\OperatorTok{{-}}\DecValTok{60}\NormalTok{,}\DecValTok{25}\NormalTok{])}
\NormalTok{ax1.grid(}\VariableTok{True}\NormalTok{, which}\OperatorTok{=}\StringTok{"both"}\NormalTok{, ls}\OperatorTok{=}\StringTok{"{-}{-}"}\NormalTok{)}
\NormalTok{ax1.legend()}

\NormalTok{ax2.semilogx(frequencies, }\DecValTok{20} \OperatorTok{*}\NormalTok{ np.log10(np.}\BuiltInTok{abs}\NormalTok{(Hs\_hp)), label}\OperatorTok{=}\StringTok{\textquotesingle{}tf\textquotesingle{}}\NormalTok{)}
\NormalTok{ax2.semilogx(sd.values[nfrequency], }\DecValTok{20} \OperatorTok{*}\NormalTok{ np.log10(}\BuiltInTok{abs}\NormalTok{(sd.values[nvoutHPF])), label}\OperatorTok{=}\StringTok{\textquotesingle{}ngspice\textquotesingle{}}\NormalTok{)}
\NormalTok{ax2.set\_title(}\StringTok{"high pass filter"}\NormalTok{)}
\CommentTok{\#ax2.set\_xlim([10,10e3])}
\CommentTok{\#ax2.set\_ylim([{-}40,25])}
\NormalTok{ax2.grid(}\VariableTok{True}\NormalTok{, which}\OperatorTok{=}\StringTok{"both"}\NormalTok{, ls}\OperatorTok{=}\StringTok{"{-}{-}"}\NormalTok{)}
\NormalTok{ax2.legend()}

\NormalTok{ax3.semilogx(frequencies, }\DecValTok{20} \OperatorTok{*}\NormalTok{ np.log10(np.}\BuiltInTok{abs}\NormalTok{(Hs\_bp)), label}\OperatorTok{=}\StringTok{\textquotesingle{}tf\textquotesingle{}}\NormalTok{)}
\NormalTok{ax3.semilogx(sd.values[nfrequency], }\DecValTok{20} \OperatorTok{*}\NormalTok{ np.log10(}\BuiltInTok{abs}\NormalTok{(sd.values[nvoutBPF])), label}\OperatorTok{=}\StringTok{\textquotesingle{}ngspice\textquotesingle{}}\NormalTok{)}
\NormalTok{ax3.set\_xlabel(}\StringTok{"frequency/Hz"}\NormalTok{)}
\NormalTok{ax3.set\_ylabel(}\StringTok{"magnitude/dB"}\NormalTok{)}
\NormalTok{ax3.set\_title(}\StringTok{"band pass filter"}\NormalTok{)}
\CommentTok{\#ax3.set\_xlim([10,10e3])}
\CommentTok{\#ax3.set\_ylim([{-}40,25])}
\NormalTok{ax3.grid(}\VariableTok{True}\NormalTok{, which}\OperatorTok{=}\StringTok{"both"}\NormalTok{, ls}\OperatorTok{=}\StringTok{"{-}{-}"}\NormalTok{)}
\NormalTok{ax3.legend()}

\NormalTok{ax4.semilogx(frequencies, }\DecValTok{20} \OperatorTok{*}\NormalTok{ np.log10(np.}\BuiltInTok{abs}\NormalTok{(Hs\_bs)), label}\OperatorTok{=}\StringTok{\textquotesingle{}tf\textquotesingle{}}\NormalTok{)}
\NormalTok{ax4.semilogx(sd.values[nfrequency], }\DecValTok{20} \OperatorTok{*}\NormalTok{ np.log10(}\BuiltInTok{abs}\NormalTok{(sd.values[nvoutBSF])), label}\OperatorTok{=}\StringTok{\textquotesingle{}ngspice\textquotesingle{}}\NormalTok{)}
\NormalTok{ax4.set\_xlabel(}\StringTok{"frequency/Hz"}\NormalTok{)}
\NormalTok{ax4.set\_title(}\StringTok{"band stop filter"}\NormalTok{)}
\CommentTok{\#ax4.set\_xlim([10,10e3])}
\CommentTok{\#ax4.set\_ylim([{-}40,25])}
\NormalTok{ax4.grid(}\VariableTok{True}\NormalTok{, which}\OperatorTok{=}\StringTok{"both"}\NormalTok{, ls}\OperatorTok{=}\StringTok{"{-}{-}"}\NormalTok{)}
\NormalTok{ax4.legend()}

\NormalTok{plt.tight\_layout()}
\CommentTok{\#plt.show()}

\NormalTok{plt.savefig(}\StringTok{"../images/sec\_characterisation/compTfIdealAmplitude.png"}\NormalTok{, }\BuiltInTok{format}\OperatorTok{=}\StringTok{"png"}\NormalTok{, dpi}\OperatorTok{=}\DecValTok{1000}\NormalTok{)}
\end{Highlighting}
\end{Shaded}

}

\end{codelisting}%

\begin{figure}

\centering{

\pandocbounded{\includegraphics[keepaspectratio]{content/../images/sec_characterisation/compTfIdealAmplitude.png}}

}

\caption{\label{fig-compTfIdealAmplitude}Comparison of the amplitude
response between the transfer functions and idealised circuit}

\end{figure}%

\begin{codelisting}

\caption{\label{lst-comptfidealphase}}

\centering{

\begin{Shaded}
\begin{Highlighting}[]
\ImportTok{import}\NormalTok{ numpy }\ImportTok{as}\NormalTok{ np}
\ImportTok{import}\NormalTok{ matplotlib.pyplot }\ImportTok{as}\NormalTok{ plt}
\ImportTok{import}\NormalTok{ sys}
\NormalTok{sys.path.insert(}\DecValTok{0}\NormalTok{, }\StringTok{\textquotesingle{}../simulation\textquotesingle{}}\NormalTok{)}
\ImportTok{import}\NormalTok{ ltspy3}

\NormalTok{sd}\OperatorTok{=}\NormalTok{ltspy3.SimData(}\StringTok{\textquotesingle{}../simulation/biquad\_univ.raw\textquotesingle{}}\NormalTok{)}

\NormalTok{nvoutLPF }\OperatorTok{=}\NormalTok{ sd.variables.index(}\StringTok{b\textquotesingle{}v(lpf)\textquotesingle{}}\NormalTok{)}
\NormalTok{nvoutHPF }\OperatorTok{=}\NormalTok{ sd.variables.index(}\StringTok{b\textquotesingle{}v(hpf)\textquotesingle{}}\NormalTok{)}
\NormalTok{nvoutBPF }\OperatorTok{=}\NormalTok{ sd.variables.index(}\StringTok{b\textquotesingle{}v(bpf)\textquotesingle{}}\NormalTok{)}
\NormalTok{nvoutBSF }\OperatorTok{=}\NormalTok{ sd.variables.index(}\StringTok{b\textquotesingle{}v(bsf)\textquotesingle{}}\NormalTok{)}
\NormalTok{nfrequency }\OperatorTok{=}\NormalTok{ sd.variables.index(}\StringTok{b\textquotesingle{}frequency\textquotesingle{}}\NormalTok{)}

\CommentTok{\#behauvioural moddling with tfs}

\CommentTok{\# Initial values}
\NormalTok{f0 }\OperatorTok{=} \DecValTok{1000}  \CommentTok{\# Resonance frequency in Hz}
\NormalTok{w0 }\OperatorTok{=} \DecValTok{2} \OperatorTok{*}\NormalTok{ np.pi }\OperatorTok{*}\NormalTok{ f0  }\CommentTok{\# Angular frequency in rad/s}
\NormalTok{Q }\OperatorTok{=} \DecValTok{10}  \CommentTok{\# Quality factor}
\NormalTok{H0 }\OperatorTok{=} \DecValTok{1}  \CommentTok{\# Play around with this later}

\CommentTok{\# Logarithmic frequency axis}
\NormalTok{frequencies }\OperatorTok{=}\NormalTok{ np.logspace(}\DecValTok{0}\NormalTok{, }\DecValTok{4}\NormalTok{, }\DecValTok{10000}\NormalTok{)  }\CommentTok{\# Frequency from 10\^{}2 to 10\^{}4 Hz}
\NormalTok{s }\OperatorTok{=} \OtherTok{1j} \OperatorTok{*} \DecValTok{2} \OperatorTok{*}\NormalTok{ np.pi }\OperatorTok{*}\NormalTok{ frequencies  }\CommentTok{\# Laplace{-}Variable s = jω}

\CommentTok{\#\#\#\#\#\#\#\#\#\#\#\#\#\#\#\#\#\#\#\#\#\#\#\#\#\#\#\#\#\#\#\#\#\#\#\#\#\#\#\#\#\#\#\#}
\CommentTok{\# Transfer functions of Active Filters}
\CommentTok{\#\#\#\#\#\#\#\#\#\#\#\#\#\#\#\#\#\#\#\#\#\#\#\#\#\#\#\#\#\#\#\#\#\#\#\#\#\#\#\#\#\#\#\#}


\CommentTok{\#\#\# Numerator}
\CommentTok{\# Low Pass Filter}
\NormalTok{b\_lp }\OperatorTok{=}\NormalTok{ H0}

\CommentTok{\# High Pass Filter}
\NormalTok{b\_hp }\OperatorTok{=}\NormalTok{ (H0 }\OperatorTok{*}\NormalTok{ (s}\OperatorTok{**}\DecValTok{2} \OperatorTok{/}\NormalTok{ w0}\OperatorTok{**}\DecValTok{2}\NormalTok{))}

\CommentTok{\# Band Pass Filter}
\NormalTok{b\_bp }\OperatorTok{=}\NormalTok{ (}\OperatorTok{{-}}\NormalTok{H0 }\OperatorTok{*}\NormalTok{ (s }\OperatorTok{/}\NormalTok{ w0))}

\CommentTok{\# Band Stop Filter}
\NormalTok{b\_bs }\OperatorTok{=} \OperatorTok{{-}}\NormalTok{(}\DecValTok{1} \OperatorTok{+}\NormalTok{ (s}\OperatorTok{**}\DecValTok{2} \OperatorTok{/}\NormalTok{ (w0}\OperatorTok{**}\DecValTok{2}\NormalTok{))) }\OperatorTok{*}\NormalTok{ H0}

\CommentTok{\# Denominator {-}\textgreater{} for all filters the same}
\NormalTok{a0 }\OperatorTok{=} \DecValTok{1}
\NormalTok{a1 }\OperatorTok{=}\NormalTok{ (s }\OperatorTok{/}\NormalTok{ (w0 }\OperatorTok{*}\NormalTok{ Q))}
\NormalTok{a2 }\OperatorTok{=}\NormalTok{ (s}\OperatorTok{**}\DecValTok{2} \OperatorTok{/}\NormalTok{ (w0}\OperatorTok{**}\DecValTok{2}\NormalTok{))}

\NormalTok{den }\OperatorTok{=}\NormalTok{ a0 }\OperatorTok{+}\NormalTok{ a1 }\OperatorTok{+}\NormalTok{ a2}

\CommentTok{\#\#\#\#\#\#\#\#\#\#\#\#\#\#\#\#\#\#\#\#\#\#\#\#\#\#\#\#\#\#\#\#\#\#\#\#\#\#\#\#\#\#\#\#}
\CommentTok{\# Calculation of the transfer functions H(s)}
\CommentTok{\#\#\#\#\#\#\#\#\#\#\#\#\#\#\#\#\#\#\#\#\#\#\#\#\#\#\#\#\#\#\#\#\#\#\#\#\#\#\#\#\#\#\#\#}
\NormalTok{Hs\_lp }\OperatorTok{=}\NormalTok{ b\_lp }\OperatorTok{/}\NormalTok{ den}
\NormalTok{Hs\_hp }\OperatorTok{=}\NormalTok{ b\_hp }\OperatorTok{/}\NormalTok{ den}
\NormalTok{Hs\_bp }\OperatorTok{=}\NormalTok{ b\_bp }\OperatorTok{/}\NormalTok{ den}
\NormalTok{Hs\_bs }\OperatorTok{=}\NormalTok{ b\_bs }\OperatorTok{/}\NormalTok{ den}

\CommentTok{\#phase plots}

\NormalTok{fig, axs }\OperatorTok{=}\NormalTok{ plt.subplots(}\DecValTok{2}\NormalTok{, }\DecValTok{2}\NormalTok{, sharex}\OperatorTok{=}\VariableTok{True}\NormalTok{, sharey}\OperatorTok{=}\VariableTok{True}\NormalTok{)}

\NormalTok{ax1, ax2, ax3, ax4 }\OperatorTok{=}\NormalTok{ axs.flatten()}

\CommentTok{\# Plotting each}
\NormalTok{ax1.semilogx(frequencies, np.angle(Hs\_lp, deg}\OperatorTok{=}\VariableTok{True}\NormalTok{), label}\OperatorTok{=}\StringTok{\textquotesingle{}tf\textquotesingle{}}\NormalTok{)}
\NormalTok{ax1.semilogx(sd.values[nfrequency],np.angle(sd.values[nvoutLPF], deg}\OperatorTok{=}\VariableTok{True}\NormalTok{),label}\OperatorTok{=}\StringTok{\textquotesingle{}ngspice\textquotesingle{}}\NormalTok{)}
\NormalTok{ax1.set\_ylabel(}\StringTok{"phase/deg"}\NormalTok{)}
\NormalTok{ax1.set\_title(}\StringTok{"low pass filter"}\NormalTok{)}
\NormalTok{ax1.set\_xlim([}\DecValTok{10}\NormalTok{,}\FloatTok{10e3}\NormalTok{])}
\CommentTok{\#ax1.set\_ylim([{-}40,25])}
\NormalTok{ax1.grid(}\VariableTok{True}\NormalTok{, which}\OperatorTok{=}\StringTok{"both"}\NormalTok{, ls}\OperatorTok{=}\StringTok{"{-}{-}"}\NormalTok{)}
\NormalTok{ax1.legend()}

\NormalTok{ax2.semilogx(frequencies, np.angle(Hs\_hp, deg}\OperatorTok{=}\VariableTok{True}\NormalTok{), label}\OperatorTok{=}\StringTok{\textquotesingle{}tf\textquotesingle{}}\NormalTok{)}
\NormalTok{ax2.semilogx(sd.values[nfrequency],np.angle(sd.values[nvoutHPF], deg}\OperatorTok{=}\VariableTok{True}\NormalTok{),label}\OperatorTok{=}\StringTok{\textquotesingle{}ngspice\textquotesingle{}}\NormalTok{)}
\NormalTok{ax2.set\_title(}\StringTok{"high pass filter"}\NormalTok{)}
\CommentTok{\#ax1.set\_xlim([10,10e3])}
\CommentTok{\#ax1.set\_ylim([{-}40,25])}
\NormalTok{ax2.grid(}\VariableTok{True}\NormalTok{, which}\OperatorTok{=}\StringTok{"both"}\NormalTok{, ls}\OperatorTok{=}\StringTok{"{-}{-}"}\NormalTok{)}
\NormalTok{ax2.legend()}

\NormalTok{ax3.semilogx(frequencies, np.angle(Hs\_bp, deg}\OperatorTok{=}\VariableTok{True}\NormalTok{), label}\OperatorTok{=}\StringTok{\textquotesingle{}tf\textquotesingle{}}\NormalTok{)}
\NormalTok{ax3.semilogx(sd.values[nfrequency],np.angle(sd.values[nvoutBPF], deg}\OperatorTok{=}\VariableTok{True}\NormalTok{),label}\OperatorTok{=}\StringTok{\textquotesingle{}ngspice\textquotesingle{}}\NormalTok{)}
\NormalTok{ax3.set\_xlabel(}\StringTok{"frequency/Hz"}\NormalTok{)}
\NormalTok{ax3.set\_ylabel(}\StringTok{"phase/deg"}\NormalTok{)}
\NormalTok{ax3.set\_title(}\StringTok{"band pass filter"}\NormalTok{)}
\CommentTok{\#ax1.set\_xlim([10,10e3])}
\CommentTok{\#ax1.set\_ylim([{-}40,25])}
\NormalTok{ax3.grid(}\VariableTok{True}\NormalTok{, which}\OperatorTok{=}\StringTok{"both"}\NormalTok{, ls}\OperatorTok{=}\StringTok{"{-}{-}"}\NormalTok{)}
\NormalTok{ax3.legend()}

\NormalTok{ax4.semilogx(frequencies, np.angle(Hs\_bs, deg}\OperatorTok{=}\VariableTok{True}\NormalTok{), label}\OperatorTok{=}\StringTok{\textquotesingle{}tf\textquotesingle{}}\NormalTok{)}
\NormalTok{ax4.semilogx(sd.values[nfrequency],np.angle(sd.values[nvoutBSF], deg}\OperatorTok{=}\VariableTok{True}\NormalTok{),label}\OperatorTok{=}\StringTok{\textquotesingle{}ngspice\textquotesingle{}}\NormalTok{)}
\NormalTok{ax4.set\_xlabel(}\StringTok{"frequency/Hz"}\NormalTok{)}
\NormalTok{ax4.set\_title(}\StringTok{"band stop filter"}\NormalTok{)}
\CommentTok{\#ax1.set\_xlim([10,10e3])}
\CommentTok{\#ax1.set\_ylim([{-}40,25])}
\NormalTok{ax4.grid(}\VariableTok{True}\NormalTok{, which}\OperatorTok{=}\StringTok{"both"}\NormalTok{, ls}\OperatorTok{=}\StringTok{"{-}{-}"}\NormalTok{)}
\NormalTok{ax4.legend()}

\NormalTok{plt.tight\_layout()}
\CommentTok{\#plt.show()}

\NormalTok{plt.savefig(}\StringTok{"../images/sec\_characterisation/compTfIdealPhase.png"}\NormalTok{, }\BuiltInTok{format}\OperatorTok{=}\StringTok{"png"}\NormalTok{, dpi}\OperatorTok{=}\DecValTok{1000}\NormalTok{)}
\end{Highlighting}
\end{Shaded}

}

\end{codelisting}%

\begin{figure}

\centering{

\pandocbounded{\includegraphics[keepaspectratio]{content/../images/sec_characterisation/compTfIdealPhase.png}}

}

\caption{\label{fig-compTfIdealPhase}Comparison of the phase response
between the transfer functions and idealised circuit}

\end{figure}%

\bookmarksetup{startatroot}

\chapter{Xschem Simulation}\label{xschem-simulation}

\section{Environment Setup}\label{environment-setup}

Proper configuration of the simulation environment is crucial for
accurate and reliable simulation results using Xschem within a Docker
container. This section details the necessary considerations regarding
path handling and Docker environment variables for a robust simulation
process.

\subsection{Absolute vs.~Relative Path
Handling}\label{absolute-vs.-relative-path-handling}

In integrated circuit simulations, particularly involving hierarchical
netlisting used by Xschem, careful management of file paths is
essential. Absolute paths were chosen to ensure consistency and
reliability, as hierarchical structures in netlisting frequently
generate new subdirectories during simulation processes. Using absolute
paths prevents issues arising from relative paths becoming invalid when
directory structures change. Relative paths, although theoretically
functional, add complexity and increase the likelihood of path errors
during simulations.

Therefore, all simulation paths have been consistently defined using
absolute references, guaranteeing stable and predictable netlisting and
simulation behavior throughout the project.

\subsection{Docker Environment Variables
Configuration}\label{docker-environment-variables-configuration}

Within the Docker container environment used for simulations, specific
variables require explicit and careful configuration. Key points
regarding Docker environment setup include:

\begin{itemize}
\item
  \textbf{DESIGNS Variable}:\\
  The core Docker environment variable, \texttt{\$DESIGNS}, is
  configured to reference the directory \texttt{root/FOSS/designs}.
  Local project paths and symbols are added explicitly to this variable
  to ensure accurate referencing and inclusion during netlisting.
\item
  \textbf{Docker Root Directory}:\\
  The root directory within the Docker environment differs from that of
  the host operating system, specifically being located at
  \texttt{root/headless}. This distinction requires careful handling
  when navigating or setting paths within Docker terminals.
\item
  \textbf{Terminal and Directory Awareness}:\\
  Users, especially those working on Windows-based systems, need to
  distinguish clearly between terminal environments (native Windows,
  Windows Subsystem for Linux, Docker) as each has fundamentally
  different root directories and file path conventions. Maintaining
  clear awareness of these distinctions prevents file referencing errors
  during simulations.
\end{itemize}

These explicit configurations and considerations ensure that the Xschem
simulation environment operates smoothly, predictably, and accurately
throughout the project's lifecycle.

\section{Symbol and Netlist
Preparation}\label{symbol-and-netlist-preparation}

Accurate preparation of symbols and netlists is fundamental for reliable
integrated circuit simulations using Xschem. This section outlines the
critical aspects regarding symbol definitions, classifications, netlist
generation, and debugging practices to ensure proper functionality and
error-free simulations.

\subsection{Symbol Definitions and Classification (Analog
Pins)}\label{symbol-definitions-and-classification-analog-pins}

All symbols used within this simulation framework must be explicitly
classified as analog components. Digital pins or incorrectly labeled
pins lead to simulation failures because they are incompatible with the
analog circuit simulation environment. Furthermore, each symbol must be
explicitly marked as a sub-circuit rather than a primitive component.
Mislabeling symbols as primitives can cause errors during hierarchical
netlist processing, as the netlister expects defined sub-circuit symbols
to correctly navigate the circuit hierarchy.

In this project, all symbols were thoroughly reviewed to confirm they
were analog and properly labeled as sub-circuits.

\subsection{Netlist Generation and
Validation}\label{netlist-generation-and-validation}

Netlist generation in Xschem follows a hierarchical methodology.
Starting from the top schematic, the netlister recursively searches for
defined symbols globally. It is essential that symbol paths and local
variables are correctly configured and matched explicitly in each
hierarchy level. If paths or symbols are misconfigured or missing, the
netlister fails to locate necessary components, leading to simulation
errors.

Therefore, explicit steps were taken to ensure:

\begin{itemize}
\tightlist
\item
  Proper global variable definitions for symbol paths.
\item
  Accurate symbol descriptions that match netlist entries.
\item
  Verification of netlists after every significant schematic
  modification.
\end{itemize}

Effective debugging practices strongly emphasize using plain-text
editors to inspect and modify symbol (\texttt{.sym}) and schematic
(\texttt{.scm}) files. The files (\texttt{.spice}) making direct text
inspection both possible and advisable.

Adopting this debugging approach significantly enhanced the efficiency
and accuracy of symbol and netlist preparation processes throughout the
simulation.

\section{OTA Circuit Integration Self-built Five-Transistor
OTA}\label{sec-self-built-ota}

Initially, our group developed our own five-transistor OTA schematic
within Xschem to gain practical insights and explore various design
parameters.

Figure~\ref{fig-self-built-ota} shows the minimalist OTA designed during
the early project phase. The circuit follows the five-transistor
topology: an NMOS differential pair (M8/M3) converts the input
difference into two anti-phase drain currents; a PMOS current mirror
(M1/M7) then forces these currents to recombine at the single-ended node
v\_out, generating voltage gain of roughly
\[ g_{m,\text{pair}}\bigl(r_{o,\text{M1}} \parallel r_{o,\text{M7}}\bigr) \]
Biasing is supplied by the tail branch (M4→M2), where the external
i\_bias current is mirrored into M2.

\begin{figure}

\centering{

\includegraphics[width=0.8\linewidth,height=\textheight,keepaspectratio]{content/../images/sec_xschem/self_build_ota.png}

}

\caption{\label{fig-self-built-ota}Self-built five-transistor OTA
schematic.}

\end{figure}%

\subsection{Validation of Self-built OTA -- unity-gain
buffer}\label{sec-ota-ug}

To verify large-signal behaviour the prototype OTA was configured as a
unity-gain follower (Figure~\ref{fig-ota-ug-tb}): v\_out is shorted to
the inverting input, and the non-inverting input is driven by a single
rail-to-rail pulse PULSE(0 1.2 10m 1u 1u 20m 2) The stimulus holds 0 V
for 10 ms, rises to 1.2 V in 1 µs, remains high for 20 ms, and then
returns to 0 V.\\
The long two-second period guarantees only one transition inside the 100
ms simulation window.\\
A 20 µA current source on i\_bias establishes the tail current, with
supplies at VDD = 1.5 V and VSS = 0 V.

\begin{figure}

\centering{

\includegraphics[width=0.8\linewidth,height=\textheight,keepaspectratio]{content/../images/sec_xschem/tb_self_tran.png}

}

\caption{\label{fig-ota-ug-tb}Unity-gain test-bench schematic.}

\end{figure}%

The essential NGSpice control block is

\begin{Shaded}
\begin{Highlighting}[]
\NormalTok{.control}
\NormalTok{save all}
\NormalTok{tran 100u 100m}
\NormalTok{*set filetype=ascii}
\NormalTok{set wr\_singlescale}
\NormalTok{*wrdata tb\_ota\_tran\_UGbuffer.txt v(v\_out) v(v\_in)}
\NormalTok{.endc}
\end{Highlighting}
\end{Shaded}

\begin{figure}

\centering{

\includegraphics[width=0.8\linewidth,height=\textheight,keepaspectratio]{content/../images/sec_xschem/self_tran_vin_vout.png}

}

\caption{\label{fig-ota-ug-wave}Transient waveforms: v\_in vs.~v\_out.}

\end{figure}%

Figure~\ref{fig-ota-ug-wave} plots the resulting traces. The output
(left) mirrors the input (right). The output does respond to the input
change, it rises up during the pulse window, indicating the OTA is doing
something. V\_out lags significantly behind v\_in and only reaches
1.1\,V, while v\_in reaches 1.2\,V. The OTA as unity-gain buffer should
closely track the input with only minimal offset or delay. This
indicates: - Too low gain - Insufficient bandwidth - Or incorrect
biasing (e.g., not enough current to drive the output stage).

After conducting simulations and analyses, and discussing the results
with the professor, the conclusion was that the five-transistor OTA
provided by Prof.~Pretel offered superior performance and reliability
for our application. Thus, the decision was made to proceed with
Prof.~Pretl's well-established OTA design.

\section{Xschem OTA Description (Five-Transistor OTA by
Prof.~Pretl)}\label{xschem-ota-description-five-transistor-ota-by-prof.-pretl}

Prof.~Pretl's five-transistor OTA is an optimized implementation,
specifically tailored for efficient analog circuit integration using the
SG13G2 CMOS process. His design, depicted in Figure~\ref{fig-pretl-ota},
maintains simplicity while enhancing performance through thoughtful
sizing and layout practices. This OTA employs a standard architecture
comprising an NMOS differential input pair, a tail current source
transistor, and a PMOS current mirror load. The input differential
voltage is transformed by the differential pair into two opposite-phase
currents, recombined by the PMOS mirror at the single-ended output node.

A dedicated transistor sets a stable bias current, ensuring predictable
and robust OTA performance over varying operating conditions. Channel
lengths and widths are carefully chosen to balance between adequate
intrinsic gain, bandwidth, and noise characteristics. Particularly, the
transistor sizes follow good IC design practices, with lengths exceeding
the minimum dimensions to suppress short-channel effects and flicker
noise. Bulk connections are properly handled to avoid body effects,
typically by tying bulks directly to their corresponding source
terminals. The accompanying layout from Prof.~Pretl (
Figure~\ref{fig-pretl-ota}) illustrates thoughtful device placement and
symmetry, reducing mismatches and parasitic coupling. Furthermore,
multiple finger transistors are employed effectively to optimize layout
density and performance, adhering to industry-standard guidelines.

Simulations and measured data presented by Prof.~Pretl validate this
design as a stable, efficient, and high-performance option for analog
signal processing. It exhibits superior gain and bandwidth compared to
simpler OTA configurations, making it suitable for more demanding
applications in integrated analog and mixed-signal circuits. Due to
these advantages, our group adopted this OTA as the baseline for
subsequent design and integration phases in our project.

\begin{figure}

\centering{

\includegraphics[width=0.8\linewidth,height=\textheight,keepaspectratio]{content/../images/sec_xschem/ota-improved.png}

}

\caption{\label{fig-pretl-ota}OTA schematic.}

\end{figure}%

\section{Sizing and Simulation}\label{sec-sizing-simulation}

Proper transistor sizing is critical when designing an OTA for use in a
biquad filter, where gain, linearity, and bandwidth all depend heavily
on MOSFET dimensions. To guide the sizing process, The \(g_m / I_D\)
methodology is adopted.

After experimenting with various combinations of transistor widths and
lengths in simulation, we found that many trade-offs emerged between
gain, speed, and voltage headroom---particularly for the cascode
devices. Following Prof.~Pretl's recommendations, assigning a reduced
channel length of \(L = 0.5u\) to the input differential pair and
cascode devices (\(M_{1/1C}\), \(M_{2/2C}\), \(M_{3/3C}\), \(M_{4/4C}\))
to maximize speed and transconductance. Meanwhile, the tail current
source transistors (\(M_5\), \(M_6\)) were given a longer channel length
of \(L = 5u\) to improve matching, suppress flicker noise, and ensure
robust common-mode rejection.

A constant sizing point of \(g_m / I_D = 13\) was applied across all
devices, which offered a good balance between transconductance
efficiency and voltage headroom. This choice ensured that even with
stacked transistors, all devices remained in saturation across the
intended signal swing and supply range.

Simulation results based on this sizing show that the OTA achieves an
open-loop gain \(A_0 > 43\,dB\), satisfying key performance goals such
as bandwidth and linear range for biquad operation. These results
confirm that the selected sizing strategy works reliably within the
supply voltage limits of the SG13G2 process, and is well-suited for
analog signal processing tasks requiring accurate and programmable
transconductance.

\subsection{Small-Signal Frequency
Response}\label{small-signal-frequency-response}

To evaluate the frequency-domain behavior of the improved OTA, an AC
simulation was performed using a differential testbench. The resulting
magnitude response, shown in Figure~\ref{fig-ota-improved-ac}, displays
the classic characteristics of a low-pass amplifier.

The curve remains flat (0\,dB) throughout the low-frequency range,
indicating that the OTA maintains constant gain for small-signal inputs
up to approximately \(f_{-3\text{dB}} \approx 80\,\text{MHz}\). This
--3\,dB point marks the unity-gain bandwidth of the amplifier in this
configuration. Beyond this corner frequency, the magnitude begins to
drop at a rate close to --20\,dB/decade, indicating a dominant
single-pole roll-off.

The high-frequency attenuation begins smoothly, confirming that the
amplifier is well-compensated and free from peaking or signs of
instability.

Overall, the OTA demonstrates a strong small-signal frequency response,
offering sufficient gain-bandwidth product for its intended use in
biquad filters and other analog signal processing tasks.

\begin{figure}

\centering{

\includegraphics[width=0.8\linewidth,height=\textheight,keepaspectratio]{content/../images/sec_xschem/ota-improved_tb-ac.png}

}

\caption{\label{fig-ota-improved-ac}AC magnitude response of the
improved OTA.}

\end{figure}%

\subsection{Large-Signal Transient
Response}\label{large-signal-transient-response}

Figure~\ref{fig-ota-improved-transient} shows the transient simulation
result of the improved OTA in response to a step input signal.

The red trace (\(v_\text{ena}\)) represents the input voltage. It
exhibits a sharp transition from 0\,V to approximately 1.5\,V at the
beginning of the simulation and remains constant for the rest of the
time window.

The blue trace (\(v_\text{out}\)) shows the corresponding OTA output
voltage. Immediately after the input step, the output rises to
approximately 0.8\,V and holds steady throughout the entire duration of
the simulation (15\,µs). There is no visible overshoot, ringing, or
delay in the output transition, and the response appears stable and
monotonic.

This result confirms that the OTA produces a constant output voltage in
response to a step input under the simulated conditions. The flat and
settled output also suggests that the biasing and feedback configuration
were correctly established.

\begin{figure}

\centering{

\includegraphics[width=0.8\linewidth,height=\textheight,keepaspectratio]{content/../images/sec_xschem/ota-improved_ota_tran.png}

}

\caption{\label{fig-ota-improved-transient}Transient response of the
improved OTA to a step input.}

\end{figure}%

\section{gm‑C Biquad Prototype -- Concept and
Implementation}\label{sec-gm-biquad}

After initially developing our own basic five-transistor OTA to explore
the design flow and simulation process in Xschem, the next step was, to
study and simulate several OTA topologies provided by Prof.~Pretl. This
included hands-on experimentation with his differential architectures
and cascode-enhanced designs, allowing us to deepen our understanding of
performance trade-offs and layout considerations in analog IC design.

In this final design step, the goal was to building our own OTA, this
time using an improved architecture specifically tailored for
integration into a gm‑C biquad filter. The goal was to implement a fully
differential, tunable low-pass filter using transconductors and
capacitors only. By controlling the OTA bias current, we aimed to
achieve continuous tuning of the filter's cut-off frequency without
modifying passive components.

The following section presents the concept, schematic implementation,
and simulation-based validation of this gm‑C biquad prototype. It marks
the transition from isolated OTA design toward the realization of a
complete analog signal processing block.

\subsection{OTA Implementation for gm‑C
Biquad}\label{ota-implementation-for-gmc-biquad}

The OTA used in our gm‑C biquad implementation is structurally inspired
by the differential architecture shown in Figure 3.15(c) of Baker (2010)
analog design textbook. As illustrated in
Figure~\ref{fig-ota-implemented}, the OTA consists of a fully
differential input stage, a folded current-mirror load, and a
programmable tail current source, making it well-suited for use as a
transconductor in gm‑C filters.

The core of the OTA includes a differential NMOS input pair (\(M_7\),
\(M_8\)), which converts the input voltage difference
\((v_\text{inp} - v_\text{inn})\) into a differential current. These
currents are then mirrored and steered by the PMOS current mirror loads
(\(M_2\) through \(M_6\)) into the output branches (\(v_\text{outp}\)
and \(v_\text{outn}\)). The top mirrors are folded downward into the
output stage, creating high output impedance and helping increase
gain---similar to the folded-cascode concept seen in the reference OTA.

The output common-mode voltage is stabilized using a biasing network
composed of transistors \(M_1\) and \(M_{11}\)--\(M_{12}\), which fix
the gate voltages for proper operation. The tail current is generated
through the NMOS branch (\(M_9\), \(M_{10}\)), which is set by the
external \(i_\text{bias}\) node and determines the OTA's
transconductance via:

\[
g_m \approx \frac{2 I_\text{bias}}{V_\text{ov}}
\]

The structure is deliberately symmetric to preserve linearity and
minimize even-order distortion, and the design reflects the need for
balanced performance across process corners in gm‑C filter applications.
Unlike the simplified example in Baker (2010), which omits layout- and
bias-stabilizing elements for clarity, our implementation explicitly
includes bias regulation to improve practical robustness and simulation
convergence.

\begin{figure}

\centering{

\includegraphics[width=0.8\linewidth,height=\textheight,keepaspectratio]{content/../images/sec_xschem/ota_gm_100u.png}

}

\caption{\label{fig-ota-implemented}Implemented OTA schematic used for
gm‑C biquad design.}

\end{figure}%

\subsection{gm‑C Biquad Schematic in
Xschem}\label{gmc-biquad-schematic-in-xschem}

Figure~\ref{fig-biquad} shows the complete schematic of the gm‑C biquad
filter implemented in Xschem. The design consists of four OTA instances,
each configured in a differential topology and biased with a shared tail
current of 100\,µA via the global \texttt{ibias} net. The power supply
is set to 1.5\,V, and all OTAs operate with a common-mode input
reference of 0.75\,V via the \texttt{vcm} node.

The circuit follows the classical gm‑C biquad topology introduced
earlier. It features two cascaded integrator stages formed by
OTA--capacitor pairs, followed by additional OTA blocks for feedforward
and feedback terms. All OTAs use the same internal schematic
(\texttt{ota\_gm\_100u}) and are referenced symmetrically.

Each OTA drives or is loaded by 10\,pF metal-insulator-metal (MIM)
capacitors (\(C_1\)--\(C_6\)), arranged in 1:2 ratios where necessary to
match the theoretical filter structure (e.g., \(C_1\) and \(C_2\)).
Differential symmetry is preserved throughout the network using mirrored
signal paths and balanced layout routing.

The input signals \texttt{vinp} and \texttt{vinn} are provided by AC
voltage sources with 0.5\,V DC offset and 0.5\,V amplitude. The output
is taken differentially as \texttt{voutp} and \texttt{voutn}, and the
simulation computes their difference using:

\begin{Shaded}
\begin{Highlighting}[]
\NormalTok{let vod = v(outp) {-} v(outn)}
\NormalTok{plot db(v(vod))}
\end{Highlighting}
\end{Shaded}

This differential output is evaluated across a frequency sweep from
1\,kHz to 100\,MHz using the NGSpice .ac dec command with 100 points per
decade:

\begin{Shaded}
\begin{Highlighting}[]
\NormalTok{ac dec 100 1k 100MEG}
\end{Highlighting}
\end{Shaded}

The NGSpice block at the top of the schematic (see
Figure~\ref{fig-biquad}) includes the device model definitions, the
simulation temperature (set to 27°C), and references to the SG13G2 PDK
libraries. Parameterized values for transistor dimensions are declared
at the beginning to allow flexibility when tuning device behavior:

\begin{Shaded}
\begin{Highlighting}[]
\NormalTok{.param lp="X"u wp="X"u l="X"u w="X"u}
\end{Highlighting}
\end{Shaded}

\begin{figure}

\centering{

\includegraphics[width=0.8\linewidth,height=\textheight,keepaspectratio]{content/../images/sec_xschem/biquad_gm_c.png}

}

\caption{\label{fig-biquad}gm‑C biquad design.}

\end{figure}%

This modular approach to simulation setup enhances reproducibility and
simplifies design iterations. The OTA instances are invoked via their
corresponding symbol files, ensuring consistency between simulation and
schematic representation.

Overall, this gm‑C biquad schematic represents a practical and scalable
implementation of a continuous-time filter using custom OTA blocks. The
approach aligns well with standard analog IC design methodology and is
suitable for integration in high-speed signal processing chains.

\subsection{Sizing the 1 kHz Low-Pass gm-C Biquad to the Baker
Topology}\label{sec-sizing-baker}

Figure~\ref{fig-biquad} reproduces the fully differential gm-C biquad
proposed by Baker (2010). In this architecture, the behavior of the
filter is governed by six small-signal parameters:

\[
\begin{aligned}
G_1 &= \frac{g_{m1}}{C_1 + C_2}, &
G_2 &= \frac{g_{m2}}{g_{m1}}, &
G_3 &= \frac{C_1}{g_{m1}}, \\
G_4 &= \frac{g_{m3}}{C_3 + C_4}, &
G_5 &= \frac{g_{m4}}{g_{m1}}, &
G_6 &= \frac{C_3}{g_{m3}}.
\end{aligned}
\]

These expressions relate the two integrator stages
\(\bigl(g_{m1}, C_{1,2}\bigr)\) and \(\bigl(g_{m3}, C_{3,4}\bigr)\) to
the feed-forward transconductors \(g_{m2}\) and \(g_{m4}\). For a
\textbf{low-pass} realisation, we adopt the convenient symmetry:

\[
g_{m1} = g_{m3} = g_m, \quad C_1 = C_2 = C_3 = C_4 = C,
\]

so that both integrator poles coincide. Under this symmetry, the pole
frequency reduces to

\[
f_0 = \frac{g_m}{4\pi C},
\]

where the factor of 4 arises from the doubled \(2C\) plates used at each
integrator output in the schematic.

To meet the low-frequency target of 1\,kHz, we must carefully select
appropriate values for \(g_m\) and \(C\), while respecting the OTA bias
constraint. Each OTA operates with a fixed tail current of
\(I_{\text{tail}} = 100\;\mu\text{A}\), resulting in
\(I_D = 50\;\mu\text{A}\) per branch. To ensure good noise performance
and strong inversion operation, we target a transconductance efficiency
of approximately

\[
\left( \frac{g_m}{I_D} \right)_{\text{target}} \approx 5\;\text{V}^{-1}.
\]

This yields a branch transconductance of

\[
g_m = \left(\frac{g_m}{I_D}\right) I_D = 5 \cdot 50\;\mu\text{A} = 0.25\;\text{mS}.
\]

Substituting this into the expression for \(f_0\), we determine the
required capacitor value for a 1\,kHz pole:

\[
C = \frac{g_m}{4\pi f_0} = \frac{0.25\;\text{mS}}{4\pi \cdot 1\;\text{kHz}} \approx 20\;\text{nF}.
\]

As each integrator output is loaded with \(2C\), the effective
capacitance at each output node is approximately 40\,nF.

\begin{quote}
\textbf{Area note:} At a typical SG13G2 MIM density of
\(2\,\text{fF}/\mu\text{m}^2\), a 20\,nF capacitor requires roughly
10\,mm² of silicon area. This is large for an on-chip design and
motivates the use of stacked MIM structures or off-chip MLCCs for
prototyping.
\end{quote}

To realise this \(g_m\) in a compact layout, we choose a MOSFET
overdrive voltage of \(V_{\text{OV}} = 0.40\,V\) and assume a
mobility-capacitance product of
\(\mu_n C_{\text{ox}} \approx 200\,\mu\text{A}/\text{V}^2\). The
required transistor aspect ratio then becomes:

\[
\frac{W}{L} = \frac{2I_D}{\mu_n C_{\text{ox}} V_{\text{OV}}^2}
= \frac{100\;\mu\text{A}}{200\;\mu\text{A}/\text{V}^2 \cdot 0.16} \approx 3.1.
\]

The final device sizing is summarised below:

\begin{longtable}[]{@{}
  >{\raggedright\arraybackslash}p{(\linewidth - 6\tabcolsep) * \real{0.1529}}
  >{\raggedleft\arraybackslash}p{(\linewidth - 6\tabcolsep) * \real{0.0824}}
  >{\raggedleft\arraybackslash}p{(\linewidth - 6\tabcolsep) * \real{0.0824}}
  >{\raggedright\arraybackslash}p{(\linewidth - 6\tabcolsep) * \real{0.6824}}@{}}
\toprule\noalign{}
\begin{minipage}[b]{\linewidth}\raggedright
Device set
\end{minipage} & \begin{minipage}[b]{\linewidth}\raggedleft
\(L\)
\end{minipage} & \begin{minipage}[b]{\linewidth}\raggedleft
\(W\)
\end{minipage} & \begin{minipage}[b]{\linewidth}\raggedright
Comment
\end{minipage} \\
\midrule\noalign{}
\endhead
\bottomrule\noalign{}
\endlastfoot
NMOS (all) & 1\,µm & 3.2\,µm & \(W/L \approx 3\), yielding
\(g_m \approx 0.25\)\,mS \\
PMOS (all) & 1\,µm & 8\,µm & Width ×2.5 for hole mobility
compensation \\
Tail pair & 5\,µm & 3.2\,µm & Long \(L\) increases \(r_o\) and
suppresses \(1/f\) noise \\
\end{longtable}

In summary, this sizing strategy realises a 1\,kHz low-pass response
while respecting the 100\,µA tail bias in theory. The chosen transistor
dimensions enable consistent \(g_m\) across all OTA branches. Although
the large on-chip capacitors pose layout challenges, they can be
addressed using multilayer MIM stacks or external capacitors during
prototyping. The unified NMOS/PMOS geometries also facilitate efficient
common-centroid layout across the four OTAs that form the biquad filter.

\section{GM/ID Methodology}\label{gmid-methodology}

One of the first questions to ask in IC design is how small or how large
are the MOSFETs used in the circuits. MOSFETs can be used in saturation
mode or in the triode state (as well as in cut-off but this is not
relevant for us). When the FET is in saturation the drain current
\(I_{D}\) is controlled primarily by the gate-source voltage \(V_{GS}\).
In this case the drain-source voltage has a smaller impact on the drain
current. For the Transistor to work in saturation the drain-source
terminals need to be driven with a voltage high enough so this
``saturates'' the FET and the highest drain-current is achieved.

On the other hand if the voltage applied across the drain-source
contacts (on a NMOS for example) is relatively low (compared to the
voltage for saturation), the FET will operate in the so called triode
mode. In triode mode the drain-source voltage \(V_{DS}\) has a
fundamentally larger impact on the drain current then in the saturation
mode. (H. Pretl and Michael Koefinger 2025)

The Methodology to solve the question, asked at the beginning of this
chapter is the \(\frac{gm}{I_{D}}\) methodology which will introduce in
a moment. There are basically three MOSFET characteristics directly
describing the behaviour of it:

\begin{itemize}
\tightlist
\item
  \(\frac{g_m}{I_D}\) : Transconductance Efficiency
\item
  \(\frac{\omega_{T}}{f_T}\) : Transit frequency
\item
  \(\frac{g_m}{g_{ds}}\) : Intrinsic Gain
\end{itemize}

To understand the first characteristic for FETs, take a look at the
different operating points which depend on the applied voltages.
Whenever applying voltages to a FET in order to control a specific drain
current \(I_{D}\), the operattion in the FET in either weak inversion,
strong inversion or moderate inversion. This behaviour is controlled by
the Overdrive Voltage \(V_{OV}\) which is defined as the difference
between the gate-source voltage and the threshold voltage. To note this
small point the drain-current is controlled by the voltage between gate
and source. Whenever an nmos is not being used as a low-side switch or
amplifier or the pmos is being used as a low-side component problems can
arise.

\[
    V_{OV} = V_{GS} - V_{TH}
\]

Important to keep in mind that the threshold voltage isn't a magical
number that can be applied to every MOSFET, it rather depends on the
geometry (with W and L for example) and other factors. For the example
nmos given in the Analog Circuit Design IHP SG13G2 Devices Table by
Professor Pretl, the threshold voltage is 0.5V. Therefore the overdrive
voltage describes how ``much'' the gate-source voltage is above the
threshold of the FET. Depending on this overdrive voltage the circuit/ic
designer can apply different \(\frac{g_{M}}{I_{D}}\) values with the
unit {[}\(\frac{1}{V}\){]}. This unit is derived in the following way:

With \(g_{M}\) defined by:

\[
g_{M} = \frac{\partial I_D}{\partial V_{GS}}
\]

and \(I_{D}\) having the unit Ampere {[}A{]} and the voltage \(V_{GS}\):

\[
\frac{\frac{A}{V}}{A} = \frac{1}{V}
\]

Before continueing with the \(\frac{g_{M}}{I_{D}}\) methode, note that
there also is the \emph{square-law} model with which circuit designers
can design MOSFET circuits. This model is usually applicable for PCB
circuits and takes the situation into account where the MOSFET is driven
in the stong inversion state. The square-law model is being applied
assuming that the FET is operating in the ``linear'' or ``triode'' mode,
however on nanometer scale FETs (down to 130 nm with the IHP-SG13G2 PDK)
this model doesn't give us precise solutions anymore. Many effects like
parasitic capacitances alter the operational behaviour of the FET and
lead to the square-law model deviating afar from the real-world
behaviour in many situations (Alan Doolittle 2025).

The square-model drain-current behavior is being described by the
following formula:

\[
I_{D} = \frac{Z \cdot \overline{\mu_n} \cdot C_{OX}}{L} [(V_{GS} - V_{T}) \cdot V_{DS} - \frac{V_{DS}^2}{2}]
\]

with the two conditions:

\(0 \leq V_{DS} \leq V_{D_{SAT}}\) and \(V_{GS} \geq V_T\)

with following definitions:

\begin{itemize}
\tightlist
\item
  \(C_{OX} = \frac{\epsilon_{ox}}{x_{ox}}\)
\item
  Z = MOSFET width
\item
  L = MOSFET Channel length
\item
  \(V_T\) = Threshold voltage
\item
  \(\overline{\mu_n}\) = effective electron mobility
\end{itemize}

The threshold voltage is defined as:

\[
V_T = 2 \phi_F + \frac{\epsilon_s}{C_{OX}} \sqrt{\frac{2q N_{A}}{\epsilon} (2\phi_F)}
\]

with:

\(\phi_F\) being the Fermi Potential (surface potential) defined by:

\[
\phi_F = \frac{kT}{q} \cdot ln (\frac{N_A}{n_i})
\]

with \(N_A\) being the acceptor doping concentration and \(n_i\) being
the intrinsic carrier concentration. The term 2 \(\phi_F\) corresponds
to the surface potential required to achieve strong inversion.

For more details the reader can consult (Alan Doolittle 2025),(Boris
Murmann 2016) or (Silveira, Flandre, and Jespers 1996).

To illustrate the problems of the square-law model when designing MOSFET
circuits at nanometer scale will look at somwe e graphs visualizing it's
limitation. First of all let's look at the formulas for the square-law
when we want to achieve more performance with our FETs:

Transconductance Efficiency: \[
\frac{g_m}{I_D} \cong \frac{2}{V_{OV}}
\]

higher efficiencs here means more transconductance for the same drain
current.

Transit Frequency:

\[
\frac{g_M}{C_{gg}} \cong \frac{3}{2}\frac{\mu V_{OV}}{L^2}
\]

higher transit frequency for the same gate-capacitance.

Intrinsic Gain:

\[
\frac{g_m}{g_{ds}} \cong \frac{2}{\lambda V_{OV}}
\]

high transconductance (at same drain-current \(I_D\)) without higher
output conductance.

The square-law model completely fails in these cases when the MOSFET is
not operation in strong inversion. In moderate and weak inversion we are
forced to use a different mathmatical model , and the
\(\frac{g_M}{I_D}\) method is a really good starting point (Ross Walker
2017).

The following figures will show the deviation between square-law and
measurements as well as the \(\frac{g_M}{I_D}\) methodology:

So first of all when we use \(g_M\) and \(I_D\) we specify that for a
specific drain-current we get a specific transconductance, for example
with a \(\frac{g_M}{I_D}\) of 10 S/A we get 10 \(\mu\) S per 1 \(\mu\) A
of bias current. And depending on how ``much'' the transistor is
operating above it's threshold voltage \(V_{th}\) (basically the
Overdrive Voltage \(V_{OV}\) ) you get different inversion levels. From
weak inveresions for low overdrive voltages to moderate inversion when
operating at approximately \(V_{OV}\) = \(V_{th}\) to high inversion
when \(V_{OV}\) \textgreater{} \(V_{th}\).

With the square-law value for transconductance efficiency we completely
deviate with that approximation in weak and moderate inversion:

\begin{figure}

\centering{

\includegraphics[width=1\linewidth,height=\textheight,keepaspectratio]{content/../images/sec_xschem/ac_inversion_vs_vov.png}

}

\caption{\label{fig-inversion-vs-vov}Inversion level vs.~Overdrive
Voltage (Ross Walker 2017)}

\end{figure}%

Another deviation from square-law to real MOSFET behaviour can be seen
when we increase the gate-source voltage of the FET and measure the
drain-current. According to square-law formual for the drain-current the
current should just increase to the square with increasing gate-source
voltage. But by taking the square of the drain-current and increasing
\(V_{GS}\) we can see that the drain-current does not magically start
flowing above the threshold-voltage and also the behaviour is also not
linear (quadratically when not taking the square of the current):

\begin{figure}

\centering{

\includegraphics[width=0.8\linewidth,height=\textheight,keepaspectratio]{content/../images/sec_xschem/ac_id_vs_gs.png}

}

\caption{\label{fig-id_vs_gs}Drain Current over Gate-Source Voltage,
Simulation vs.~square-law (Ross Walker 2017)}

\end{figure}%

This simulation is done for a n-channel MOSFET with a drain-source
voltage of 1.8 V and a size of L = 180 nm and W = 5 \(\mu\)m.

The drain-current behaviour at sub-threshold gate voltages is completely
inaccurate for the square-law too, and the following graph visualizes
the limitation of the square-law at this point again:

\begin{figure}

\centering{

\includegraphics[width=0.9\linewidth,height=\textheight,keepaspectratio]{content/../images/sec_xschem/ac_gmid_vs_gs.png}

}

\caption{\label{fig-gmid_vs_gs}Drain Current over Gate-Source Voltage,
comparison between an nmos, a bjt and the square-law (Ross Walker 2017)}

\end{figure}%

These three examples show that the approach using square-law to size
MOSFETs is not sufficient when the transistor is operating in weak or
moderate inversion and when driving the FET (nmos for example) with a
low gate-source (or overdrive-) voltage. To cite Mr.~Walker on this
topic: ``This means that the square law equation (which assumes 100\%
drift current) does not work unless the gate overdrive is several
\(\frac{kT}{q}\), (Ross Walker 2017)''.

To conclude this, we can keep in mind that there is no simple formula
that can describe the drain-current behaviour in all situations and be
universally used. So using the \(\frac{g_m}{I_D}\) methodology is the
way to go in our project.

Now with that out of the way we can design our circuits using the
\(\frac{g_m}{I_D}\) methodology. The main properties of our MOSFETs we
can manipulate in xschem are the lenght of the channel L, the width W
and the bias current \(I_D\). The common way to use this method is to
first characterize nmos and pmos field effect transistors and then use
this data to design the circuits. In the chapter ``MOSFET
characterization Testbench'' chapter in (H. Pretl and Michael Koefinger
2025) we can see how the values for the \(\frac{g_m}{I_D}\) methodology
are being simulated for later use.

The lenght of the MOSFET channel also has a large influence on it's
frequency characteristic as it can be seen in this simulation:

\begin{figure}

\centering{

\includegraphics[width=0.8\linewidth,height=\textheight,keepaspectratio]{images/sec_xschem/ac_eaa_ft_vs_L.png}

}

\caption{\label{fig-ft_vs_L}transit frequency vs.~channel length L (H.
Pretl and Michael Koefinger 2025)}

\end{figure}%

The operating areas of interest for us are the saturation region (when
using the FET as an amplifier for example) and the region when the FET
is being used to ``just'' work as a switch. With setting \(V_{DS}\) to
\(\frac{V_{DD}}{2}\) we keep the FET in saturation. Reminding ourselves
again that with larger \(g_M\) we have more ``gain'' and with a smaller
\(I_D\) we have higher efficiency we try to hit the sweetspot between
size (as every square milimeter has it's cost) and current consumption
(if we have wearable battery powered devices for example). Keeping also
in mind that temperature has a large effect we cannot use arbitratily
large drain currents.

Following plot visualizes the dependancy of \(\frac{g_M}{I_D}\) to the
gate-source voltage and shows the transit frequency behaviour too:

\begin{figure}

\centering{

\includegraphics[width=0.8\linewidth,height=\textheight,keepaspectratio]{images/sec_xschem/ac_ft_vs_gm_vgs.png}

}

\caption{\label{fig-gm_vs_gs_vs_ft}\(\frac{g_M}{I_D}\) and \(f_T\) over
the gate-source voltage (H. Pretl and Michael Koefinger 2025)}

\end{figure}%

\subsection{GM/ID Overview}\label{gmid-overview}

Device sizing is a fundamental step in analog circuit design.Proper
sizing ensures the desired trade-off between gain, bandwidth, power
consumption, and linearity. A structured approach to sizing guarantees
that the transistors operate in their optimal region---typically strong
inversion and saturation for analog applications.

This section presents a basic OTA sizing methodology based on the work
of Prepl, whose notebook is publicly available at
\href{https://github.com/iic-jku/analog-circuit-design/blob/main/sizing/sizing_basic_ota.ipynb}{GitHub:
analog-circuit-design}.

The notebook demonstrates a Python-based framework for calculating
initial transistor dimensions in a basic OTA. Below is an explanation of
the key steps, equations, and important code snippets from the sizing
method, along with visual aids to clarify the process.

The sizing workflow follows these general steps:

\begin{enumerate}
\def\labelenumi{\arabic{enumi}.}
\tightlist
\item
  Define process parameters.
\item
  Specify bias currents and overdrive voltages.
\item
  Calculate width-to-length ratios (W/L) for each transistor.
\item
  Determine transconductance, output resistance, and gain.
\item
  Evaluate key performance metrics like slew rate and gain-bandwidth
  product.
\end{enumerate}

At the beginning of the notebook, fundamental technology parameters are
defined, including supply voltages, threshold voltages, mobility, and
channel-length modulation factors for NMOS and PMOS transistors.

\begin{Shaded}
\begin{Highlighting}[]
\CommentTok{\# Technology and model parameters}
\NormalTok{VDD }\OperatorTok{=} \FloatTok{1.8}     \CommentTok{\# Supply voltage [V]}
\NormalTok{VTHN }\OperatorTok{=} \FloatTok{0.4}    \CommentTok{\# NMOS threshold voltage [V]}
\NormalTok{VTHP }\OperatorTok{=} \OperatorTok{{-}}\FloatTok{0.4}   \CommentTok{\# PMOS threshold voltage [V]}
\NormalTok{mu\_n\_Cox }\OperatorTok{=} \FloatTok{200e{-}6}  \CommentTok{\# NMOS process transconductance parameter [A/V\^{}2]}
\NormalTok{mu\_p\_Cox }\OperatorTok{=} \FloatTok{100e{-}6}  \CommentTok{\# PMOS process transconductance parameter [A/V\^{}2]}
\NormalTok{lambda\_n }\OperatorTok{=} \FloatTok{0.1}     \CommentTok{\# Channel length modulation NMOS [1/V]}
\NormalTok{lambda\_p }\OperatorTok{=} \FloatTok{0.1}     \CommentTok{\# Channel length modulation PMOS [1/V]}
\end{Highlighting}
\end{Shaded}

These parameters are critical for calculating transistor drain current
in saturation:

\[
I_D = \frac{1}{2} \mu C_{ox} \frac{W}{L} (V_{GS} - V_{TH})^2 
\]

The differential pair (M1 and M2) defines the input transconductance of
the OTA. Sizing begins by selecting bias current and overdrive voltage
(V\(_{OV}\)):

\begin{Shaded}
\begin{Highlighting}[]
\CommentTok{\# Bias current for each NMOS transistor}
\NormalTok{ID1 }\OperatorTok{=} \FloatTok{100e{-}6}   \CommentTok{\# [A]}

\CommentTok{\# Overdrive voltage for M1 and M2}
\NormalTok{VOV1 }\OperatorTok{=} \FloatTok{0.2}     \CommentTok{\# [V]}

\CommentTok{\# Calculate W/L ratio for M1 and M2}
\NormalTok{WL1 }\OperatorTok{=} \DecValTok{2} \OperatorTok{*}\NormalTok{ ID1 }\OperatorTok{/}\NormalTok{ (mu\_n\_Cox }\OperatorTok{*}\NormalTok{ VOV1}\OperatorTok{**}\DecValTok{2}\NormalTok{)}
\BuiltInTok{print}\NormalTok{(}\SpecialStringTok{f"W/L for M1 and M2: }\SpecialCharTok{\{}\NormalTok{WL1}\SpecialCharTok{:.2f\}}\SpecialStringTok{"}\NormalTok{)}
\end{Highlighting}
\end{Shaded}

\begin{verbatim}
W/L for M1 and M2: 25.00
\end{verbatim}

In this example: - Bias current (ID1) is 100 µA per transistor. -
Overdrive voltage (VOV1) is 0.2 V. - Calculated W/L ensures M1 and M2
operate in saturation with desired transconductance.

The tail current source (M5) provides total bias current to the
differential pair:

\begin{Shaded}
\begin{Highlighting}[]
\CommentTok{\# Tail current source sizing (M5)}
\NormalTok{ID5 }\OperatorTok{=} \FloatTok{200e{-}6}   \CommentTok{\# Total bias current [A]}
\NormalTok{VOV5 }\OperatorTok{=} \FloatTok{0.2}     \CommentTok{\# Overdrive voltage [V]}

\NormalTok{WL5 }\OperatorTok{=} \DecValTok{2} \OperatorTok{*}\NormalTok{ ID5 }\OperatorTok{/}\NormalTok{ (mu\_n\_Cox }\OperatorTok{*}\NormalTok{ VOV5}\OperatorTok{**}\DecValTok{2}\NormalTok{)}
\BuiltInTok{print}\NormalTok{(}\SpecialStringTok{f"W/L for M5: }\SpecialCharTok{\{}\NormalTok{WL5}\SpecialCharTok{:.2f\}}\SpecialStringTok{"}\NormalTok{)}
\end{Highlighting}
\end{Shaded}

\begin{verbatim}
W/L for M5: 50.00
\end{verbatim}

M5 carries the combined current of M1 and M2, typically twice ID1.

PMOS current mirrors (M3 and M4) act as active loads for the
differential pair, impacting gain and output resistance:

\begin{Shaded}
\begin{Highlighting}[]
\CommentTok{\# Load transistors sizing (M3 and M4)}
\NormalTok{ID3 }\OperatorTok{=} \FloatTok{100e{-}6}   \CommentTok{\# [A]}
\NormalTok{VOV3 }\OperatorTok{=} \FloatTok{0.2}     \CommentTok{\# [V]}

\NormalTok{WL3 }\OperatorTok{=} \DecValTok{2} \OperatorTok{*}\NormalTok{ ID3 }\OperatorTok{/}\NormalTok{ (mu\_p\_Cox }\OperatorTok{*}\NormalTok{ VOV3}\OperatorTok{**}\DecValTok{2}\NormalTok{)}
\BuiltInTok{print}\NormalTok{(}\SpecialStringTok{f"W/L for M3 and M4: }\SpecialCharTok{\{}\NormalTok{WL3}\SpecialCharTok{:.2f\}}\SpecialStringTok{"}\NormalTok{)}
\end{Highlighting}
\end{Shaded}

\begin{verbatim}
W/L for M3 and M4: 50.00
\end{verbatim}

Lowering overdrive voltage (V\(_{OV3}\)) increases output resistance,
improving OTA voltage gain.

Once W/L ratios are calculated, device transconductance (g\(_m\)) and
output resistance (r\(_o\)) are derived:

\begin{Shaded}
\begin{Highlighting}[]
\NormalTok{gm1 }\OperatorTok{=}\NormalTok{ mu\_n\_Cox }\OperatorTok{*}\NormalTok{ WL1 }\OperatorTok{*}\NormalTok{ VOV1 }\OperatorTok{/} \DecValTok{2}
\NormalTok{ro1 }\OperatorTok{=} \DecValTok{1} \OperatorTok{/}\NormalTok{ (lambda\_n }\OperatorTok{*}\NormalTok{ ID1)}
\NormalTok{ro3 }\OperatorTok{=} \DecValTok{1} \OperatorTok{/}\NormalTok{ (lambda\_p }\OperatorTok{*}\NormalTok{ ID3)}

\BuiltInTok{print}\NormalTok{(}\SpecialStringTok{f"gm1: }\SpecialCharTok{\{}\NormalTok{gm1}\SpecialCharTok{:.2e\}}\SpecialStringTok{ S"}\NormalTok{)}
\BuiltInTok{print}\NormalTok{(}\SpecialStringTok{f"ro1: }\SpecialCharTok{\{}\NormalTok{ro1}\SpecialCharTok{:.2e\}}\SpecialStringTok{ Ω"}\NormalTok{)}
\BuiltInTok{print}\NormalTok{(}\SpecialStringTok{f"ro3: }\SpecialCharTok{\{}\NormalTok{ro3}\SpecialCharTok{:.2e\}}\SpecialStringTok{ Ω"}\NormalTok{)}
\end{Highlighting}
\end{Shaded}

\begin{verbatim}
gm1: 5.00e-04 S
ro1: 1.00e+05 Ω
ro3: 1.00e+05 Ω
\end{verbatim}

The OTA's DC gain is:

\[
A_v = g_{m1} \cdot (r_{o1} \parallel r_{o3})
\]

In Python:

\begin{Shaded}
\begin{Highlighting}[]
\NormalTok{Av }\OperatorTok{=}\NormalTok{ gm1 }\OperatorTok{*}\NormalTok{ (ro1 }\OperatorTok{*}\NormalTok{ ro3) }\OperatorTok{/}\NormalTok{ (ro1 }\OperatorTok{+}\NormalTok{ ro3)}
\BuiltInTok{print}\NormalTok{(}\SpecialStringTok{f"Voltage gain A\_v: }\SpecialCharTok{\{}\NormalTok{Av}\SpecialCharTok{:.2f\}}\SpecialStringTok{"}\NormalTok{)}
\end{Highlighting}
\end{Shaded}

\begin{verbatim}
Voltage gain A_v: 25.00
\end{verbatim}

\textbf{Evaluating Performance Metrics - Slew Rate and GBW:}

Slew rate and gain-bandwidth product (GBW) predict OTA dynamic
performance:

\begin{Shaded}
\begin{Highlighting}[]
\ImportTok{import}\NormalTok{ numpy }\ImportTok{as}\NormalTok{ np}

\NormalTok{CL }\OperatorTok{=} \FloatTok{1e{-}12}   \CommentTok{\# Load capacitance [F]}

\CommentTok{\# Slew rate calculation}
\NormalTok{SR }\OperatorTok{=}\NormalTok{ ID5 }\OperatorTok{/}\NormalTok{ CL}
\BuiltInTok{print}\NormalTok{(}\SpecialStringTok{f"Slew Rate: }\SpecialCharTok{\{}\NormalTok{SR}\SpecialCharTok{:.2e\}}\SpecialStringTok{ V/s"}\NormalTok{)}

\CommentTok{\# Gain Bandwidth Product (GBW)}
\NormalTok{GBW }\OperatorTok{=}\NormalTok{ gm1 }\OperatorTok{/}\NormalTok{ (}\DecValTok{2} \OperatorTok{*}\NormalTok{ np.pi }\OperatorTok{*}\NormalTok{ CL)}
\BuiltInTok{print}\NormalTok{(}\SpecialStringTok{f"Gain Bandwidth Product: }\SpecialCharTok{\{}\NormalTok{GBW}\SpecialCharTok{:.2e\}}\SpecialStringTok{ Hz"}\NormalTok{)}
\end{Highlighting}
\end{Shaded}

\begin{verbatim}
Slew Rate: 2.00e+08 V/s
Gain Bandwidth Product: 7.96e+07 Hz
\end{verbatim}

\begin{itemize}
\tightlist
\item
  \textbf{Slew Rate (SR)}: OTA's ability to respond to large signals.
\item
  \textbf{Gain Bandwidth Product (GBW)}: Small-signal frequency
  response.
\end{itemize}

The presented methodology offers a systematic, reproducible approach to
OTA design. Starting with hand calculations and validating through
simulation helps designers optimize OTA performance while balancing
speed, power, and gain. This lays groundwork for addressing noise,
mismatch, and layout parasitics. Further information is publicly
available at
\href{https://github.com/iic-jku/analog-circuit-design/blob/main/sizing/sizing_basic_ota.ipynb}{GitHub:
analog-circuit-design}.

\bookmarksetup{startatroot}

\chapter{Integrated Layouting}\label{integrated-layouting}

The following chapters describe how to design an integrated circuit
using the software KLayout. An introduction and setup for the software
will be given as well as some tips and tricks to easier work with the
software.

\section{Introduction into the fundamentals of integrated
layouting.}\label{introduction-into-the-fundamentals-of-integrated-layouting.}

Before designing on the wafer, the basics of the design process must be
understood. To that end an explanation of the software as well as the
physical layers is necessary to comprehend the design challenges.

When designing an integrated chip, the designer first needs to choose a
technology. Each technology has different design rules and properties,
which makes it suitable for the application. A technology usually comes
with a process design kit (PDK). A PDK contains a device library,
verification checks and technology data.

The device library contains the symbols used in the schematic and the
PCells/devices used in the layout. The verification checks consist of
the design rule check (DRC) as well as the layout vs.~schematic (LVS).
The technology data represents the basis of DRC and defines the physical
constraints of the technology (see Wikipedia contributors (2025))

The used technology was the sg13g2 technology from IHP. Which is a
BiCMOS technology with a 130nm minimum gate length. This PDK is open
source and in a preview status (see IHP‑GmbH (2025a))

To design the layout itself one first needs to be familiar with the
different masks and the structure of a waver. The following is a
simplified explanation of the different layers, their use and the
manufacturing process behind it. For this report this is kept simple and
will only focus on the layers used in the design. The simplified
explenation was heavily inspired by the chapters 2 to 5 in CMOS by
Jacaob Baker (see Baker (2010)).

The base is the bulk and a field oxide (FOX) layer on top of it. In this
technology the bulk is a p-substrate. The FOX layer is an insulator,
which separates devices from one another.

To make a device first the FOX layer needs to be opened (see
Figure~\ref{fig-doping_substrate}). That is done by etching away the
FOX. The active mask allows the designer to specify the areas where the
FOX will be opened. This is necessary for the doping process, which is
usually done by an n- or p-select mask. In sg12g2 the n-select mask does
not need to be specified. It is necessary to remove the FOX, otherwise
the doping specified in the n-select mask will be stopped by the FOX.
The select mask also needs to be bigger than the Active mask due to
misalignment during the manufacturing process.

\begin{figure}

\centering{

\includegraphics[width=6.22505in,height=2.27264in]{content/../images/sec_finns_property/image2.png}

}

\caption{\label{fig-doping_substrate}Doping process of the substrate}

\end{figure}%

The next layer is the poly layer. The poly layer forms the gate of the
transistors, the name is an abbreviation for Polysilicon. The poly can
be routed like normal metal layers, with the exception that it is for
more restive than the metal layers, so caution is advised for long poly
traces.

In the manufacturing process the thick poly layer prevents the n-select
mask to dope the part below the gate, therefore creating the drain and
source regions of the nmos (see Figure~\ref{fig-poly_manu}).

\begin{figure}

\centering{

\includegraphics[width=4.06736in,height=3.24028in]{content/../images/sec_finns_property/image4.png}

}

\caption{\label{fig-poly_manu}Figure : Manufacturing the poly layer in
the process}

\end{figure}%

Also note that the drain and source contacts are interchangeable, as
they are identical regions in the active area. The last two pictures
showed how an NMOS device is formed. As mentioned before, the n-select
mask doesn't exist in the used PDK, but the p-select mask does exist
with the label pSD.

As the PMOS is a negated NMOS, the designer first needs to place the p+
regions in an n-substrate. This is done by enveloping the transistor in
an NWell. The NWell is a region that can be specified by the designer
and works like an implant in the substrate. The FOX is also opened by
the active layer but pSD layer is used to implant the p+ region (see
Figure~\ref{fig-pmos_manu}).

\begin{figure}

\centering{

\includegraphics[width=6.92938in,height=2.36599in]{content/../images/sec_finns_property/image5.png}

}

\caption{\label{fig-pmos_manu}Side view of an PMOS transistor}

\end{figure}%

The N- or PMOS is formed in the bulk, but there are no electrical
connections to the pins of the device. Connections between devices can
be established using the poly or active layer, but both options are not
advised. Therefore, the first metal layer is used for connections
between devices. Connections from an active or poly layer region to the
metal layer can´t be made by overlaying the layers as they are separated
by an insulator. To connect these different layers the contact layer
(Cnt) is used. It connects active or poly layer to the metal 1 layer, by
removing the insulator over these regions (see
Figure~\ref{fig-metal_connection}). Note that the first connection
between these regions is always over the metal 1 layer. To switch
between the different Metal layers, vias can be used, similar to a
multilayer PCB design. These vias are formed by opening the insulator
and placing in an conductor like tungsten the opening. In the figure
this process is simplied.

\begin{figure}

\centering{

\includegraphics[width=6.10915in,height=3.93131in]{content/../images/sec_finns_property/image6.png}

}

\caption{\label{fig-metal_connection}Metal connection to the drain and
source contacts}

\end{figure}%

The MOSFET is a 4-pin device, consisting of source, gate, drain and
body. While the body often is omitted on a schematic level, in the
integrated layout the body of the MOSFET need to be tied to a defined
potential. The body is also called the substrate or bulk, in case of a
PMOS it is the NWell.

For the nmos this means connecting the substrate to VSS, while the NWell
of a PMOS will be pulled to VDD. To have multiple locations to where the
bulk or NWell connected to their potentials is advised.

The later chapters will work in the sg13g2 technology, the knowledge
acquired in this chapter can be mapped to this technology. The figure
(see Figure~\ref{fig-layer_comp}) compares the technology on the left,
with the learned knowledge on the right. Note that the PWell is not used
for NMOS devices only for ties.

\begin{figure}

\centering{

\includegraphics[width=6.98958in,height=1.42674in]{content/../images/sec_finns_property/image8.png}

}

\caption{\label{fig-layer_comp}ayer comparison between IHP and theory}

\end{figure}%

\section{KLayout}\label{klayout}

The previous chapter described the fundamentals of the integrated
layout, in this chapter these fundamentals are put into practice to
create the layout for the 5T-OTA outlined by the schematic by Haraled
Pretl(see Pretl et al. (2025))

The schematic outlined by Professor Pretl consists of 13 N- and PMOS
devices. These devices are defined within the PDK and are the
fundamental blocks, which make up the layout within the technology.
Within the layout software these are defined as PCells. As the schematic
designer works out the optimal design, one of the main variables to
change the behaviour of the MOSFETs, is the w/l ratio. This ratio
defines the width and length of the gate over the active area. These w/l
ratios are crucial to the layout process, as they define the physical
size of the devices. In the layout the w/l values were given by the
calculation of Professor Pretel. These w/l ratios will be used to
configure the PCells.

The layout program used is called KLayout. Similar to Xschem it works in
a strict hierarchical system. As the layout of the 5T-OTA will be used
in the Top-Level layout of the Biquad, the 5T-OTA itself is a cell, that
can be called multiple times, throughout the layout of the Top-Level. To
make it easier for the designer to connect the individual components of
the Top-Level together, the 5T-OTA just uses the metal 1 layer. The
interface points are the vias stacks. Therefore, making connections on
the higher layers easier, as the individual cells for the 5T-OTA only
reside on the most bottom layer.

In the following chapters DRC and LVS will be important concepts. DRC
(Desing Rule Check) describes the process of comparing the physical
layout to the constrains of the technology. This will throw an error,
for example if a ``trace'' on the Metal1 layer is too thin. LVS (Layout
Versus Schematic) describes the process of comparing the schematic to
the layout. This will throw an error when two nets are shorted or a
device it not correctly connected.

\subsection{The Setup of KLayout}\label{the-setup-of-klayout}

While the previous chapters showed only a cross section of the waver,
the designer only sees the waver from a top view, but needs to recall
the cross section, while working on the layout. This especially true for
the layer that can be used as conductors.

This chapter will explain how to set up KLayout and the tools necessary
to run DRC and LVS. It also shows an example way for a simple workflow.

Before the layout process begins it is paramount aggravate all the data
files necessary. To do that create a new folder labelled Layout and copy
your schematic design into this folder. Consider renaming your schematic
into letters only. Do not use special character similar to UNIX symbols,
only underscores are permitted.

To create a valid netlist, open the schematic in Xschem. Under
\emph{simulation-\textgreater LVS} make sure that ``\emph{LVS netlist +
Top level is a .subckt}'' and ``\emph{Upper case .SUBCKT and .ENDS}'' is
selected.\\
After selecting these options, run the netlist and simulation. In the
folder structure a new folder labelled simulations with file named
``\emph{YOUR\_PROJECT.spice}'' should appear. Copy that file into your
main layout directory.

To begin designing create a new layout file called the same as the
schematic by running ``Klayout -e''. The argument -e opens Klayout in
editing mode. To make the editing mode default, select
\emph{File-\textgreater Setup-\textgreater Application-\textgreater Editing
Mode-\textgreater Use} \emph{editing mode by default}.

To create a new layout, select ``\emph{New Layout}''. A wizard should
appear that allows the selection of the technology. It is important to
rename the Top Cell to match the name of your schematic so
``YOUR\_PROJECT''.\\
Please note that all the scripts later used are case sensitive, so
consider this when renaming files.

Upon first saving, select the GDS file standard and name your file
\emph{YOUR\_PROJECT.GDS}. The main layout directory should now contain
\emph{YOUR\_PROJECT.sch/.spice/.gds} files. As well as a simulation
folder. Now create a folder called LVS in this main layout directory.

For ease of access, create shell-script in your main layout folder to
run LVS. The command saved with in the file should be:\\

\begin{Shaded}
\begin{Highlighting}[]
\NormalTok{python /foss/pdks/ihp{-}sg13g2/libs.tech/klayout/tech/lvs/run\_lvs.py {-}{-}netlist=YOUR\_PROJECT.spice {-}{-}layout=YOUR\_PROJECT.gds {-}{-}run\_dir=LVS/}
\end{Highlighting}
\end{Shaded}

Try running this command, to ensure it runs correctly. If it fails,
saying something about ``use\_same\_circuit'', check the names of your
files, including the cell name of your layout. It's crucial that all
files are named the same.

This script will produce a. lvsdb file in the LVS directory. This file
contains the comparison between the netlist of the schematic, as well as
the extracted netlist of the layout.

As the design process is a constant back and forth with design and
running LVS and DRC, these steps are necessary to ensure a good
workflow.

Before designing we need to set up the tools in Klayout. For the DRC,
open \emph{Macros-\textgreater Macro Development} and select
\emph{DRC-\textgreater{[}Technology
sg13g2{]}-\textgreater sg13g2\_maximal}. Copy the contents of this
script and paste them into a new DRC script. Running this script will
provide the ``\emph{Marker Database Brower}'' from which the individual
DRC errors can be selected. As the error names and a semi-detailed
description can be found the ``\emph{SG13G2\_os\_layout\_rules.pdf}'' on
the Github of IHP (see IHP‑GmbH (2025b)).

To run LVS run the shell script and open the netlist browser under
\emph{Tool-\textgreater Netlist Browser}. To see the LVS errors click on
\emph{file-\textgreater open} and then navigate to
\emph{/LVS/YOUR\_PROJECT.lvsdb}. This only needs to be done one time, as
the LVS shell script will override the current .lvsdb file, when run
anew. Select \emph{File-\textgreater Reload} to get the current
LVS-data.

After loading the .lvsdb file the cross-reference tab will show the
discrepancies between the netlists. This is because no layout has been
created yet.

\begin{tcolorbox}[enhanced jigsaw, title=\textcolor{quarto-callout-tip-color}{\faLightbulb}\hspace{0.5em}{Some quality-of-life options}, colback=white, bottomrule=.15mm, coltitle=black, colbacktitle=quarto-callout-tip-color!10!white, toptitle=1mm, colframe=quarto-callout-tip-color-frame, rightrule=.15mm, arc=.35mm, leftrule=.75mm, bottomtitle=1mm, titlerule=0mm, toprule=.15mm, opacitybacktitle=0.6, opacityback=0, left=2mm, breakable]

\begin{itemize}
\item
  To provide a better overview of the used layers, right click on the
  right-hand side layer tab to open a pop up menu and select
  ``\emph{hide empty layers}''.
\item
  Disabling navigation by holding the middle mouse button: Check
  \emph{File-\textgreater Setup-\textgreater Navigation-\textgreater Zoom
  and Pan-\textgreater Mouse} alt mode. This allows the control more
  similar to Autodesk Fusion or KiCad.
\item
  Disabling the zoom when pasting:
  \emph{File-\textgreater Setup-\textgreater Navigation-\textgreater Zoom}
  and select \emph{Pan-\textgreater On paste-\textgreater Pan to pasted
  objects.} This option disables the full zoom to object, which can
  cause disorientation.
\item
  To show all the hierarchy and layers select
  \emph{Display-\textgreater Full Hierarchy}
\item
  To change to dark mode select
  \emph{File-\textgreater Setup-\textgreater Display-\textgreater Background-\textgreater Background
  Color-\textgreater\#42}
\end{itemize}

\end{tcolorbox}

The real layout process can now begin in Klayout. Different from PCB
design, the symbols on the schematic side do not need to be matched to a
footprint. Rather the symbol are the fundamental devices in the
technology. To work with the devices of the technology open the library
on the down-left under \emph{Libraries-\textgreater SG13\_dev}. The list
below shows the usable devices. By dragging and dropping them into the
design, the devices are placed in the layout.

In previous chapters mentioned the w/l ratios will now play a role as
the individual devices are configured. To configure a device or PCells,
double click it and navigate to PCell parameters and enter the W/L
ratios. Updating the PCell will change the device to the desired
parameters.

\section{The Layout Process}\label{the-layout-process}

The workflow is much more incremental than in PCB design, as one
implements device by device.\\
To that end, it is practical to copy the original schematic and delete
the already placed devices, to maintain an overview of the progress.
After every two to three cells, it is advised to run LVS and DRC.

Placing and configuring a device was already described, connecting
devices is usually done by the Metal1 layer. It is advised to work with
the metal layer alongside the hierarchy. So for the top level use the
higher metal layers like Metal3 to Metal6, while for the lower levels
use Metal1 to Metal2.

To connect different pins, use the LAYER.drawing layer. Conductors can
either be the GatPoly or the Metal layers. Keep in mind that
GatPoly.drawing has a higher resistance that Metal.drawing. To switch
between the metal layers, use the device called ``via\_stack'' and
configure it accordingly.

To connect the gate of a MOS device the GatPoly needs to be connected to
Metal1 through the Cnt layer. In contrast to the source and drain
contacts this need to be done by hand. See
(Figure~\ref{fig-gat_connection}) for a DRC clean connection of the gate
of an NMOS device.

\begin{figure}

\centering{

\includegraphics[width=4.96736in,height=3.18661in]{content/../images/sec_finns_property/image9.png}

}

\caption{\label{fig-gat_connection}Example DRC clean connection of the
gate}

\end{figure}%

In general, connect each device separately from each other. It might be
tempting to design the layout close together, with overlapping NWells,
Active and GatPoly areas, but this will lead to both DRC and LVS errors,
which will be difficult to address as each device is close together.
This is especially true, when first learning the process.

The layout of the 5T-OTA was done in multiple iterations, each iteration
taught valuable lessons. The next section presents these lessons,
outlining a specific workflow and offering tips and guidance for
beginners.

\begin{tcolorbox}[enhanced jigsaw, title=\textcolor{quarto-callout-tip-color}{\faLightbulb}\hspace{0.5em}{Tips and tricks for the Layout Process.}, colback=white, bottomrule=.15mm, coltitle=black, colbacktitle=quarto-callout-tip-color!10!white, toptitle=1mm, colframe=quarto-callout-tip-color-frame, rightrule=.15mm, arc=.35mm, leftrule=.75mm, bottomtitle=1mm, titlerule=0mm, toprule=.15mm, opacitybacktitle=0.6, opacityback=0, left=2mm, breakable]

\begin{itemize}
\item
  Run LVS and DRC every 1 to 3 devices!
\item
  Aligning structures is easier when using a crosshair.\\
  Enable this by checking View-\textgreater Crosshair Cursor
\item
  Use the path tool whenever possible.

  \begin{itemize}
  \item
    The tool has an adjustable width, which can be set on the bottom
    right menu.
  \item
    Set this width to 0.16 um to route DRC free on most layers.
  \item
    Do not use free form, as jagged shaped will create DRC errors.
  \end{itemize}
\item
  Use the shift key to snap to constrain vertically or horizontally.
\item
  To change the layer of structure, select the desired layer, select the
  structure and press shift+L
\item
  Copying and pasting a structure will result in duplicating the
  structure in place. Use the move tool to separate both structures.
\item
  If things get too cluttered use the layer menu to declutter the
  interface.

  \begin{itemize}
  \item
    Right clicking and selecting ``Visibility follows selection'' will
    only show the selected layer.
  \item
    Setting different tabs for routing, DRC checking or overviewing the
    general layout is a good option to keep things tidy and readable
  \item
    Changing the appearance of different layers is a good idea,
    especially for layers that overlap, like pSD or NWell.
  \end{itemize}
\item
  Use the polygon only on special occasions as it is harder to clear of
  DRC errors.
\end{itemize}

\begin{quote}
\textbf{Coming from discrete EDA software}

\begin{itemize}
\item
  The devices will not be placed automatically.
\item
  There is no way of changing the grid, as it is defined in technology.
\item
  There is no ``Ratsnest'' or ``Airal Connections''.
\item
  There is no DRC in the background, short circuits will only be caught
  by running the LVS.
\item
  There is no net highlighting.
\item
  There are no zones, but their absence will cause DRC errors.
\item
  There is no dragging, moving a device always means moving the
  ``traces'' by hand with it.
\item
  There are no footprints, the PCells are the fundamental footprints
  defined by the technology.
\item
  Text size can't be set but will adjust to the zoom level.
\item
  The GUI is sometimes not working, especially true for LVS.
\end{itemize}
\end{quote}

\end{tcolorbox}

\subsection{Working with LVS during
layouting}\label{working-with-lvs-during-layouting}

The Netlist Database Browser will provide an overview of all the nets
both in the layout and the schematic or ``Reference''. Each net has a
number in brackets behind it. This number describes all the pins it is
connected to. The netlists will only then match if these numbers are
matched.

During the layout process it is paramount to use the LAYER.text layers
and the text tool to assign labels to the GatPoly or Metal1 layers. To
select these layers, one may need to show all layers under
Layer-\textgreater Show all. To keep the comparison in the netlist
browser simple, name the nets of the layout according to the names of
the nets in the schematic. Please refer to the designer's etiquette,
while labeling the layout or schematic (see Pretl, Koefinger, and Dorrer
(2025)).

When working out the LVS it is advised to check mainly the devices in
the list. This will provide a comprehensive overview of the devices, and
the nets connected to them. Running LVS and DRC frequently will help
with the overview.

It is normal that even the correctly wired devices are shown as errors.
Especially the drain and source pins often swap places, which can't be
fixed by rotating the device. To clear the errors all nets connected to
the device need to be cleared first. Only when all nets are correctly
connected, will the device be clear of errors.

Having two nets on the same pit is always short circuit and should be
handled with the highest priory.

After a few MOS devices are placed, LVS will throw an error as their
body ``pin'' is not connected to a defined potential. In an prior
chapter the process of tying down the substrate was already explained.
In this technology the substrate is a p-substrate, this means the NWells
need to be tied down and there must be multiple tying down points for
the p-substrate (see Figure~\ref{fig-tiedown}). Please note that the
size of the masks shown has been altered for visibility reasons. Also
note that each PMOS needs its own tie to the VDD potential, while the
ties for the NMOS can be placed anywhere on the layout, as the substrate
is the bulk.

\begin{figure}

\centering{

\includegraphics[width=6.12245in,height=3.12569in]{content/../images/sec_finns_property/image11.png}

}

\caption{\label{fig-tiedown}Tying down different types of bulks
contacts. Not to scale.}

\end{figure}%

While using the netlist browser allows also for a rudimentary net
highlighting in the layout. The practicality depends on the number of
devices connected to the net. In general troubleshooting get more
difficult, the more pins are connected to the same net.

\subsection{Working with DRC during
layouting.}\label{working-with-drc-during-layouting.}

Running the DRC often and addressing them as soon as possible, allows
for a denser and generally more compact layout. It also reduces the time
spend reworking the layout after all devices have been connected.

As described in set up chapter the ``Marker Database Browser'' is a
built-in tool for addressing DRC errors. This tool shows the amount of
DRC errors as well as the type. After following the steps described in
the setup, one can open the Brower under Tool-\textgreater Marker
Browser. And running the DRC by selecting the run button.

The DRC will categorise in cells. Opening the DRC error log up will show
the name of the error. Names link to the
``\emph{SG13G2\_os\_layout\_rules.pdf}'' by IHP. While KLayout also
provides a description of the errors, it is best to use the DRC in
conjunction with the DRC-PDF. The DRC-PDF provides context and
dimensions for each error.

Working dually with the DRC-PDF and the DRC-Browser, allows the user to
first clarify the DRC errors and than locate it by selecting it in the
browser.

A very useful tool to clear DRC errors is the ``partial'' tool in the
toolbar. This allows to modify structures after they are placed. This
speed up the workflow immensely as the structures don't need to redrawn.

Some DRC errors can't be fixed at a given point of the design process.
These DRC errors are usually due to density errors. But as the layout is
not fully completed, cleaning these DRC errors has at best no impact on
the design. Nevertheless, these errors must be cleared before the tape
out.

\section{The layout of the 5T-OTA}\label{the-layout-of-the-5t-ota}

To design the 5T-OTA the schematic and the values of Professor Pretel
were used (see Pretl et al. (2025)).

As the export functionality of KLayout is rather limited, it is best to
open the .gds file in KLayout directly, but the is also a SVG in the
layout folder.

While not very space efficient the layout followed the location of the
MOSFETS in the schematic to ensure readability. This helped during the
layout process, as it gave a clear structure to the layout. This also
benefitted the understanding of the LVS and DRC during the layout
immensely.

As space is one of main drivers of cost in IC design, a final version
should be more mindful of the space used by the devices and traces. This
would also clear the density errors, due to a higher density, because of
the more compact design.

The design features multiple NWell and PWell Ties, to ensure that both
bulks and NWell are connected to their potentials. Each PMOS-device has
and NWell tie, which is closely located to the device. The PWell ties
surround the design, as described in the introductory chapter to the
integrated layout, the bulk shares one potential, but it is good
practice to not use a star formation, to ensure this potential is evenly
distributed.

One of the design goals of the design is the connection only on the
first metal layer, this should ensure that the connection layers above
are not influenced by the design. In addition, it provides a clear
connection point for the top-level schematic.

The 5T-OTA.gds could easily be imported as cell into an arbitrary top-
or mid-level design. Which concludes the design goal of this report.

\bookmarksetup{startatroot}

\chapter{Conclusion}\label{conclusion}

This project took us from a theoretical modelling of the filter with
transfer functions, to the first five-transistor prototype through two
reference designs by Prof.~Pretl and, finally, to a ``advanced'' OTA
driving a gm-C biquad filter. On the way several problems were
encountered:

\begin{itemize}
\item
  \textbf{Prototype OTA} - The first try of implementing an OTA resulted
  in a schematic which was very error prone and took a long time in the
  debug process. During this time valueable insights were gained
  regarding transistor usage, hierarchical netlisting and transient and
  ac testing. In the end the prototype OTA was abandoned for a design
  from Prof.~Pretl, as the required specifications could not be achieved
  with the prototype.
\item
  \textbf{Pretl's 5T-OTAs} -- Using the design for the 5T-OTA provided a
  good starting point for implementing the universal biquad filter. The
  width and length of the transistor channels were sized using the
  \(g_m/I_D\) method. Analysing this filter in the frequency domain
  revealed that the behaviour expected from the filter was not being
  simulated. The hypothesis here is that the OTAs could not provide
  enough gain for the circuit to work correctly.
\item
  \textbf{Gm-C biquad} -- With that insight experience the next step was
  sizing a new OTA, embedded four copies in a Baker-style biquad, and
  verified low-pass behaviour in Ngspice. Although the filter has yet to
  reach textbook performance, the design flow until the AC analysis ran
  without critical errors.
\end{itemize}

Beyond the circuits themselves other issue were tackled in the IC design
flow: absolute path handling inside Docker, strict symbol
classification, plain-text netlist debugging, and (just as important)
team communication. Working through those issues was as valuable as the
schematics that were produced.

\section{Outlook}\label{outlook}

At the current status of the project the expected behaviour of an
universal biquad could not be reproduced with self-build OTAs and the
Gm-c differential OTA setup is also not producing an expected low pass
filter. If the schematics are capable of producing an expected AC
behaviour, those simulation could be compared to modelled behaviour to
ensure correct operation behauviour. Additionally, the component
stabiltiy analysis would be redone with the used sizing of the
transistors to verify that the circuit is not prone to oszillations. All
these checks are important to be done before tape-out, as the actual
process of manufacturing an IC is time and ressource consuming and it is
important to be sure, that the behaviour of the IC is what is wanted.

As the layout of the 5T-OTA is finished, the next steps are the import
of the 5T-OTA Cell into the top-level layout. This can easily be
achieved, as KLayout provides and import function, that was already
tested. The design of the top-level should include the resistors and
capacitors needed for the biquad. In addition, bonding pads need to be
placed to connect the layout to a pin of the chosen package. LVS in the
top-level will be a challenge, as the exported netlist must account for
the hierarchy of the schematic. This will also extend the time needed to
run the LVS-script. When the top-level is layouted, the remaining
density DRC-violations need to be addressed. As the layout should be
complete at this point this can be done by the fill tool provided in
KLayout. In total only a little more work needs to be done on the
layout, to get ready for an eventual tape out.

\bookmarksetup{startatroot}

\chapter*{Bibliography}\label{bibliography}
\addcontentsline{toc}{chapter}{Bibliography}

\markboth{Bibliography}{Bibliography}

\phantomsection\label{refs}
\begin{CSLReferences}{1}{0}
\bibitem[\citeproctext]{ref-SQRLaw}
Alan Doolittle. 2025. {``{Lecture 25 - MOSFET Basics (Understanding with
Math}.''}
\url{https://alan.ece.gatech.edu/ECE3040/Lectures/Lecture25-MOSTransQuantitativeId-Vd-Vg.pdf}.

\bibitem[\citeproctext]{ref-analog_university_current_mirror}
Analog University. 2024. {``Chapter 11: Current Mirror.''}
\href{https://wiki.analog.com/university/courses/electronics/text/chapter-11\%20}{https://wiki.analog.com/university/courses/electronics/text/chapter-11
}.

\bibitem[\citeproctext]{ref-baker2010cmos}
Baker, R. Jacob. 2010. \emph{CMOS: Circuit Design, Layout, and
Simulation}. 3rd ed. IEEE Press / Wiley.

\bibitem[\citeproctext]{ref-MurmannGmId}
Boris Murmann. 2016. {``{Systematic Design of Analog Circuits Using
Pre-Computed Lookup Tables}.''}
\url{https://www.ieeetoronto.ca/wp-content/uploads/2020/06/20160226toronto_sscs.pdf}.

\bibitem[\citeproctext]{ref-fliege1991}
Fliege, Norbert. 1991. \emph{Systhemtheorie}. 1st ed. Stuttgart:
Teubner.

\bibitem[\citeproctext]{ref-hpretlacd}
H. Pretl and Michael Koefinger. 2025. {``{Analog Circuit Design}.''}
\url{https://iic-jku.github.io/analog-circuit-design/}.

\bibitem[\citeproctext]{ref-ihpOpenPDK}
IHP‑GmbH. 2025a. {``IHP Open‑source PDK: 130\,Nm BiCMOS (SG13G2) Open
Source Process Design Kit.''}
\url{https://github.com/IHP-GmbH/IHP-Open-PDK}.

\bibitem[\citeproctext]{ref-IHP_SG13G2_layout_rules}
---------. 2025b. {``SG13G2\_open‑source\_layout\_rules.pdf.''}
\url{https://github.com/IHP‑GmbH/IHP‑Open‑PDK/blob/main/ihp‑sg13g2/libs.doc/doc/SG13G2_os_layout_rules.pdf}.

\bibitem[\citeproctext]{ref-gm_acakids}
Kids, Academic. 2025. {``Transconductance.''}
\url{https://academickids.com/encyclopedia/index.php/Transconductance}.

\bibitem[\citeproctext]{ref-nerdsdostuff_current_mirror}
Nerds Do Stuff. 2024. {``Current Mirror: Working and Circuit.''}
\href{https://nerdsdostuff.com/electronic_circuits/current-mirror-working-and-circuit/\%20}{https://nerdsdostuff.com/electronic\_circuits/current-mirror-working-and-circuit/
}.

\bibitem[\citeproctext]{ref-pretl2025}
Pretl, Harald, Michael Koefinger, and Simon Dorrer. 2025. {``Analog
{Circuit} {Design}.''} \url{https://doi.org/10.5281/zenodo.14387481}.

\bibitem[\citeproctext]{ref-pretl2025_ota5t_schematic}
Pretl, Harald, Michael Koefinger, Simon Dorrer, and IIC‑JKU. 2025.
{``Ota‑5t.svg -- 5‑transistor OTA Schematic from the Analog Circuit
Design Course.''}
\url{https://github.com/iic-jku/analog-circuit-design/blob/main/xschem/ota-5t.svg}.

\bibitem[\citeproctext]{ref-pretl_git2025}
Pretl, Harald, and Georg Zachl. 2025. {``GitHub Repository of the
IIC-OSIC-TOOLS.''} Zenodo.
\url{https://doi.org/10.5281/zenodo.14634518}.

\bibitem[\citeproctext]{ref-aslkpro2012}
Rao, K. R. K., and C. P. Ravikumar. 2012. \emph{Analog System Lab Kit
PRO MANUAL}. Texas Instruments.

\bibitem[\citeproctext]{ref-razavi2018}
Razavi, Behzad. 2018. {``The Biquadratic Filter.''} \emph{IEEE
Solid-State Circuit Magazine}, 11--16.

\bibitem[\citeproctext]{ref-razavi2024}
---------. 2024. {``The Design of a Biquadratic Filter.''} \emph{IEEE
Solid-State Circuit Magazine}, 6--13.

\bibitem[\citeproctext]{ref-reisch2007}
Reisch, Michael. 2007. \emph{Elektronische Bauelemente - Funktion,
Grundschaltungen, Modellierung Mit SPICE}. 2nd ed. Stuttgart:
Springer-Verlag.

\bibitem[\citeproctext]{ref-renner2025}
Renner, Nils. 2025. {``Biquadratic IIR (SOS) Filters.''}

\bibitem[\citeproctext]{ref-RWalker}
Ross Walker. 2017. {``{Chapter 5, gm/ID - Based Design}.''}
\url{https://web02.gonzaga.edu/faculty/talarico/EE406/documents/gmid.pdf}.

\bibitem[\citeproctext]{ref-Approximating_OPAMP}
Schmid, H. 2000. {``Approximating the Universal Active Element.''}
\emph{IEEE Transactions on Circuits and Systems II: Analog and Digital
Signal Processing} 47 (11): 1160--69.
\url{https://doi.org/10.1109/82.885124}.

\bibitem[\citeproctext]{ref-gmId}
Silveira, F., D. Flandre, and P. G. A. Jespers. 1996. {``A Gm/Id Based
Methodology for the Design of CMOS Analog Circuits and Its Application
to the Synthesis of a Silicon-on-Insulator Micropower OTA.''} \emph{IEEE
Journal of Solid-State Circuits} 31 (9): 1314--19.
\url{https://doi.org/10.1109/4.535416}.

\bibitem[\citeproctext]{ref-Wangenheim2007}
Wangenheim, Lutz v. 2007. \emph{Aktive Filter Und Oszillatoren: Entwurf
Und Schaltungstechnik Mit Integrierten Bausteinen}. Berlin:
Springer-Verlag.

\bibitem[\citeproctext]{ref-wiki_opamps}
Wikipedia. 2025. {``Operationsverstärker.''}
\url{https://de.wikipedia.org/wiki/Operationsverst\%C3\%A4rker}.

\bibitem[\citeproctext]{ref-PDK_WIKI}
Wikipedia contributors. 2025. {``Process Design Kit.''}
\url{https://de.wikipedia.org/wiki/Process_Design_Kit}.

\end{CSLReferences}


\backmatter


\end{document}
